% Random Variables and Discrete Random Variables — Pure Mathematics with Statistics Upper Sixth — Cours signé OnBuch+
% Official MINESEC syllabus. Compiler avec : tectonic PMS11-random-variables-and-discrete-random-variables.tex
\newcommand{\DOCMATIERE}{Pure Mathematics with Statistics}
\newcommand{\DOCNIVEAU}{Upper Sixth}
\newcommand{\DOCMODULE}{Upper Sixth — Pure mathematics with statistics}
\newcommand{\DOCLECON}{11}
\newcommand{\DOCTITRE}{Random Variables and Discrete Random Variables}
% OnBuch+ — préambule « cours signés » ENRICHI (compilé avec tectonic / XeLaTeX)
% Système visuel « Pop » d'OnBuch+ : fond crème (jamais blanc), bordures encre
% épaisses, ombres « dures » décalées, titres Archivo Black en capitales.
% Version MATHÉMATIQUES : graphes (pgfplots/tikz), boîtes pédagogiques
% (définition, propriété, méthode, attention, le savais-tu, exercice résolu,
% exercice, TP, bilan), page de garde et signature OnBuch+ ancrée en bas.
%
% Macros attendues dans main.tex AVANT \input{preamble} :
%   \DOCMATIERE  \DOCNIVEAU  \DOCMODULE  \DOCLECON (numéro)  \DOCTITRE
\documentclass[11pt]{article}

\usepackage{fontspec}
\usepackage[english]{babel}
\usepackage[a4paper,margin=2cm,top=3.4cm,bottom=2.8cm,headheight=1.4cm,footskip=1.3cm]{geometry}
\usepackage{amsmath,amssymb,amsfonts}
\usepackage{mathtools} % \xmapsto, \coloneqq... souvent utilisés par les modèles
\usepackage{cancel} % simplifications barrées (bilans de forces)
% \abs{z} (module, valeur absolue) et \norm{u} (norme), utilisés par les modèles
\providecommand{\abs}{}\let\abs\relax\DeclarePairedDelimiter{\abs}{\lvert}{\rvert}
\providecommand{\norm}{}\let\norm\relax\DeclarePairedDelimiter{\norm}{\lVert}{\rVert}
\usepackage[table]{xcolor} % [table] : fournit aussi \rowcolors (alternance de lignes)
\usepackage{pagecolor}
\usepackage{tikz}
\usetikzlibrary{arrows.meta,positioning,calc,shapes.geometric,shapes.misc,shapes.arrows,decorations.pathmorphing,decorations.pathreplacing,decorations.markings,patterns,shadows,mindmap,trees,fit,backgrounds,matrix,intersections,angles,quotes}
\usepackage{pgfplots}
\usepackage[version=4]{mhchem} % équations nucléaires \ce{^{235}_{92}U}
\usepackage[european,siunitx]{circuitikz} % schémas électriques
\pgfplotsset{compat=1.18}
\usepgfplotslibrary{fillbetween,groupplots} % aires entre courbes (calcul intégral)
\usepackage{enumitem}
\usepackage{fancyhdr}
% [most] charge toutes les bibliothèques tcolorbox (listings, minted, raster,
% documentation...) et gonfle inutilement la mémoire TeX sur un document long
% (~300 boîtes) : on ne charge que ce qui est réellement utilisé.
\usepackage[skins,breakable]{tcolorbox}
\usepackage{graphicx}
\usepackage{colortbl}
\usepackage{booktabs}
\usepackage{tabularx}
\usepackage{diagbox} % case d'en-tête diagonale des tableaux à double entrée
% Colonnes extensibles centrée / gauche / droite (C, L, R), souvent utilisées par les modèles
\newcolumntype{C}{>{\centering\arraybackslash}X}
\newcolumntype{L}{>{\raggedright\arraybackslash}X}
\newcolumntype{R}{>{\raggedleft\arraybackslash}X}
\usepackage{array}
\usepackage{multicol}
\usepackage{siunitx}
% parse-numbers=false : accepte aussi \SI{2,5 \times 10^{-3}}{...}
\sisetup{locale=UK,output-decimal-marker={.},group-minimum-digits=4,parse-numbers=false}
\DeclareSIUnit\ohm{\ensuremath{\symup{\Omega}}} % U+2126 absent de la police
\DeclareSIUnit\division{div} % réglages d'oscilloscope (ms/div, V/div)
\DeclareSIUnit\tr{tr} % tours (vitesse de rotation, ex. \SI{1500}{\tr\per\minute})
\usepackage{qrcode}
% \degree : commande fréquemment utilisée par les modèles pour un angle ou une
% température en dehors de \SI{}{} (non fournie par les paquets chargés).
\providecommand{\degree}{^{\circ}}
% \qed (fin de démonstration), utilisé par les modèles sans amsthm
\providecommand{\qed}{\hfill$\square$}
% \barre : double barre des valeurs interdites dans un tableau de variations (tkz-tab)
\providecommand{\barre}{\ensuremath{\|}}
% environnement proof (amsthm non chargé) utilisé par les modèles
\ifdefined\proof\else\newenvironment{proof}[1][Proof]{\par\noindent\textit{#1.}\ }{\qed\par}\fi
\usepackage{lastpage}
\usepackage{hyperref}
\hypersetup{hidelinks}

% ============================================================
% Palette « Pop » OnBuch+ — mêmes hex que la classe P (plus_ui.dart)
% ============================================================
\definecolor{poporange}{HTML}{F59321}
\definecolor{popdark}{HTML}{E07A0C}
\definecolor{popcream}{HTML}{FFF6E8}
\definecolor{popink}{HTML}{1C1714}
\definecolor{poppurple}{HTML}{7A5AE0}
\definecolor{popgreen}{HTML}{2E9E6A}
\definecolor{poppink}{HTML}{E0457B}
\definecolor{popblue}{HTML}{2F7DE1}
\definecolor{popgold}{HTML}{C98A12}
\definecolor{popmuted}{HTML}{8A7F76}
\definecolor{popline}{HTML}{EADFD2}
\definecolor{popgray}{HTML}{8A7F76}
\colorlet{popgrey}{popgray}
% Alias tolérants (le modèle nomme parfois les couleurs différemment)
\colorlet{popwhite}{white}
\colorlet{popblack}{popink}
\colorlet{popred}{poppink}
\colorlet{popyellow}{popgold}
% Teintes claires pour fonds de tableaux / zones
\colorlet{popblueL}{popblue!10!white}
\colorlet{poporangeL}{poporange!14!white}
\colorlet{popgreenL}{popgreen!12!white}
\colorlet{poppurpleL}{poppurple!12!white}
\colorlet{poppinkL}{poppink!12!white}
\colorlet{popgoldL}{popgold!14!white}

\pagecolor{popcream}

% ============================================================
% Typographie
% ============================================================
\defaultfontfeatures{Ligatures=TeX}
\setmainfont{PlusJakartaSans-Medium.ttf}[Path=../../../fonts/,BoldFont=PlusJakartaSans-Bold.ttf]
% Maths en sans-serif (Fira Math) avec chiffres et lettres droites de Plus
% Jakarta, pour que les formules se fondent dans le texte.
\usepackage[math-style=ISO,bold-style=ISO]{unicode-math}
\setmathfont{FiraMath-Regular.otf}[Path=../../../fonts/]
\setmathfont{PlusJakartaSans-Medium.ttf}[Path=../../../fonts/,range={up/num,up/{Latin,latin},\mathup}]
% (icomma retiré : décimales avec point en anglais)
% Caractères mathématiques Unicode absents des polices de texte (Plus Jakarta,
% Archivo Black) : rendus par leur équivalent mathématique, en texte comme en
% formule — sinon ils s'affichent en carré vide (ex. « Arithmétique dans ℤ »).
\usepackage{newunicodechar}
\newunicodechar{ℝ}{\ensuremath{\mathbb{R}}}
\newunicodechar{ℤ}{\ensuremath{\mathbb{Z}}}
\newunicodechar{ℂ}{\ensuremath{\mathbb{C}}}
\newunicodechar{ℕ}{\ensuremath{\mathbb{N}}}
\newunicodechar{ℚ}{\ensuremath{\mathbb{Q}}}
\newunicodechar{𝔻}{\ensuremath{\mathbb{D}}}
\newunicodechar{α}{\ensuremath{\alpha}}
\newunicodechar{β}{\ensuremath{\beta}}
\newunicodechar{γ}{\ensuremath{\gamma}}
\newunicodechar{δ}{\ensuremath{\delta}}
\newunicodechar{ε}{\ensuremath{\varepsilon}}
\newunicodechar{θ}{\ensuremath{\theta}}
\newunicodechar{λ}{\ensuremath{\lambda}}
\newunicodechar{μ}{\ensuremath{\mu}}
\newunicodechar{σ}{\ensuremath{\sigma}}
\newunicodechar{τ}{\ensuremath{\tau}}
\newunicodechar{φ}{\ensuremath{\varphi}}
\newunicodechar{ψ}{\ensuremath{\psi}}
\newunicodechar{ω}{\ensuremath{\omega}}
\newunicodechar{Σ}{\ensuremath{\Sigma}}
% majuscules produites par \MakeUppercase dans les titres (α-aminés -> Α)
\newunicodechar{Α}{\ensuremath{\alpha}}
\newunicodechar{Β}{\ensuremath{\beta}}
\newunicodechar{Γ}{\ensuremath{\Gamma}}
\newunicodechar{Δ}{\ensuremath{\Delta}}
\newunicodechar{Ε}{\ensuremath{\varepsilon}}
\newunicodechar{Θ}{\ensuremath{\Theta}}
\newunicodechar{Λ}{\ensuremath{\Lambda}}
\newunicodechar{Μ}{\ensuremath{\mu}}
\newunicodechar{Τ}{\ensuremath{\tau}}
\newunicodechar{Φ}{\ensuremath{\Phi}}
\newunicodechar{Ψ}{\ensuremath{\Psi}}
\newunicodechar{Ω}{\ensuremath{\Omega}}
\newunicodechar{↦}{\ensuremath{\mapsto}}
\newunicodechar{∈}{\ensuremath{\in}}
\newunicodechar{∉}{\ensuremath{\notin}}
\newunicodechar{∧}{\ensuremath{\wedge}}
\newunicodechar{⇔}{\ensuremath{\Leftrightarrow}}
\newunicodechar{⟺}{\ensuremath{\Longleftrightarrow}}
\newunicodechar{⇒}{\ensuremath{\Rightarrow}}
\newunicodechar{⟹}{\ensuremath{\Longrightarrow}}
\newunicodechar{∘}{\ensuremath{\circ}}
\newunicodechar{∪}{\ensuremath{\cup}}
\newunicodechar{∩}{\ensuremath{\cap}}
\newunicodechar{≡}{\ensuremath{\equiv}}
\newunicodechar{⊂}{\ensuremath{\subset}}
\newunicodechar{⊥}{\ensuremath{\perp}}
\newunicodechar{∃}{\ensuremath{\exists}}
\newunicodechar{∀}{\ensuremath{\forall}}
\newunicodechar{∛}{\ensuremath{\sqrt[3]{\,}}}
\newunicodechar{∜}{\ensuremath{\sqrt[4]{\,}}}
\newunicodechar{⃗}{\ensuremath{{}^{\to}}}
\newunicodechar{ⁿ}{\ifmmode^{n}\else\textsuperscript{\ensuremath{n}}\fi}
\newunicodechar{ˣ}{\ifmmode^{x}\else\textsuperscript{\ensuremath{x}}\fi}
\newunicodechar{ᵇ}{\ifmmode^{b}\else\textsuperscript{\ensuremath{b}}\fi}
\newunicodechar{ʳ}{\ifmmode^{r}\else\textsuperscript{\ensuremath{r}}\fi}
\newunicodechar{ᵘ}{\ifmmode^{u}\else\textsuperscript{\ensuremath{u}}\fi}
\newunicodechar{ʸ}{\ifmmode^{y}\else\textsuperscript{\ensuremath{y}}\fi}
\newunicodechar{ᵃ}{\ifmmode^{a}\else\textsuperscript{\ensuremath{a}}\fi}
\newunicodechar{ᵉ}{\ifmmode^{e}\else\textsuperscript{\ensuremath{e}}\fi}
\newunicodechar{ᵏ}{\ifmmode^{k}\else\textsuperscript{\ensuremath{k}}\fi}
\newunicodechar{ᵖ}{\ifmmode^{p}\else\textsuperscript{\ensuremath{p}}\fi}
\newunicodechar{ᴺ}{\ifmmode^{N}\else\textsuperscript{\ensuremath{N}}\fi}
\newunicodechar{ᶜ}{\ifmmode^{c}\else\textsuperscript{\ensuremath{c}}\fi}
\newunicodechar{ʰ}{\ifmmode^{h}\else\textsuperscript{\ensuremath{h}}\fi}
\newunicodechar{ᵠ}{\ifmmode^{q}\else\textsuperscript{\ensuremath{q}}\fi}
\newunicodechar{ₙ}{\ifmmode_{n}\else\textsubscript{\ensuremath{n}}\fi}
\newunicodechar{ₐ}{\ifmmode_{a}\else\textsubscript{\ensuremath{a}}\fi}
\newunicodechar{ᵢ}{\ifmmode_{i}\else\textsubscript{\ensuremath{i}}\fi}
\newunicodechar{ᵣ}{\ifmmode_{r}\else\textsubscript{\ensuremath{r}}\fi}
\newunicodechar{ₖ}{\ifmmode_{k}\else\textsubscript{\ensuremath{k}}\fi}
\newunicodechar{ᵦ}{\ifmmode_{\beta}\else\textsubscript{\ensuremath{\beta}}\fi}
\newunicodechar{ₓ}{\ifmmode_{x}\else\textsubscript{\ensuremath{x}}\fi}
\newfontfamily\popdisplay{ArchivoBlack-Regular.ttf}[Path=../../../fonts/]
\newfontfamily\poplogo{SpaceGrotesk-Bold.ttf}[Path=../../../fonts/]
\newfontfamily\popsemi{PlusJakartaSans-SemiBold.ttf}[Path=../../../fonts/]
\newfontfamily\popxbold{PlusJakartaSans-ExtraBold.ttf}[Path=../../../fonts/]
\newfontfamily\popxboldls{PlusJakartaSans-ExtraBold.ttf}[Path=../../../fonts/,LetterSpace=3.0]

\setlist[enumerate]{leftmargin=1.7em,itemsep=5pt,topsep=4pt}
\setlist[itemize]{leftmargin=1.7em,itemsep=4pt,topsep=4pt,label={\color{poporange}\textbullet}}
\setlength{\parindent}{0pt}
\setlength{\parskip}{5pt}
\renewcommand{\arraystretch}{1.25}

% pgfplots : style Pop par défaut
\pgfplotsset{
  every axis/.append style={axis line style={popink,line width=1pt},tick style={popink},
    grid style={popline,line width=0.5pt},label style={font=\small},tick label style={font=\footnotesize}},
  popcurve/.style={poporange,line width=1.8pt,no markers,smooth},
}
% Même style disponible dans un \draw TikZ simple (hors pgfplots)
\tikzset{popcurve/.style={poporange,line width=1.8pt}}
\tikzset{popgrid/.style={popline,very thin}}
\colorlet{popcurve}{poporange} % le nom du style est parfois utilisé comme couleur

% ============================================================
% Pill / squiggle
% ============================================================
\newcommand{\poppill}[3]{%
  \tikz[baseline=-0.55ex]\node[draw=popink,line width=1.2pt,rounded corners=2.6pt,%
    fill=#2,inner xsep=6.5pt,inner ysep=3pt]{%
    \popxboldls\fontsize{7}{7}\selectfont\color{#3}\MakeUppercase{#1}};%
}
\newcommand{\popsquiggle}[1]{%
  \tikz[baseline=-0.4ex]{
    \draw[#1,line width=2.6pt,line cap=round]
      plot [smooth] coordinates {(0,0.09) (0.34,0.01) (0.68,0.21) (1.02,0.01) (1.36,0.10)};
  }%
}
% Mot-clé surligné (marqueur orange sous le texte)
\newcommand{\cle}[1]{{\popxbold\color{popdark}#1}}

% ============================================================
% En-tête / pied
% ============================================================
\pagestyle{fancy}
\fancyhf{}
\renewcommand{\headrulewidth}{0pt}
\renewcommand{\footrulewidth}{0pt}
\lhead{\raisebox{-0.15ex}{\poplogo\fontsize{13}{13}\selectfont\color{popink}OnBuch\color{poporange}+}}
\rhead{\poppill{\DOCMATIERE\ \textbullet\ \DOCNIVEAU\ \textbullet\ Chapter \DOCLECON}{white}{popink}}
\lfoot{\popxbold\fontsize{7.6}{7.6}\selectfont\color{popmuted}\MakeUppercase{Signed course OnBuch+}\ \textbullet\ {\popsemi\DOCTITRE}}
\rfoot{\tikz[baseline=-0.6ex]\node[rounded corners=8.5pt,fill=popink,minimum height=17pt,minimum width=17pt,inner xsep=5pt,inner sep=0]{\popxbold\fontsize{8}{8}\selectfont\color{popcream}\thepage\,/\,\pageref*{LastPage}};}

% ============================================================
% Titres
% ============================================================
\newcommand{\chaptitre}[2]{%
  \begin{tcolorbox}[enhanced,colback=poporange,colframe=popink,boxrule=2.6pt,arc=11pt,
      drop shadow=popink,left=15pt,right=15pt,top=12pt,bottom=11pt,before skip=3pt,after skip=18pt]
    \poppill{#2}{popink}{popcream}\par\vspace{9pt}
    {\raggedright\hyphenpenalty=10000\exhyphenpenalty=10000\popdisplay\fontsize{22}{24}\selectfont\color{popcream}\MakeUppercase{#1}\par}
  \end{tcolorbox}
}
% Section numérotée : pastille orange avec le numéro + titre Archivo
\newcounter{popsec}
\newcommand{\coursec}[1]{%
  \stepcounter{popsec}\setcounter{popsub}{0}%
  \par\vspace{16pt}\noindent
  \begin{minipage}{\linewidth}
  \tikz[baseline=-0.7ex]\node[draw=popink,line width=1.6pt,fill=poporange,rounded corners=4pt,minimum size=22pt,inner sep=0]{\popdisplay\fontsize{12}{12}\selectfont\color{popcream}\thepopsec};%
  \hspace{8pt}{\popdisplay\fontsize{15}{17}\selectfont\color{popink}\MakeUppercase{#1}}\par
  \vspace{4pt}\noindent{\color{poporange}\rule{36pt}{3.6pt}}
  \end{minipage}\par\nopagebreak\vspace{8pt}
}
\newcounter{popsub}
\newcommand{\courssub}[1]{%
  \stepcounter{popsub}%
  \par\vspace{9pt}\noindent{\popxbold\fontsize{11.5}{13}\selectfont\color{popink}{\color{poporange}\thepopsec.\thepopsub}\hspace{6pt}#1}\par\nopagebreak\vspace{4pt}
}

% ============================================================
% Boîtes pédagogiques — même recette : fond blanc, bordure épaisse colorée,
% ombre dure encre, badge plein. Titre optionnel [#1].
% ============================================================
\tcbset{popbase/.style={breakable,enhanced,colback=white,boxrule=2.3pt,arc=9pt,
  shadow={3pt}{-3pt}{0pt}{popink},left=14pt,right=14pt,top=11pt,bottom=11pt,before skip=11pt,after skip=15pt,
  fontupper=\popsemi\fontsize{10.6}{14.5}\selectfont\color{popink},
  fontlower=\popsemi\fontsize{10.6}{14.5}\selectfont\color{popink}}}
% Coche dessinée (le glyphe ✓ n'existe pas dans les polices du document)
\newcommand{\popcheck}[1]{\tikz[baseline=-0.5ex]\draw[#1,line width=1.6pt,line cap=round,line join=round](-0.13,0.0)--(-0.04,-0.09)--(0.13,0.1);}
\providecommand{\coche}{\popcheck{popgreen}}
\providecommand{\resultat}[1]{\par\smallskip\noindent\fcolorbox{popgreen}{popgreenL}{\parbox{\dimexpr\linewidth-2\fboxsep-2\fboxrule\relax}{\textbf{Result.} #1}}\par\smallskip}
\newcommand{\popboxhead}[3]{\poppill{#1}{#2}{white}\ifx\relax#3\relax\else\hspace{6pt}{\popxbold\fontsize{10.6}{12}\selectfont\color{popink}#3}\fi\par\vspace{7pt}}

\newtcolorbox{aretenir}[1][]{popbase,colframe=popgreen,before upper={\popboxhead{Key points}{popgreen}{#1}}}
\newtcolorbox{exemplebox}[1][]{popbase,colframe=poporange,before upper={\popboxhead{Exemple}{poporange}{#1}}}
\newtcolorbox{definition}[1][]{popbase,colframe=popblue,before upper={\popboxhead{Definition}{popblue}{#1}}}
\newtcolorbox{propriete}[1][]{popbase,colframe=poppurple,before upper={\popboxhead{Property}{poppurple}{#1}}}
\newtcolorbox{methode}[1][]{popbase,colframe=popgold,before upper={\popboxhead{Method}{popgold}{#1}}}
\newtcolorbox{attention}[1][]{popbase,colframe=poppink,before upper={\popboxhead{Warning}{poppink}{#1}}}
\newtcolorbox{experience}[1][]{popbase,colframe=popblue,colback=popblueL,before upper={\popboxhead{Experiment}{popblue}{#1}}}
% « Le savais-tu ? » : carte violette pleine (contexte, culture, Cameroun)
\newtcolorbox{savaistu}[1][]{popbase,colback=poppurple,colframe=popink,
  fontupper=\popsemi\fontsize{10.4}{14.2}\selectfont\color{white},
  before upper={\poppill{Did you know?}{white}{poppurple}\ifx\relax#1\relax\else\hspace{6pt}{\popxbold\color{white}#1}\fi\par\vspace{7pt}}}
% Exercice résolu : énoncé en haut, \tcblower puis la solution sur fond crème
\newcounter{popexo}\newcounter{popexr}
\newtcolorbox{exoresolu}[1][]{popbase,colframe=poporange,colbacklower=popcream,
  segmentation style={popink,line width=1.2pt,dashed},
  before upper={\stepcounter{popexr}\popboxhead{Worked example \thepopexr}{poporange}{#1}},
  before lower={\poppill{Solution}{popink}{popcream}\par\vspace{6pt}}}
% Exercice (non résolu en place) — la correction est dans la partie Corrigés
\newtcolorbox{exercice}[1][]{popbase,colframe=popink,
  before upper={\stepcounter{popexo}\popboxhead{Exercise \thepopexo}{popink}{#1}}}
% Corrigé
\newtcolorbox{corrige}[1][]{popbase,colframe=popgreen,colback=popgreenL,
  before upper={\popboxhead{Solution}{popgreen}{#1}}}
% Objectifs / prérequis / bilan
\newtcolorbox{objectifs}{popbase,colframe=popink,colback=poporangeL,
  before upper={\popboxhead{Chapter objectives}{popink}{}}}
\newtcolorbox{prerequis}{popbase,colframe=popink,colback=popblueL,
  before upper={\popboxhead{Prerequisites}{popblue}{}}}
\newtcolorbox{bilan}{popbase,colframe=popink,colback=white,boxrule=2.8pt,
  before upper={\popboxhead{Fiche bilan}{poporange}{L'essentiel à connaître pour l'examen}}}
% Figure encadrée avec légende
\newtcolorbox{popfigure}[1][]{enhanced,breakable=false,colback=white,colframe=popline,boxrule=1.4pt,arc=8pt,
  left=10pt,right=10pt,top=10pt,bottom=8pt,before skip=10pt,after skip=12pt,halign=center,
  lower separated=false,halign lower=center,
  fontlower=\popsemi\fontsize{9}{11.5}\selectfont\color{popmuted},
  #1}
\newcommand{\legende}[1]{\par\vspace{6pt}{\popsemi\fontsize{9}{11.5}\selectfont\color{popmuted}#1}}

% ============================================================
% Signature OnBuch+
% ============================================================
\newcommand{\poplogoinline}[1]{{\poplogo\fontsize{#1}{#1}\selectfont\color{popink}OnBuch\color{poporange}+}}
\newcommand{\signatureonbuch}{%
  \begin{tcolorbox}[enhanced,colback=popink,colframe=popink,boxrule=2.2pt,arc=10pt,
      drop shadow=poporange,left=14pt,right=14pt,top=10pt,bottom=10pt,before skip=0pt,after skip=0pt]
    \tikz[baseline=-0.7ex]\node[circle,fill=popgreen,draw=popcream,line width=1.3pt,minimum size=22pt,inner sep=0]{\popcheck{white}};%
    \hspace{9pt}\begin{minipage}[c]{0.62\linewidth}
      {\popxbold\fontsize{12}{13}\selectfont\color{popcream}Signed\ \textbullet\ {\poplogo OnBuch\color{poporange}+}}\par\vspace{2pt}
      {\raggedright\popsemi\fontsize{8}{10.5}\selectfont\color{popline}\MakeUppercase{OnBuch+ teaching team}\\\MakeUppercase{Follows the official MINESEC syllabus}\par}
    \end{minipage}\hfill
    {\poplogo\fontsize{22}{22}\selectfont\color{popcream}OnBuch\color{poporange}+}
  \end{tcolorbox}%
}

% ============================================================
% Page de garde
% #1 titre · #2 matière · #3 niveau · #4 module · #5 n° leçon · #6 durée ·
% #7 description / objectifs (texte ou itemize)
% ============================================================
\newcommand{\pagegarde}[7]{%
  \thispagestyle{empty}
  \newgeometry{a4paper,margin=1.7cm,top=1.6cm,bottom=1.4cm}%
  \begin{tikzpicture}[remember picture,overlay]
    \fill[poporange!18] ([xshift=-3.2cm,yshift=-3.6cm]current page.north east) circle (4.6cm);
    \fill[poppurple!14] ([xshift=2.4cm,yshift=7.6cm]current page.south west) circle (3.8cm);
  \end{tikzpicture}%
  {\poplogo\fontsize{30}{30}\selectfont\color{popink}OnBuch\color{poporange}+}\hfill
  \raisebox{0.6ex}{\poppill{Signed course \textbullet\ Premium}{popink}{popcream}}\par
  \vspace{1pt}\popsquiggle{poppurple}\par
  \vspace{8pt}
  \begin{tcolorbox}[enhanced,colback=poporange,colframe=popink,boxrule=2.8pt,arc=13pt,
      drop shadow=popink,left=18pt,right=18pt,top=12pt,bottom=13pt,after skip=10pt]
    \poppill{#2 \textbullet\ #3}{popink}{popcream}\hspace{5pt}\poppill{#4}{white}{popink}\par\vspace{8pt}
    {\popdisplay\fontsize{14}{14}\selectfont\color{popink}CHAPTER #5}\par\vspace{4pt}
    {\raggedright\hyphenpenalty=10000\exhyphenpenalty=10000\popdisplay\fontsize{25}{27}\selectfont\color{popcream}\MakeUppercase{#1}\par}
    \vspace{8pt}
    {\popxbold\fontsize{9.5}{9.5}\selectfont\color{popink}\MakeUppercase{Suggested time: #6}}
  \end{tcolorbox}
  \begin{tcolorbox}[enhanced,colback=white,colframe=popink,boxrule=2.2pt,arc=10pt,
      drop shadow=popink,left=16pt,right=16pt,top=10pt,bottom=10pt,after skip=8pt]
    \poppill{In this chapter}{poppurple}{white}\par\vspace{5pt}
    \popsemi\fontsize{9.3}{12.6}\selectfont\color{popink}#7
  \end{tcolorbox}
  \noindent\begin{minipage}[t]{0.487\linewidth}
    \begin{tcolorbox}[enhanced,colback=popgreen,colframe=popink,boxrule=2.2pt,arc=10pt,
        drop shadow=popink,left=11pt,right=11pt,top=10pt,bottom=10pt,height=3.1cm,valign=center]
      \tcbox[colback=white,colframe=popink,boxrule=1.3pt,arc=5pt,boxsep=0pt,nobeforeafter,
        left=3pt,right=3pt,top=3pt,bottom=3pt,valign=center]{\qrcode[height=1.9cm]{https://play.google.com/store/apps/details?id=com.luvvix.onbuch&hl=en}}%
      \hspace{8pt}\begin{minipage}[c]{0.5\linewidth}
        {\popdisplay\fontsize{11}{12.5}\selectfont\color{white}\MakeUppercase{Download OnBuch}}\par\vspace{4pt}
        {\raggedright\popsemi\fontsize{8.4}{11}\selectfont\color{white}Free \textbullet\ Google Play\\AI tutor, past papers, results\par}
      \end{minipage}
    \end{tcolorbox}
  \end{minipage}\hfill
  \begin{minipage}[t]{0.487\linewidth}
    \begin{tcolorbox}[enhanced,colback=poppurple,colframe=popink,boxrule=2.2pt,arc=10pt,
        drop shadow=popink,left=11pt,right=11pt,top=10pt,bottom=10pt,height=3.1cm,valign=center]
      \tikz[baseline=-0.7ex]\node[circle,fill=white,draw=popink,line width=1.3pt,minimum size=1.6cm,inner sep=0]{\popdisplay\fontsize{24}{24}\selectfont\color{poppurple}+};%
      \hspace{8pt}\begin{minipage}[c]{0.6\linewidth}
        {\popdisplay\fontsize{11}{12.5}\selectfont\color{white}\MakeUppercase{Go OnBuch+}}\par\vspace{4pt}
        {\raggedright\popsemi\fontsize{8.4}{11}\selectfont\color{white}All signed courses, unlimited AI tutor, PDF sheets — OnBuch+ tab in the app\par}
      \end{minipage}
    \end{tcolorbox}
  \end{minipage}
  \vspace{8pt}
  \popcontenu
  \vfill
  \signatureonbuch
  \restoregeometry
  \newpage
}

% Rangée « Ce cours contient » (page de garde)
\newcommand{\poptile}[3]{% #1 couleur #2 gros texte #3 légende
  \begin{tcolorbox}[enhanced,colback=#1,colframe=popink,boxrule=2pt,arc=9pt,shadow={2.5pt}{-2.5pt}{0pt}{popink},
      left=6pt,right=6pt,top=9pt,bottom=9pt,halign=center,valign=center,height=1.75cm,width=\linewidth,nobeforeafter]
    {\popdisplay\fontsize{12.5}{13}\selectfont\color{white}\MakeUppercase{#2}}\par\vspace{3pt}
    {\centering\popsemi\fontsize{7.8}{9.5}\selectfont\color{white}#3\par}
  \end{tcolorbox}}
\newcommand{\popcontenu}{%
  \par\noindent{\popxboldls\fontsize{7.5}{7.5}\selectfont\color{popmuted}THIS SIGNED COURSE CONTAINS}\par\vspace{7pt}
  \noindent\begin{minipage}[t]{0.235\linewidth}\poptile{popblue}{Course}{complete, illustrated, step by step}\end{minipage}\hfill
  \begin{minipage}[t]{0.235\linewidth}\poptile{poporange}{Methods}{and worked examples}\end{minipage}\hfill
  \begin{minipage}[t]{0.235\linewidth}\poptile{poppink}{Exercises}{exam style + solutions}\end{minipage}\hfill
  \begin{minipage}[t]{0.235\linewidth}\poptile{popgreen}{Summary sheet}{the essentials to remember}\end{minipage}\par}

% Sommaire
\newtcolorbox{sommaire}{enhanced,colback=white,colframe=popink,boxrule=2.2pt,arc=10pt,
  drop shadow=popink,left=18pt,right=18pt,top=13pt,bottom=13pt,before skip=6pt,after skip=20pt,
  before upper={\poppill{Contents}{poppurple}{white}\par\vspace{9pt}\popsemi\fontsize{10.6}{14.5}\selectfont\color{popink}}}

% Signature de fin de cours — poussée tout en bas de la dernière page
\newcommand{\findecours}{%
  \par\vfill\nopagebreak
  \begin{center}\popsquiggle{poporange}\end{center}\vspace{4pt}
  \signatureonbuch
}
\AtBeginDocument{\def\llbracket{\mathopen{[\mkern-3.5mu[}}\def\rrbracket{\mathclose{]\mkern-3.5mu]}}} % intervalles d'entiers (glyphe ⟦ absent de la police)
% Glyphes absents de Fira Math : redéfinis à partir de glyphes disponibles
\AtBeginDocument{%
  \def\setminus{\mathbin{\backslash}}\let\smallsetminus\setminus
  \def\vdots{\mathord{\vbox{\baselineskip3.2pt\lineskiplimit0pt\kern4pt\hbox{.}\hbox{.}\hbox{.}}}}%
  \def\ddots{\mathinner{\mkern1mu\raise6.5pt\hbox{.}\mkern2mu\raise3.6pt\hbox{.}\mkern2mu\raise0.7pt\hbox{.}\mkern1mu}}%
  \def\triangle{\mathord{\scalebox{0.9}{$\symup{\Delta}$}}}%
  \def\nabla{\mathord{\raisebox{\depth}{\scalebox{1}[-1]{$\symup{\Delta}$}}}}%
  \def\checkmark{\popcheck{popdark}}%
  \def\star{\mathbin{\ast}}\def\bigstar{\mathord{\ast}}%
  \def\diamond{\mathbin{\scalebox{0.6}{$\Diamond$}}}%
  \def\bigcup{\mathop{\mathchoice{\raisebox{-0.3em}{\scalebox{1.7}{$\cup$}}}{\scalebox{1.25}{$\cup$}}{\cup}{\cup}}\displaylimits}%
  \def\bigcap{\mathop{\mathchoice{\raisebox{-0.3em}{\scalebox{1.7}{$\cap$}}}{\scalebox{1.25}{$\cap$}}{\cap}{\cap}}\displaylimits}%
  \def\lesseqqgtr{\mathrel{\lessgtr}}\def\bigoplus{\mathop{\oplus}\displaylimits}%
}
\providecommand{\ssi}{if and only if} % abréviation fréquente
\AtBeginDocument{\providecommand{\wideparen}{\overparen}} % arc de cercle

% Fonctions usuelles que les modèles utilisent sans que amsmath les fournisse
\providecommand{\arsinh}{\operatorname{arsinh}}\providecommand{\arcosh}{\operatorname{arcosh}}
\providecommand{\artanh}{\operatorname{artanh}}\providecommand{\arsech}{\operatorname{arsech}}
\providecommand{\arcsch}{\operatorname{arcsch}}\providecommand{\arcoth}{\operatorname{arcoth}}
\providecommand{\arcsinh}{\operatorname{arsinh}}\providecommand{\arccosh}{\operatorname{arcosh}}
\providecommand{\arctanh}{\operatorname{artanh}}\providecommand{\sech}{\operatorname{sech}}
\providecommand{\csch}{\operatorname{csch}}\providecommand{\arcsec}{\operatorname{arcsec}}
\providecommand{\arccsc}{\operatorname{arccsc}}\providecommand{\arccot}{\operatorname{arccot}}
\providecommand{\sgn}{\operatorname{sgn}}\providecommand{\lcm}{\operatorname{lcm}}
\providecommand{\Var}{\operatorname{Var}}\providecommand{\Cov}{\operatorname{Cov}}

%% FIGFIX-BEGIN v4 — correctifs globaux des figures (docs/onbuch-generation/tools/figfix)
\usepackage{adjustbox}\usepackage{etoolbox}
% toute figure TikZ/pgfplots plus large que la ligne est réduite à la largeur disponible
\BeforeBeginEnvironment{tikzpicture}{\begin{adjustbox}{max width=\linewidth}}
\AfterEndEnvironment{tikzpicture}{\end{adjustbox}}
% légendes pgfplots lisibles par-dessus les courbes
\pgfplotsset{every axis/.append style={legend style={fill=white,fill opacity=0.92,text opacity=1,draw=black!15,inner sep=3pt}}}
% étiquettes d'arêtes (syntaxe edge["..."]) avec fond blanc
\tikzset{every edge quotes/.append style={fill=white,inner sep=1.2pt}}
%% FIGFIX-END
\begin{document}

\pagegarde{Random Variables and Discrete Random Variables}{Pure Mathematics with Statistics}{Upper Sixth}{Upper Sixth — Pure mathematics with statistics}{11}{9 periods}{This chapter introduces random variables as numerical summaries of random experiments and develops the full toolkit for discrete random variables: probability distributions, expectation, variance, cumulative distribution functions, mode, median, sums of independent variables and probability generating functions. It provides the essential foundation for the study of standard distributions (binomial, Poisson, normal) later in the course.
\vspace{4pt}
\begin{multicols}{2}\small
\begin{itemize}
\item By the end of this chapter I can define a random variable and distinguish between discrete and continuous random variables.
\item By the end of this chapter I can construct the probability distribution (table of values and probabilities) of a discrete random variable from a described experiment.
\item By the end of this chapter I can verify and use the conditions for a probability mass function, including finding unknown constants.
\item By the end of this chapter I can compute the expectation E(X), the variance Var(X) and the standard deviation of a discrete random variable, and use the linearity rules E(aX+b) and Var(aX+b).
\item By the end of this chapter I can build and interpret the cumulative distribution function F(x) and use it to compute probabilities.
\item By the end of this chapter I can determine the mode and the median of a discrete probability distribution.
\end{itemize}\end{multicols}}

\begin{sommaire}
\begin{multicols}{2}
\begin{enumerate}
\item Random variables and discrete random variables
\item Probability distributions and the probability mass function
\item Expectation, variance and standard deviation
\item The cumulative distribution function, mode and median
\item The distribution of X₁ + X₂
\item Probability generating functions
\item Integration activity
\item Methods and worked examples
\item Exercises
\item Solutions to the exercises
\item Summary sheet
\end{enumerate}
\end{multicols}
\end{sommaire}

\begin{objectifs}
\begin{itemize}
\item By the end of this chapter I can define a random variable and distinguish between discrete and continuous random variables.
\item By the end of this chapter I can construct the probability distribution (table of values and probabilities) of a discrete random variable from a described experiment.
\item By the end of this chapter I can verify and use the conditions for a probability mass function, including finding unknown constants.
\item By the end of this chapter I can compute the expectation E(X), the variance Var(X) and the standard deviation of a discrete random variable, and use the linearity rules E(aX+b) and Var(aX+b).
\item By the end of this chapter I can build and interpret the cumulative distribution function F(x) and use it to compute probabilities.
\item By the end of this chapter I can determine the mode and the median of a discrete probability distribution.
\item By the end of this chapter I can find the distribution of the sum of two independent discrete random variables.
\item By the end of this chapter I can define and use a probability generating function to recover probabilities, expectation and variance.
\end{itemize}
\end{objectifs}

\begin{prerequis}
\begin{itemize}
\item Elementary probability: sample spaces, events, P(A), complementary and combined events (Lower Sixth).
\item Conditional probability and independence of events, P(A∩B) = P(A)P(B) for independent events.
\item Counting techniques: arrangements, permutations and combinations.
\item Summation notation Σ and manipulation of finite sums.
\item Basic differentiation of polynomials and power series (useful for probability generating functions).
\end{itemize}
\end{prerequis}

\newpage

\coursec{Random variables and discrete random variables}

At the Mokolo market in Yaoundé, a trader sells phone credit cards and never knows in advance how many customers will come each day: some days 3, other days 15. Yet, to avoid losses, she must decide every morning how many cards to stock. Stock too few, and she loses sales; stock too many, and her money sleeps on the shelf. How can mathematics help her make a good decision when the outcome is uncertain? The answer lies in modelling the number of customers by a \cle{random variable}, summarising its behaviour by a \cle{probability distribution}, and computing its \cle{expectation}. This chapter gives you the tools that banks, insurers and even the GCE examiners use to reason about uncertainty.

\courssub{Random experiments revisited}

\textbf{Sample space.} Recall from your study of probability that a \cle{random experiment} is a process whose outcome cannot be predicted with certainty, even though we know the list of all possible outcomes. The set of all possible outcomes is called the \cle{sample space}, denoted $\Omega$. For example:

\begin{itemize}
\item Tossing one coin: $\Omega = \{H, T\}$.
\item Tossing two coins: $\Omega = \{HH, HT, TH, TT\}$.
\item Rolling a fair die: $\Omega = \{1, 2, 3, 4, 5, 6\}$.
\item Counting the trader's daily customers: $\Omega = \{0, 1, 2, 3, \ldots\}$.
\end{itemize}

\textbf{The need to attach numbers to outcomes.} Outcomes are not always numbers: $HH$, $HT$ are pairs of letters, not quantities. Yet to compute averages, to compare experiments, to build models for the trader, we \emph{need numbers}. We therefore invent a rule that attaches a real number to every outcome. For the two coins, a natural rule is: ``count the number of heads''. Then
\[ HH \mapsto 2, \qquad HT \mapsto 1, \qquad TH \mapsto 1, \qquad TT \mapsto 0. \]
This rule is exactly what we call a random variable.

\courssub{Definition of a random variable}

\begin{definition}[Random variable]
A \cle{random variable} $X$ is a real-valued function defined on the sample space $\Omega$ of a random experiment:
\[ X : \Omega \longrightarrow \mathbb{R}, \qquad \omega \longmapsto X(\omega). \]
To each outcome $\omega$ of the experiment, $X$ associates exactly one real number $X(\omega)$.
\end{definition}

\begin{attention}[A random variable is neither random nor a variable]
The name is misleading. A random variable is \textbf{not random}: the rule $X$ itself is perfectly fixed. It is \textbf{not a variable} in the algebraic sense (like the $x$ in $2x+1=5$): it is a \emph{function}. The randomness lies in the experiment, not in $X$. Also, do not confuse $X$ (the function, written with a capital letter) with $x$ (one particular value it can take, written with a small letter). Examiners penalise this confusion.
\end{attention}

\textbf{The event $\{X = x\}$.} Since $X$ is a function on $\Omega$, the statement ``$X$ takes the value $x$'' describes a set of outcomes, hence an \emph{event}:
\[ \{X = x\} = \{\omega \in \Omega : X(\omega) = x\}. \]
Its probability is written $P(X = x)$. For the two coins, $\{X = 1\} = \{HT, TH\}$, so $P(X=1) = \tfrac{2}{4} = \tfrac12$.

\begin{popfigure}
\begin{center}
\begin{tikzpicture}[scale=1.0, every node/.style={font=\small}]
% Left: outcomes
\node[font=\bfseries, popblue] at (0,3.4) {Outcomes $\omega \in \Omega$};
\node[draw, popblue, rounded corners, fill=popblueL] (hh) at (0,2.4) {$HH$};
\node[draw, popblue, rounded corners, fill=popblueL] (ht) at (0,1.4) {$HT$};
\node[draw, popblue, rounded corners, fill=popblueL] (th) at (0,0.4) {$TH$};
\node[draw, popblue, rounded corners, fill=popblueL] (tt) at (0,-0.6) {$TT$};
% Right: values
\node[font=\bfseries, poporange] at (5,3.4) {Values $X(\omega) \in \mathbb{R}$};
\node[draw, poporange, rounded corners, fill=white] (v2) at (5,2.4) {$2$};
\node[draw, poporange, rounded corners, fill=white] (v1) at (5,1.0) {$1$};
\node[draw, poporange, rounded corners, fill=white] (v0) at (5,-0.6) {$0$};
% Arrows
\draw[->, thick, popdark] (hh) -- (v2);
\draw[->, thick, popdark] (ht) -- (v1);
\draw[->, thick, popdark] (th) -- (v1);
\draw[->, thick, popdark] (tt) -- (v0);
\node[popmuted] at (2.5,3.0) {$X$ = number of heads};
\end{tikzpicture}
\end{center}
\legende{The random variable $X$ ``number of heads'' maps each outcome of tossing two coins to a real number. Note that two outcomes, $HT$ and $TH$, share the same image $1$.}
\end{popfigure}

\courssub{Discrete and continuous random variables}

\begin{definition}[Discrete random variable]
A random variable $X$ is \cle{discrete} if the set of values it can take is \textbf{finite} or \textbf{countably infinite} (i.e.\ its values can be listed as $x_1, x_2, x_3, \ldots$).
\end{definition}

By contrast, a \cle{continuous} random variable can take any value in an interval of $\mathbb{R}$ (its set of values is uncountable). Continuous variables will be studied later; this chapter deals only with the discrete case.

\begin{exemplebox}[Discrete or continuous?]
\begin{itemize}
\item The number of goals scored by the Indomitable Lions in a match: values $0, 1, 2, \ldots$ — \textbf{discrete}.
\item The number of customers of the Mokolo trader in a day: values $0, 1, 2, \ldots$ — \textbf{discrete}.
\item The score on a fair die: values $\{1,2,3,4,5,6\}$ (finite) — \textbf{discrete}.
\item The height of a randomly chosen Upper Sixth student: any value in an interval, say $[1.40\ \mathrm{m}, 2.00\ \mathrm{m}]$ — \textbf{continuous}.
\item The waiting time for a taxi at the Ngousso junction: any non-negative real number — \textbf{continuous}.
\end{itemize}
\end{exemplebox}

\begin{aretenir}[Finite versus countably infinite]
\begin{itemize}
\item \textbf{Finite:} the die score ($6$ values), the number of heads in two tosses ($3$ values).
\item \textbf{Countably infinite:} the number of tosses of a coin until the first head appears: $1, 2, 3, \ldots$ with no upper bound — the list never ends, but it can still be enumerated.
\end{itemize}
\end{aretenir}

\courssub{Probability distribution and probability mass function}

\begin{definition}[Probability distribution of a discrete random variable]
Let $X$ be a discrete random variable taking values $x_1, x_2, \ldots$ (finite or countably infinite). The \cle{probability distribution} of $X$ is the complete list of pairs $(x_i, p_i)$ where
\[ p_i = P(X = x_i). \]
It is usually presented as a table:
\begin{center}
\begin{tabularx}{\linewidth}{|l|X|X|X|X|}
\hline
\rowcolor{popblueL} $x_i$ & $x_1$ & $x_2$ & $\cdots$ & $x_n$ \\
\hline
$P(X = x_i)$ & $p_1$ & $p_2$ & $\cdots$ & $p_n$ \\
\hline
\end{tabularx}
\end{center}
Since the events $\{X = x_i\}$ are mutually exclusive and exhaustive, we always have
\[ 0 \le p_i \le 1 \quad \text{for every } i, \qquad \text{and} \qquad \sum_{i} p_i = 1. \]
\end{definition}

\begin{definition}[Probability mass function (PMF)]
The function $p_X(x) = P(X = x)$ defined for all real $x$ (with $p_X(x)=0$ if $x$ is not a possible value of $X$) is called the \cle{probability mass function} of $X$. For a discrete random variable, the PMF completely characterises the distribution.
\end{definition}

\begin{aretenir}[Support of a discrete random variable]
The set of values $x$ for which $P(X=x) > 0$ is called the \cle{support} of $X$. For a discrete variable, the support is exactly the set of values that $X$ can actually take.
\end{aretenir}

\begin{methode}[Building the distribution of $X$]
\begin{enumerate}
\item List \textbf{all} the outcomes of $\Omega$ and their probabilities.
\item Compute the value $X(\omega)$ for \textbf{each} outcome.
\item \textbf{Group} the outcomes giving the same value $x_i$: they form the event $\{X = x_i\}$.
\item \textbf{Add} the probabilities of the outcomes inside each group to get $p_i = P(X = x_i)$.
\item \textbf{Check} that $\sum p_i = 1$. If not, an outcome has been forgotten or counted twice.
\end{enumerate}
\end{methode}

\begin{exoresolu}[Number of heads in two tosses of a fair coin]
Two fair coins are tossed. Let $X$ be the number of heads obtained. Determine the probability distribution of $X$.
\tcblower
\textbf{Solution.} The sample space is $\Omega = \{HH, HT, TH, TT\}$, each outcome having probability $\tfrac14$ since the coins are fair and independent.

Applying the method: $X(HH) = 2$, $X(HT) = 1$, $X(TH) = 1$, $X(TT) = 0$. Grouping:
\[ P(X=0) = P(\{TT\}) = \tfrac14, \quad P(X=1) = P(\{HT, TH\}) = \tfrac24 = \tfrac12, \quad P(X=2) = P(\{HH\}) = \tfrac14. \]
\begin{center}
\begin{tabularx}{\linewidth}{|l|X|X|X|}
\hline
\rowcolor{popblueL} $x_i$ & $0$ & $1$ & $2$ \\
\hline
$P(X = x_i)$ & $\dfrac14$ & $\dfrac12$ & $\dfrac14$ \\
\hline
\end{tabularx}
\end{center}
Check: $\tfrac14 + \tfrac12 + \tfrac14 = 1$. \checkmark
\end{exoresolu}

\begin{exoresolu}[Score on a fair die]
A fair die is rolled and $Y$ is the score shown. Give the distribution of $Y$.
\tcblower
\textbf{Solution.} Here each outcome of $\Omega = \{1,2,3,4,5,6\}$ is already a number, so $Y(\omega) = \omega$. Each face has probability $\tfrac16$:
\begin{center}
\begin{tabularx}{\linewidth}{|l|X|X|X|X|X|X|}
\hline
\rowcolor{popblueL} $y_i$ & $1$ & $2$ & $3$ & $4$ & $5$ & $6$ \\
\hline
$P(Y = y_i)$ & $\dfrac16$ & $\dfrac16$ & $\dfrac16$ & $\dfrac16$ & $\dfrac16$ & $\dfrac16$ \\
\hline
\end{tabularx}
\end{center}
This is the \emph{uniform} distribution on $\{1, \ldots, 6\}$: all values are equally likely.
\end{exoresolu}

\begin{exoresolu}[Defective items in a sample]
A box contains $5$ phone chargers, of which $2$ are defective. A quality controller picks $2$ chargers at random, without replacement. Let $Z$ be the number of defective chargers in the sample. Find the distribution of $Z$.
\tcblower
\textbf{Solution.} The possible values of $Z$ are $0, 1, 2$. The number of ways to choose $2$ chargers out of $5$ is $\binom{5}{2} = 10$, all equally likely.
\begin{align*}
P(Z = 0) &= \frac{\binom{2}{0}\binom{3}{2}}{10} = \frac{3}{10}, \\
P(Z = 1) &= \frac{\binom{2}{1}\binom{3}{1}}{10} = \frac{6}{10}, \\
P(Z = 2) &= \frac{\binom{2}{2}\binom{3}{0}}{10} = \frac{1}{10}.
\end{align*}
\begin{center}
\begin{tabularx}{\linewidth}{|l|X|X|X|}
\hline
\rowcolor{popblueL} $z_i$ & $0$ & $1$ & $2$ \\
\hline
$P(Z = z_i)$ & $\dfrac{3}{10}$ & $\dfrac{6}{10}$ & $\dfrac{1}{10}$ \\
\hline
\end{tabularx}
\end{center}
Check: $\tfrac{3}{10} + \tfrac{6}{10} + \tfrac{1}{10} = 1$. \checkmark
\end{exoresolu}

\courssub{Representing a distribution by a bar chart}

A probability distribution is conveniently pictured by a \cle{bar chart}: the values $x_i$ are placed on the horizontal axis and the heights of the bars equal the probabilities $p_i$. The total of the heights is always $1$. (For a discrete variable, a histogram with bars of equal width is also acceptable, but the bar chart emphasises the gaps between possible values.)

\begin{popfigure}
\begin{center}
\begin{tikzpicture}
\begin{axis}[
    ybar, bar width=22pt,
    width=9cm, height=6cm,
    xlabel={Number of heads $x$},
    ylabel={$P(X = x)$},
    xtick={0,1,2},
    ymin=0, ymax=0.65,
    ytick={0,0.25,0.5},
    yticklabels={$0$,$0.25$,$0.5$},
    axis lines=left,
    every axis plot/.append style={fill=popblue, draw=popdark},
]
\addplot coordinates {(0,0.25) (1,0.5) (2,0.25)};
\end{axis}
\end{tikzpicture}
\end{center}
\legende{Bar chart of the probability distribution of $X$, the number of heads when two fair coins are tossed. The symmetry reflects the fairness of the coins.}
\end{popfigure}

\begin{attention}[Common mistakes]
\begin{itemize}
\item Forgetting to \textbf{add} probabilities when several outcomes give the same value of $X$ (e.g.\ writing $P(X=1) = \tfrac14$ instead of $\tfrac12$ for the two coins).
\item Giving a table whose probabilities do not sum to $1$ — always check!
\item Confusing $P(X = x)$ with $P(X \le x)$: the latter is the \emph{cumulative} probability, studied later in this chapter.
\item Believing a discrete variable can take ``any'' value: between two consecutive values of a discrete variable, no value is possible (a die never shows $3.5$).
\item Writing $P(X=x)$ for a value $x$ outside the support; remember $P(X=x)=0$ for such $x$.
\end{itemize}
\end{attention}

\begin{savaistu}[Why insurers love random variables]
When an insurance company in Douala fixes the price of a policy, it models the number of accidents a driver will have in a year as a discrete random variable, estimates its distribution from past data, and computes its expectation. The premium must exceed this expected cost, otherwise the company loses money on average. The same mathematics powers mobile-money risk management and the design of fair games of chance.
\end{savaistu}

\begin{exercice}[Practice]
\begin{enumerate}
\item Three fair coins are tossed and $X$ counts the number of heads. Construct the probability distribution of $X$ and draw its bar chart.
\item A bag contains $4$ red and $6$ black balls. Two balls are drawn at random without replacement, and $Y$ counts the red balls drawn. Find the distribution of $Y$.
\item State, with a reason, whether each variable is discrete or continuous: (a) the number of students absent in your class tomorrow; (b) the rainfall in Bamenda next Saturday, in millimetres.
\end{enumerate}
\end{exercice}

\begin{corrige}[Exercise 1]
\textbf{1.} $\Omega$ has $8$ equally likely outcomes. Grouping by number of heads: $P(X=0) = \tfrac18$, $P(X=1) = \tfrac38$, $P(X=2) = \tfrac38$, $P(X=3) = \tfrac18$. Sum $= 1$. \checkmark\ The bar chart has bars of heights $\tfrac18, \tfrac38, \tfrac38, \tfrac18$ at $x = 0,1,2,3$.

\textbf{2.} $\binom{10}{2} = 45$ equally likely pairs. $P(Y=0) = \tfrac{\binom{6}{2}}{45} = \tfrac{15}{45} = \tfrac13$; $P(Y=1) = \tfrac{\binom{4}{1}\binom{6}{1}}{45} = \tfrac{24}{45} = \tfrac{8}{15}$; $P(Y=2) = \tfrac{\binom{4}{2}}{45} = \tfrac{6}{45} = \tfrac{2}{15}$. Sum $= 1$. \checkmark

\textbf{3.} (a) Discrete: the values are whole numbers $0, 1, 2, \ldots$. (b) Continuous: rainfall can take any non-negative real value within an interval.
\end{corrige}

\coursec{Probability distributions and the probability mass function}

In the previous section we met the idea of a \cle{random variable}: a numerical quantity $X$ whose value depends on the outcome of a random experiment, and we saw that a \cle{discrete random variable} takes only isolated values, which we can list as $x_1, x_2, x_3, \dots$ Knowing the \emph{possible values} of $X$, however, is only half of the story. A fair die and a loaded die both take the values $1, 2, 3, 4, 5, 6$; what distinguishes them is \emph{how likely} each value is. The complete description of a discrete random variable is therefore a list of its values \emph{together with} the probability attached to each value. That complete description is what we now call a \cle{probability distribution}, and the function that delivers the probabilities is the \cle{probability mass function}. Mastering this object is essential: nearly every GCE question on discrete random variables begins by giving you a distribution and asking you to work with it.

\courssub{The probability mass function}

\begin{definition}[Probability mass function]
Let $X$ be a discrete random variable taking values $x_1, x_2, \dots, x_n$ (or an infinite list $x_1, x_2, \dots$). The \cle{probability mass function} (p.m.f.) of $X$ is the function $p$ defined by
\[ p(x) = P(X = x), \]
that is, $p(x_i)$ is the probability that $X$ takes the value $x_i$. We also write $p_i = P(X = x_i)$. For any $x$ which is not a possible value of $X$, $p(x) = 0$.
\end{definition}

The word \emph{mass} is suggestive: imagine a total mass of $1$ kilogram of probability, chopped into lumps and placed at the points $x_1, x_2, \dots$ on a number line. The lump at $x_i$ has mass $p_i$. This picture immediately tells us what conditions a genuine p.m.f. must satisfy.

\begin{propriete}[Defining conditions of a p.m.f.]
A function $p$ is the probability mass function of a discrete random variable if and only if
\begin{enumerate}
\item $p(x_i) \geq 0$ for every $i$ (in fact $0 \leq p(x_i) \leq 1$);
\item $\displaystyle\sum_{i} p(x_i) = 1$, the sum being taken over \emph{all} possible values of $X$.
\end{enumerate}
\end{propriete}

\textbf{Justification.} Condition 1 holds because every probability lies between $0$ and $1$. For condition 2, the events $\{X = x_1\}, \{X = x_2\}, \dots$ are \emph{mutually exclusive} ($X$ cannot take two values at once) and \emph{exhaustive} ($X$ must take \emph{some} value). Hence, by the addition law for mutually exclusive events,
\[ \sum_i p(x_i) = P(X = x_1) + P(X = x_2) + \cdots = P(X \text{ takes one of its values}) = 1. \]
This simple argument is worth being able to reproduce: it is the reason behind the single most used technique in this topic.

\begin{aretenir}[Presenting a distribution]
A probability distribution is usually presented as a \textbf{complete table}:
\begin{center}
\begin{tabular}{|c|ccccc|}
\hline
\rowcolor{popblueL} $x$ & $x_1$ & $x_2$ & $x_3$ & $\cdots$ & $x_n$ \\ \hline
$P(X = x)$ & $p_1$ & $p_2$ & $p_3$ & $\cdots$ & $p_n$ \\ \hline
\end{tabular}
\end{center}
The table must show \emph{every} possible value of $X$ and the probabilities must sum to $1$. Always check this sum before using a distribution.
\end{aretenir}

\begin{savaistu}[Terminology note]
The syllabus mentions ``probability density or mass function''. For \emph{discrete} random variables the correct term is \textbf{probability mass function} (p.m.f.). The term ``probability density function'' (p.d.f.) is reserved for \emph{continuous} random variables, which you will meet later. In the GCE examination, always write ``probability mass function'' or ``p.m.f.'' when dealing with discrete variables.
\end{savaistu}

\courssub{Finding an unknown constant: the classic GCE question}

Very often the GCE gives a distribution containing an unknown constant $k$ and asks you to find it. The key is always the same: the probabilities must add up to $1$.

\begin{methode}[Finding an unknown constant in a distribution]
\begin{enumerate}
\item Write down the equation $\sum_i p(x_i) = 1$ using the given formula or table.
\item Solve the resulting equation for $k$.
\item \textbf{Check} that with this value of $k$, every probability satisfies $0 \leq p(x_i) \leq 1$. (If an equation gives two solutions, one of them usually fails this check and must be rejected.)
\item Rewrite the distribution as a complete numerical table.
\end{enumerate}
\end{methode}

\begin{exoresolu}[A distribution given by a formula]
A discrete random variable $X$ has probability mass function
\[ P(X = x) = kx, \qquad x = 1, 2, 3, 4. \]
Find the value of $k$ and write out the probability distribution of $X$ in a table.
\tcblower
\textbf{Solution.} The probabilities must sum to $1$:
\[ \sum_{x=1}^{4} kx = k(1) + k(2) + k(3) + k(4) = 10k. \]
Setting $10k = 1$ gives $k = \dfrac{1}{10} = 0.1$.
\par
\textbf{Check:} the probabilities are $0.1, 0.2, 0.3, 0.4$, all between $0$ and $1$. \checkmark
\par
The distribution is:
\begin{center}
\begin{tabular}{|c|cccc|}
\hline
\rowcolor{popblueL} $x$ & $1$ & $2$ & $3$ & $4$ \\ \hline
$P(X = x)$ & $0.1$ & $0.2$ & $0.3$ & $0.4$ \\ \hline
\end{tabular}
\end{center}
\end{exoresolu}

\begin{exoresolu}[A quadratic equation for $k$]
A random variable $X$ takes the values $0, 1, 2$ with
\[ P(X = 0) = k, \qquad P(X = 1) = 2k, \qquad P(X = 2) = 4k^2 + k. \]
Find the value of $k$.
\tcblower
\textbf{Solution.} Summing the probabilities:
\[ k + 2k + 4k^2 + k = 4k^2 + 4k = 1
\quad\Longrightarrow\quad 4k^2 + 4k - 1 = 0. \]
By the quadratic formula,
\[ k = \frac{-4 \pm \sqrt{16 + 16}}{8} = \frac{-4 \pm \sqrt{32}}{8} = \frac{-1 \pm \sqrt{2}}{2}. \]
Now $\dfrac{-1 - \sqrt{2}}{2} \approx -1.207 < 0$, which would make $P(X=0)$ negative: rejected. Hence
\[ k = \frac{\sqrt{2} - 1}{2} \approx 0.207. \]
\textbf{Check:} $P(X=0) \approx 0.207$, $P(X=1) \approx 0.414$, $P(X=2) \approx 4(0.0429) + 0.207 \approx 0.379$; all in $[0,1]$ and $0.207 + 0.414 + 0.379 = 1$. \checkmark
\end{exoresolu}

\begin{exoresolu}[An infinite distribution: geometric-type]
A discrete random variable $X$ takes values $1, 2, 3, \dots$ with
\[ P(X = x) = k \left(\frac{2}{3}\right)^x. \]
Find the value of $k$.
\tcblower
\textbf{Solution.} The sum of probabilities is an infinite geometric series:
\[ \sum_{x=1}^{\infty} k \left(\frac{2}{3}\right)^x = k \cdot \frac{\frac{2}{3}}{1 - \frac{2}{3}} = k \cdot \frac{2/3}{1/3} = 2k. \]
Setting $2k = 1$ gives $k = \frac{1}{2}$.
\par
\textbf{Check:} $P(X=x) = \frac{1}{2}\left(\frac{2}{3}\right)^x > 0$ for all $x \ge 1$, and the sum is $1$. \checkmark
\end{exoresolu}

\begin{attention}[Do not forget the check]
When the equation for $k$ is quadratic, one root is almost always invalid because it produces a negative probability or a probability greater than $1$. Examiners award a mark for explicitly rejecting the invalid root. Also, never assume $k$ itself is a probability: in $P(X = x) = kx$, the constant $k$ is just a scaling factor.
\end{attention}

\courssub{Computing probabilities of events from the distribution}

Once the distribution is known, the probability of any event concerning $X$ is found by \textbf{adding the masses} at the values that satisfy the event, because the events $\{X = x_i\}$ are mutually exclusive:
\[ P(a \leq X \leq b) = \sum_{a \leq x_i \leq b} p(x_i). \]

\begin{exemplebox}[Compound events]
Using the distribution $P(X = x) = 0.1x$ for $x = 1, 2, 3, 4$ found above:
\begin{itemize}
\item $P(X \leq 2) = P(X=1) + P(X=2) = 0.1 + 0.2 = 0.3$;
\item $P(X > 2) = P(X=3) + P(X=4) = 0.3 + 0.4 = 0.7$ (note: $= 1 - P(X \leq 2)$);
\item $P(2 \leq X \leq 3) = 0.2 + 0.3 = 0.5$;
\item $P(X \text{ is even}) = P(X=2) + P(X=4) = 0.2 + 0.4 = 0.6$.
\end{itemize}
\end{exemplebox}

\begin{attention}[Strict versus weak inequalities]
For a \emph{discrete} variable, $P(X > 2)$ and $P(X \geq 2)$ are \textbf{different}: the first excludes $x = 2$, the second includes it. Read the inequality sign with great care — this is one of the most frequent sources of lost marks. (For continuous variables, met later, the two coincide.)
\end{attention}

\courssub{The discrete uniform distribution}

Our first \emph{standard family} of distributions is the simplest imaginable: all values equally likely.

\begin{definition}[Discrete uniform distribution]
A discrete random variable $X$ has the \cle{discrete uniform distribution} on $\{1, 2, \dots, n\}$ if
\[ P(X = x) = \frac{1}{n} \quad \text{for } x = 1, 2, \dots, n. \]
This is a valid p.m.f. since each probability is positive and $n \times \dfrac{1}{n} = 1$.
\end{definition}

The score on a fair die is uniform on $\{1, \dots, 6\}$; the number drawn in a fair tombola with tickets $1$ to $n$ is uniform on $\{1, \dots, n\}$. Uniformity is a \emph{property of the experiment}, not of the set of values: the loaded die below has the same values but a very different distribution.

\begin{popfigure}
\begin{center}
\begin{tikzpicture}
\begin{axis}[
  ybar, bar width=0.35cm,
  width=0.48\linewidth, height=5.2cm,
  xlabel={$x$}, ylabel={$P(X=x)$},
  xtick={1,2,3,4,5,6},
  ymin=0, ymax=0.35,
  title={Uniform (fair die)},
  title style={font=\small},
  label style={font=\small},
  tick label style={font=\small},
  ymajorgrids, grid style={popline},
]
\addplot[fill=popblue, draw=popdark] coordinates {(1,0.1667) (2,0.1667) (3,0.1667) (4,0.1667) (5,0.1667) (6,0.1667)};
\end{axis}
\end{tikzpicture}
\hfill
\begin{tikzpicture}
\begin{axis}[
  ybar, bar width=0.35cm,
  width=0.48\linewidth, height=5.2cm,
  xlabel={$x$}, ylabel={$P(X=x)$},
  xtick={1,2,3,4,5,6},
  ymin=0, ymax=0.35,
  title={Non-uniform (loaded die)},
  title style={font=\small},
  label style={font=\small},
  tick label style={font=\small},
  ymajorgrids, grid style={popline},
]
\addplot[fill=poporange, draw=popdark] coordinates {(1,0.05) (2,0.10) (3,0.15) (4,0.20) (5,0.20) (6,0.30)};
\end{axis}
\end{tikzpicture}
\end{center}
\legende{Two distributions on $\{1,2,3,4,5,6\}$: the uniform distribution (every bar has height $1/6$) and a non-uniform distribution. In both cases the heights sum to $1$.}
\end{popfigure}

\courssub{Functions of a random variable: the distribution of $Y = g(X)$}

Suppose $X$ is the score on a die and you win $Y = 2X + 1$ francs, or you square the outcome, $Y = X^2$. Then $Y$ is itself a discrete random variable, and we can derive its distribution from that of $X$. The principle is:
\[ P(Y = y) = \sum_{x \,:\, g(x) = y} P(X = x), \]
that is, \textbf{add the probabilities of all values of $X$ that are sent to the same value $y$}. When $g$ is one-to-one on the values of $X$ (like $g(x) = 2x + 1$), each $y$ comes from a single $x$ and the probabilities are simply carried over. When $g$ is \emph{not} one-to-one (like $g(x) = x^2$ when $X$ takes negative values), several values of $X$ may collapse onto one value of $Y$, and their probabilities must be added.

\begin{exoresolu}[Distribution of $Y = X^2$]
Let $X$ have the distribution
\begin{center}
\begin{tabular}{|c|ccccc|}
\hline
\rowcolor{popblueL} $x$ & $-2$ & $-1$ & $0$ & $1$ & $2$ \\ \hline
$P(X = x)$ & $0.1$ & $0.2$ & $0.3$ & $0.25$ & $0.15$ \\ \hline
\end{tabular}
\end{center}
Find the probability distribution of $Y = X^2$.
\tcblower
\textbf{Solution.} The possible values of $Y$ are $0, 1, 4$. Grouping the values of $X$ with the same square:
\begin{align*}
P(Y = 0) &= P(X = 0) = 0.3,\\
P(Y = 1) &= P(X = -1) + P(X = 1) = 0.2 + 0.25 = 0.45,\\
P(Y = 4) &= P(X = -2) + P(X = 2) = 0.1 + 0.15 = 0.25.
\end{align*}
Hence the distribution of $Y$ is
\begin{center}
\begin{tabular}{|c|ccc|}
\hline
\rowcolor{popblueL} $y$ & $0$ & $1$ & $4$ \\ \hline
$P(Y = y)$ & $0.3$ & $0.45$ & $0.25$ \\ \hline
\end{tabular}
\end{center}
\textbf{Check:} $0.3 + 0.45 + 0.25 = 1$. \checkmark
\end{exoresolu}

\begin{popfigure}
\begin{center}
\begin{tikzpicture}[font=\small]
% X table
\node[font=\small\bfseries] at (1.5,2.6) {Distribution of $X$};
\draw[popdark] (0,1.6) rectangle (3,2.2);
\draw[popdark] (0,1.0) rectangle (3,1.6);
\foreach \i/\v in {0/-2, 0.6/-1, 1.2/0, 1.8/1, 2.4/2} {
  \draw[popdark] (\i,1.0) -- (\i,2.2);
  \node at (\i+0.3,1.9) {$\v$};
}
\foreach \i/\v in {0/0.1, 0.6/0.2, 1.2/0.3, 1.8/0.25, 2.4/0.15} {
  \node at (\i+0.3,1.3) {$\v$};
}
% arrows
\draw[->, thick, popblue] (0.3,1.0) .. controls (1.2,0.2) .. (5.1,1.15);
\draw[->, thick, popblue] (0.9,1.0) .. controls (1.8,0.1) .. (5.7,1.05);
\draw[->, thick, popblue] (1.5,1.0) .. controls (2.6,0.0) .. (4.5,1.05);
\draw[->, thick, poporange] (2.1,1.0) .. controls (3.2,0.1) .. (5.85,1.15);
\draw[->, thick, poporange] (2.7,1.0) .. controls (3.6,0.3) .. (6.45,1.2);
\node[font=\small, align=center] at (3.2,0.55) {square: $y = x^2$};
% Y table
\node[font=\small\bfseries] at (5.4,2.6) {Distribution of $Y = X^2$};
\draw[popdark] (4.2,1.6) rectangle (6.6,2.2);
\draw[popdark] (4.2,1.0) rectangle (6.6,1.6);
\foreach \i/\v in {4.2/0, 5.0/1, 5.8/4} {
  \draw[popdark] (\i,1.0) -- (\i,2.2);
  \node at (\i+0.4,1.9) {$\v$};
}
\node at (4.6,1.3) {$0.3$};
\node at (5.4,1.3) {$0.45$};
\node at (6.2,1.3) {$0.25$};
\node[font=\footnotesize, text=popmuted, align=center] at (5.4,3.1) {$0.2+0.25$ at $y=1$,\ $0.1+0.15$ at $y=4$};
\end{tikzpicture}
\end{center}
\legende{Transforming the distribution of $X$ into that of $Y = X^2$: values of $X$ with the same square merge, and their probabilities are added.}
\end{popfigure}

\begin{attention}[Merging values under $Y = g(X)$]
When finding the distribution of $Y = X^2$, values of $X$ with the same square (e.g. $-2$ and $2$) must have their probabilities \textbf{added}. A very common error is to list $y = 4$ twice, or to copy only one of the two probabilities. Similarly, for $Y = |X|$ or $Y = (X-1)^2$, always ask: \emph{which values of $X$ give the same $Y$?}
\end{attention}

\begin{exemplebox}[A one-to-one transformation: $Y = 2X + 1$]
If $X$ is uniform on $\{1, 2, 3\}$ (each probability $\tfrac13$), then $Y = 2X + 1$ takes the values $3, 5, 7$, each with probability $\tfrac13$: the distribution of $Y$ is uniform on $\{3, 5, 7\}$. No merging occurs because $g(x) = 2x + 1$ is strictly increasing.
\end{exemplebox}

\begin{exemplebox}[A Cameroonian context: Loterie Nationale]
In a simplified version of the \emph{Loterie Nationale}, a ticket bears a number $X$ from $1$ to $5$ with probabilities:
\begin{center}
\begin{tabular}{|c|ccccc|}
\hline
\rowcolor{popblueL} $x$ & $1$ & $2$ & $3$ & $4$ & $5$ \\ \hline
$P(X = x)$ & $0.1$ & $0.2$ & $0.3$ & $0.25$ & $0.15$ \\ \hline
\end{tabular}
\end{center}
The prize is $Y = 1000X - 2000$ FCFA. Find the distribution of $Y$.
\textbf{Solution.} $g(x) = 1000x - 2000$ is one-to-one. The values of $Y$ are $-1000, 0, 1000, 2000, 3000$ with the same probabilities $0.1, 0.2, 0.3, 0.25, 0.15$ respectively.
\end{exemplebox}

\begin{exercice}[Practice]
\begin{enumerate}
\item A random variable $X$ has p.m.f. $P(X = x) = kx^2$ for $x = 1, 2, 3$. Find $k$, tabulate the distribution, and compute $P(X \geq 2)$ and $P(X \text{ is odd})$.
\item The distribution of $X$ is given by $P(X = 0) = 0.5k$, $P(X = 1) = k$, $P(X = 2) = k^2$, $P(X = 3) = 0.1$. Find $k$.
\item $X$ is uniform on $\{-2, -1, 0, 1, 2\}$. Find the distributions of $Y = X^2$ and of $Z = 3X - 1$.
\item A discrete random variable $X$ takes values $0, 1, 2, 3$ with $P(X = x) = k(4-x)$. Find $k$ and the distribution of $Y = (X-1)^2$.
\end{enumerate}
\end{exercice}

\begin{corrige}[Exercise 1]
\textbf{1.} $\sum kx^2 = k(1 + 4 + 9) = 14k = 1$, so $k = \tfrac{1}{14}$. The probabilities are $\tfrac{1}{14}, \tfrac{4}{14}, \tfrac{9}{14}$ for $x = 1, 2, 3$. Then $P(X \geq 2) = \tfrac{4}{14} + \tfrac{9}{14} = \tfrac{13}{14}$ and $P(X \text{ odd}) = P(X=1) + P(X=3) = \tfrac{1}{14} + \tfrac{9}{14} = \tfrac{10}{14} = \tfrac{5}{7}$.
\par
\textbf{2.} $0.5k + k + k^2 + 0.1 = 1 \Rightarrow k^2 + 1.5k - 0.9 = 0 \Rightarrow 10k^2 + 15k - 9 = 0$, giving $k = \dfrac{-15 \pm \sqrt{225 + 360}}{20} = \dfrac{-15 \pm \sqrt{585}}{20}$. Only the positive root is valid: $k = \dfrac{-15 + \sqrt{585}}{20} \approx 0.459$. (Check: probabilities $\approx 0.230,\ 0.459,\ 0.211,\ 0.1$, sum $= 1$.)
\par
\textbf{3.} Each value of $X$ has probability $\tfrac15 = 0.2$. For $Y = X^2$: $P(Y=0) = 0.2$, $P(Y=1) = 0.2 + 0.2 = 0.4$, $P(Y=4) = 0.2 + 0.2 = 0.4$. For $Z = 3X - 1$: the values are $-7, -4, -1, 2, 5$, each with probability $0.2$ (one-to-one map, no merging).
\par
\textbf{4.} $\sum_{x=0}^3 k(4-x) = k(4+3+2+1) = 10k = 1 \Rightarrow k = 0.1$. Distribution of $X$:
\begin{center}
\begin{tabular}{|c|cccc|}
\hline
\rowcolor{popblueL} $x$ & $0$ & $1$ & $2$ & $3$ \\ \hline
$P(X=x)$ & $0.4$ & $0.3$ & $0.2$ & $0.1$ \\ \hline
\end{tabular}
\end{center}
For $Y = (X-1)^2$: $X=0 \to Y=1$, $X=1 \to Y=0$, $X=2 \to Y=1$, $X=3 \to Y=4$.
$P(Y=0) = P(X=1) = 0.3$; $P(Y=1) = P(X=0)+P(X=2) = 0.4+0.2 = 0.6$; $P(Y=4) = P(X=3) = 0.1$.
Distribution of $Y$:
\begin{center}
\begin{tabular}{|c|ccc|}
\hline
\rowcolor{popblueL} $y$ & $0$ & $1$ & $4$ \\ \hline
$P(Y=y)$ & $0.3$ & $0.6$ & $0.1$ \\ \hline
\end{tabular}
\end{center}
\end{corrige}

\begin{attention}[Exam tip — GCE]
Always present the final distribution as a \textbf{complete table} showing every value of the variable and its probability. Marks are awarded for the full set of values \emph{and} the full set of probabilities; a formula alone, or a table missing a value, loses marks. Finish with the check $\sum p_i = 1$ — it takes one line and often earns a method mark.
\end{attention}

\begin{aretenir}[Summary of the section]
\begin{itemize}
\item The \cle{probability mass function} of a discrete random variable is $p(x) = P(X = x)$.
\item It must satisfy $p(x_i) \geq 0$ for all $i$ and $\sum_i p(x_i) = 1$; the second condition is the tool for finding unknown constants.
\item Probabilities of events are sums of masses: $P(a \leq X \leq b) = \sum_{a \leq x_i \leq b} p(x_i)$; watch strict versus weak inequalities.
\item The \cle{discrete uniform distribution} on $\{1, \dots, n\}$ has $P(X = x) = \tfrac{1}{n}$.
\item For $Y = g(X)$, $P(Y = y) = \sum_{g(x) = y} P(X = x)$: add probabilities of values of $X$ that merge.
\end{itemize}
\end{aretenir}

With the distribution of a discrete random variable now firmly in hand, we can ask the next natural question: how do we summarise a whole distribution by one or two numbers? That is the role of the \emph{expectation}, \emph{variance} and \emph{standard deviation}, the subject of the next section.

\coursec{Expectation, variance and standard deviation}

In the previous section we learned how to describe a discrete random variable by its probability mass function (p.m.f.) or its probability distribution table. That table gives a complete picture of the possible values and their likelihoods. However, in practice we often need a single number that summarises the \emph{centre} of the distribution (the ``average'' outcome we would expect in the long run) and another number that measures how much the values \emph{spread} around that centre. These are the \textbf{expectation} (or mean) and the \textbf{variance} (with its square root, the standard deviation). They are the cornerstones of statistical decision‑making — from evaluating a game of chance to comparing investment portfolios.

\courssub{The expectation (mean) of a discrete random variable}

Imagine a market trader in Douala who sells roasted plantains (\emph{beignets de plantain}). Over many days she records her daily profit (in CFA francs). Some days she makes 5000, some days 10000, some days 2000, and occasionally she makes a loss. If we build the probability distribution of her daily profit $X$ from a long history, the \textbf{expectation} $E(X)$ is exactly the long‑run average profit per day. It is the value she can ``expect'' on a typical day, not because it will occur every day, but because the fluctuations balance out over time.

\begin{definition}[Expectation of a discrete random variable]
Let $X$ be a discrete random variable taking values $x_1, x_2, \dots, x_n$ with probabilities $p_1, p_2, \dots, p_n$ (where $p_i = P(X=x_i)$ and $\sum_{i=1}^n p_i = 1$). The \textbf{expectation} (or \textbf{mean}) of $X$, denoted $E(X)$ or $\mu$, is defined by
\[
E(X) = \sum_{i=1}^{n} x_i\,p_i = \sum_{i=1}^{n} x_i\,P(X=x_i).
\]
It is the \textbf{weighted average} of the possible values, the weights being the probabilities.
\end{definition}

\begin{exemplebox}[Expectation of a fair die score]
Let $X$ be the score when a fair six‑sided die is rolled. The distribution is
\[
\begin{array}{c|cccccc}
x & 1 & 2 & 3 & 4 & 5 & 6 \\ \hline
P(X=x) & \frac16 & \frac16 & \frac16 & \frac16 & \frac16 & \frac16
\end{array}
\]
Then
\[
E(X) = \sum_{x=1}^{6} x\cdot\frac16 = \frac{1+2+3+4+5+6}{6} = \frac{21}{6} = 3.5.
\]
Notice that $3.5$ is not a possible outcome of a single roll; it is the theoretical centre of the distribution.
\end{exemplebox}

\begin{exemplebox}[Expectation of a discrete uniform distribution]
If $X$ takes values $1,2,\dots,n$ each with probability $\frac1n$, then
\[
E(X) = \frac{1}{n}\sum_{x=1}^n x = \frac{n(n+1)}{2n} = \frac{n+1}{2}.
\]
For a fair die ($n=6$), this gives $\frac{7}{2}=3.5$ as before.
\end{exemplebox}

\textbf{Expectation of a function of $X$.} Often we need the average of a transformed variable, e.g. $X^2$ or $aX+b$. The same weighting principle applies:

\begin{propriete}[Expectation of a function]
For any real‑valued function $g$,
\[
E\bigl(g(X)\bigr) = \sum_{i} g(x_i)\,P(X=x_i).
\]
In particular,
\[
E(X^2) = \sum_{i} x_i^2\,P(X=x_i).
\]
\end{propriete}

\begin{exemplebox}[Computing $E(X^2)$ for the die]
Using the same die,
\[
E(X^2) = \frac{1^2+2^2+3^2+4^2+5^2+6^2}{6} = \frac{91}{6} \approx 15.1667.
\]
\end{exemplebox}

\textbf{Linearity of expectation.} One of the most powerful properties of expectation is that it is a \textbf{linear operator}. This means we can pull constants out and split sums — a fact that simplifies countless calculations.

\begin{propriete}[Linearity of expectation]
For any discrete random variable $X$ and any constants $a,b \in \mathbb{R}$,
\[
E(aX + b) = a\,E(X) + b.
\]
More generally, for any two random variables $X$ and $Y$ (discrete or not, independent or not),
\[
E(X+Y) = E(X) + E(Y).
\]
\end{propriete}

\begin{propriete}[Proof of linearity]
Let $X$ take values $x_i$ with probabilities $p_i$. Then $aX+b$ takes values $ax_i+b$ with the same probabilities $p_i$. Hence
\begin{align*}
E(aX+b) &= \sum_i (ax_i+b)\,p_i = a\sum_i x_i p_i + b\sum_i p_i \\
&= a\,E(X) + b\cdot 1 = a\,E(X) + b.
\end{align*}
The additivity $E(X+Y)=E(X)+E(Y)$ follows similarly by summing over the joint distribution (no independence required).
\end{propriete}

\begin{exemplebox}[Linear transformation of the die score]
Suppose a game pays $Y = 2X + 100$ CFA francs, where $X$ is the die score. Then
\[
E(Y) = 2E(X) + 100 = 2\times 3.5 + 100 = 107 \text{ CFA francs}.
\]
We did not need to build the full distribution of $Y$; linearity did the work instantly.
\end{exemplebox}

\begin{exemplebox}[Expectation of a non‑linear function]
Let $X$ be the score of a fair die. Compute $E(1/X)$.
\[
E(1/X) = \sum_{x=1}^6 \frac{1}{x}\cdot\frac16 = \frac16\left(1+\frac12+\frac13+\frac14+\frac15+\frac16\right) = \frac{49}{120} \approx 0.4083.
\]
Note that $E(1/X) \neq 1/E(X) = 2/7 \approx 0.2857$. Expectation does not commute with non‑linear functions.
\end{exemplebox}

\courssub{Variance and standard deviation}

Expectation tells us \emph{where} the distribution is centred, but not how \emph{spread out} it is. Two distributions can have the same mean but very different variability. Consider two games:
\begin{itemize}
\item Game A: you win 1000 CFA with probability 0.5, lose 1000 CFA with probability 0.5. Mean = 0.
\item Game B: you win 1,000,000 CFA with probability 0.5, lose 1,000,000 CFA with probability 0.5. Mean = 0.
\end{itemize}
Both have $E(X)=0$, but Game B is enormously riskier. The \textbf{variance} quantifies this spread.

\begin{definition}[Variance and standard deviation]
The \textbf{variance} of a discrete random variable $X$, denoted $\operatorname{Var}(X)$ or $\sigma^2$, is the expected squared deviation from the mean:
\[
\operatorname{Var}(X) = E\bigl[(X - \mu)^2\bigr] = \sum_i (x_i - \mu)^2\,p_i,
\]
where $\mu = E(X)$. An equivalent and often more convenient formula is
\[
\operatorname{Var}(X) = E(X^2) - \bigl[E(X)\bigr]^2.
\]
The \textbf{standard deviation} is $\sigma = \sqrt{\operatorname{Var}(X)}$. It has the same units as $X$ and measures the typical distance of outcomes from the mean.
\end{definition}

\begin{propriete}[Proof of equivalence of variance formulas]
\begin{align*}
E\bigl[(X-\mu)^2\bigr] &= E(X^2 - 2\mu X + \mu^2) \\
&= E(X^2) - 2\mu E(X) + \mu^2 \\
&= E(X^2) - 2\mu^2 + \mu^2 = E(X^2) - \mu^2.
\end{align*}
\end{propriete}

\begin{exemplebox}[Variance of a fair die score]
We already have $E(X)=3.5$ and $E(X^2)=91/6$. Hence
\[
\operatorname{Var}(X) = \frac{91}{6} - \left(\frac{7}{2}\right)^2 = \frac{91}{6} - \frac{49}{4} = \frac{182-147}{12} = \frac{35}{12} \approx 2.9167,
\]
and $\sigma = \sqrt{35/12} \approx 1.708$.
\end{exemplebox}

\begin{exemplebox}[Variance of a discrete uniform distribution]
For $X$ uniform on $\{1,2,\dots,n\}$, $E(X)=\frac{n+1}{2}$ and $E(X^2)=\frac{(n+1)(2n+1)}{6}$. Then
\[
\operatorname{Var}(X) = \frac{(n+1)(2n+1)}{6} - \left(\frac{n+1}{2}\right)^2 = \frac{n^2-1}{12}.
\]
For a fair die ($n=6$), $\operatorname{Var}(X)=\frac{35}{12}$ as above.
\end{exemplebox}

\textbf{Effect of linear transformations on variance.} Adding a constant shifts the whole distribution but does not change its spread; multiplying by a constant stretches or compresses the spread.

\begin{propriete}[Variance of a linear transformation]
For any constants $a,b \in \mathbb{R}$,
\[
\operatorname{Var}(aX + b) = a^2 \operatorname{Var}(X).
\]
Consequently, $\sigma_{aX+b} = |a|\,\sigma_X$.
\end{propriete}

\begin{propriete}[Proof]
Let $\mu = E(X)$. Then $E(aX+b) = a\mu+b$. Now
\begin{align*}
\operatorname{Var}(aX+b) &= E\bigl[(aX+b - (a\mu+b))^2\bigr] \\
&= E\bigl[a^2(X-\mu)^2\bigr] = a^2 E\bigl[(X-\mu)^2\bigr] = a^2 \operatorname{Var}(X).
\end{align*}
\end{propriete}

\begin{exemplebox}[Variance of the game $Y = 2X+100$]
From the die example, $\operatorname{Var}(X)=35/12$. Then
\[
\operatorname{Var}(Y) = 2^2 \times \frac{35}{12} = \frac{140}{12} = \frac{35}{3} \approx 11.667,
\]
and $\sigma_Y = 2\sigma_X \approx 3.416$.
\end{exemplebox}

\begin{exemplebox}[Variance of a shifted variable]
Let $X$ be the die score and $Z = X - 3.5$ (centred variable). Then $E(Z)=0$ and $\operatorname{Var}(Z) = \operatorname{Var}(X) = 35/12$. Shifting by the mean does not change variance.
\end{exemplebox}

\courssub{Efficient computation: the working table}

When computing $E(X)$ and $\operatorname{Var}(X)$ by hand, a structured table avoids errors and saves time.

\begin{methode}[Working table for mean and variance]
\begin{enumerate}
\item List all possible values $x_i$ in the first column.
\item Write the corresponding probabilities $p_i = P(X=x_i)$ in the second column.
\item Compute $x_i p_i$ (third column) and sum to get $E(X) = \sum x_i p_i$.
\item Compute $x_i^2 p_i$ (fourth column) and sum to get $E(X^2) = \sum x_i^2 p_i$.
\item Finally, $\operatorname{Var}(X) = E(X^2) - [E(X)]^2$ and $\sigma = \sqrt{\operatorname{Var}(X)}$.
\end{enumerate}
\end{methode}

\begin{exemplebox}[Number of heads in three fair coin tosses]
Let $X$ be the number of heads. The distribution is binomial with $n=3, p=0.5$:
\[
\begin{array}{c|cccc}
x & 0 & 1 & 2 & 3 \\ \hline
p_i & 1/8 & 3/8 & 3/8 & 1/8
\end{array}
\]
Working table:
\[
\begin{array}{c|c|c|c}
x_i & p_i & x_i p_i & x_i^2 p_i \\ \hline
0 & 1/8 & 0 & 0 \\
1 & 3/8 & 3/8 & 3/8 \\
2 & 3/8 & 6/8 & 12/8 \\
3 & 1/8 & 3/8 & 9/8 \\ \hline
\text{Sum} & 1 & 12/8 = 1.5 & 24/8 = 3
\end{array}
\]
Thus $E(X)=1.5$, $E(X^2)=3$, $\operatorname{Var}(X)=3 - (1.5)^2 = 0.75$, $\sigma = \sqrt{0.75} \approx 0.866$.
\end{exemplebox}

\begin{exemplebox}[A distribution with negative values]
A random variable $X$ has distribution:
\[
\begin{array}{c|ccccc}
x & -2 & -1 & 0 & 1 & 2 \\ \hline
P(X=x) & 0.1 & 0.2 & 0.4 & 0.2 & 0.1
\end{array}
\]
Working table:
\[
\begin{array}{c|c|c|c}
x_i & p_i & x_i p_i & x_i^2 p_i \\ \hline
-2 & 0.1 & -0.2 & 0.4 \\
-1 & 0.2 & -0.2 & 0.2 \\
0 & 0.4 & 0 & 0 \\
1 & 0.2 & 0.2 & 0.2 \\
2 & 0.1 & 0.2 & 0.4 \\ \hline
\text{Sum} & 1 & 0 & 1.2
\end{array}
\]
$E(X)=0$, $E(X^2)=1.2$, $\operatorname{Var}(X)=1.2$, $\sigma = \sqrt{1.2} \approx 1.095$.
\end{exemplebox}

\courssub{Applications: fair games and decision making}

A \textbf{fair game} is one in which the expected net gain is zero. If you pay an entry fee $c$ to play a game that pays $X$, your net gain is $Y = X - c$. The game is fair iff $E(Y)=0$, i.e. $c = E(X)$.

\begin{exoresolu}[Evaluating a game of chance]
A game consists of rolling a fair die. You receive CFA francs equal to the square of the score (i.e. $X^2$). The entry fee is 16 CFA. Is the game fair? If not, who has the advantage?
\tcblower
\textbf{Solution.}
We already computed $E(X^2) = 91/6 \approx 15.1667$ CFA.
Net gain $Y = X^2 - 16$. Then $E(Y) = E(X^2) - 16 = 91/6 - 16 = (91-96)/6 = -5/6 \approx -0.833$ CFA.
Since $E(Y) < 0$, the game is \textbf{not fair}; the player loses on average $0.83$ CFA per game. The organiser has the advantage.
To make it fair, the entry fee should be $E(X^2) = 91/6 \approx 15.17$ CFA.

\end{exoresolu}

\begin{exoresolu}[Finding the fair entry fee]
A bag contains 3 red balls and 2 blue balls. You draw one ball at random. If it is red you win 200 CFA; if blue you lose 100 CFA. What entry fee would make this a fair game?
\tcblower
\textbf{Solution.}
Let $X$ be the net gain before entry fee. $P(\text{red})=3/5$, $P(\text{blue})=2/5$.
$E(X) = 200\times\frac35 + (-100)\times\frac25 = 120 - 40 = 80$ CFA.
For a fair game, entry fee $c = E(X) = 80$ CFA.

\end{exoresolu}

\paragraph{Comparing two investments using mean and standard deviation.}
Suppose you must choose between two investment plans whose returns (in millions of CFA) are discrete random variables $R_1$ and $R_2$ with the following distributions:

\begin{center}
\begin{tabularx}{\linewidth}{|c|X|X|}\hline
\textbf{Return} & \textbf{Plan A ($R_1$)} & \textbf{Plan B ($R_2$)} \\ \hline
2 & 0.2 & 0.1 \\
5 & 0.6 & 0.2 \\
8 & 0.2 & 0.7 \\ \hline
\end{tabularx}
\end{center}

\begin{exoresolu}[Comparing Plan A and Plan B]
Compute $E(R)$ and $\sigma_R$ for each plan and comment.
\tcblower
\textbf{Plan A:}
\begin{align*}
E(R_1) &= 2(0.2) + 5(0.6) + 8(0.2) = 0.4 + 3.0 + 1.6 = 5.0 \\
E(R_1^2) &= 4(0.2) + 25(0.6) + 64(0.2) = 0.8 + 15.0 + 12.8 = 28.6 \\
\operatorname{Var}(R_1) &= 28.6 - 5.0^2 = 3.6,\quad \sigma_1 = \sqrt{3.6} \approx 1.897.
\end{align*}
\textbf{Plan B:}
\begin{align*}
E(R_2) &= 2(0.1) + 5(0.2) + 8(0.7) = 0.2 + 1.0 + 5.6 = 6.8 \\
E(R_2^2) &= 4(0.1) + 25(0.2) + 64(0.7) = 0.4 + 5.0 + 44.8 = 50.2 \\
\operatorname{Var}(R_2) &= 50.2 - 6.8^2 = 50.2 - 46.24 = 3.96,\quad \sigma_2 = \sqrt{3.96} \approx 1.990.
\end{align*}
\textbf{Interpretation:} Plan B has a higher expected return (6.8 vs 5.0 million CFA) but also slightly higher risk (standard deviation 1.99 vs 1.90 million CFA). A risk‑averse investor might prefer Plan A; a risk‑tolerant one would choose Plan B. The decision depends on the investor's utility function, but $E$ and $\sigma$ provide the essential summary.

\end{exoresolu}

\courssub{Illustrations}

\begin{popfigure}
\centering
\begin{tikzpicture}
\begin{axis}[
    ybar,
    bar width=12pt,
    width=0.9\linewidth,
    height=6cm,
    xlabel={Value $x$},
    ylabel={Probability $P(X=x)$},
    xtick={1,2,3,4,5,6,7,8,9},
    xticklabels={1,2,3,4,5,6,7,8,9},
    ymin=0, ymax=0.35,
    legend style={at={(0.5,-0.15)},anchor=north,legend columns=2},
    enlarge x limits=0.15,
]
\addplot[fill=popblue!70, draw=popblue] coordinates {
    (2,0.1) (5,0.2) (8,0.7)
};
\addplot[fill=poporange!70, draw=poporange] coordinates {
    (3,0.1) (5,0.8) (7,0.1)
};
\legend{Distribution B ($\mu=6.8$), Distribution A ($\mu=5.0$)}
\end{axis}
\end{tikzpicture}
\legende{Two distributions with different means and variances. Distribution B has a higher mean but also a slightly larger spread.}
\end{popfigure}

\begin{popfigure}
\centering
\begin{tikzpicture}
\begin{axis}[
    ybar,
    bar width=12pt,
    width=0.9\linewidth,
    height=6cm,
    xlabel={Value $x$},
    ylabel={Probability $P(X=x)$},
    xtick={1,2,3,4,5,6,7,8,9},
    xticklabels={1,2,3,4,5,6,7,8,9},
    ymin=0, ymax=0.5,
    legend style={at={(0.5,-0.15)},anchor=north,legend columns=2},
    enlarge x limits=0.15,
]
\addplot[fill=popgreen!70, draw=popgreen] coordinates {
    (3,0.1) (5,0.8) (7,0.1)
};
\addplot[fill=poppurple!70, draw=poppurple] coordinates {
    (1,0.1) (5,0.8) (9,0.1)
};
\legend{Low variance ($\sigma^2=0.4$), High variance ($\sigma^2=6.4$)}
\end{axis}
\end{tikzpicture}
\legende{Two distributions with the \textbf{same mean} ($\mu=5$) but \textbf{different variances}. The purple distribution is much more spread out.}
\end{popfigure}

\begin{popfigure}
\centering
\begin{tikzpicture}[scale=0.9]
% Number line
\draw[thick, ->] (-0.5,0) -- (10.5,0) node[right] {$x$};
\foreach \x in {1,2,3,4,5,6,7,8,9} {
    \draw (\x,0.1) -- (\x,-0.1) node[below] {\x};
}
% Bars for distribution: values 3,5,7 with probs 0.1, 0.8, 0.1
\draw[fill=popblue!50, draw=popblue] (3,0) rectangle (3.4,0.5);
\draw[fill=popblue!50, draw=popblue] (5,0) rectangle (5.4,4.0);
\draw[fill=popblue!50, draw=popblue] (7,0) rectangle (7.4,0.5);
% Mean marker
\draw[popred, very thick] (5, -0.3) -- (5, 4.3) node[above] {$\mu = E(X) = 5$};
\draw[popred, decorate, decoration={brace, amplitude=5pt}] (4.5, -0.6) -- (5.5, -0.6) node[midway, below=5pt] {Balance point};
% Labels
\node[align=center] at (5, 4.8) {Probability mass acts like weights on a ruler;\\ the mean is the centre of mass.};
\end{tikzpicture}
\legende{The expectation $E(X)$ is the \textbf{balance point} (centre of mass) of the probability bar chart. If the bars were physical weights on a light ruler, the ruler would balance exactly at $x = \mu$.}
\end{popfigure}

\courssub{Common mistakes and warnings}

\begin{attention}[Variance formula traps]
\begin{enumerate}
\item \textbf{Wrong:} $\operatorname{Var}(X) = E(X^2) - E(X)^2$ \emph{without brackets}. The correct formula is $\operatorname{Var}(X) = E(X^2) - \bigl[E(X)\bigr]^2$. The square is \emph{after} taking the expectation of $X$, not before.
\item \textbf{Wrong:} $E(X^2) = \bigl[E(X)\bigr]^2$. This is false in general; it holds only when $X$ is constant (variance zero). For the die, $E(X^2)=91/6 \approx 15.17$ while $[E(X)]^2 = 12.25$.
\item \textbf{Wrong:} $\operatorname{Var}(aX+b) = a\operatorname{Var}(X)+b$. The correct rule is $\operatorname{Var}(aX+b) = a^2\operatorname{Var}(X)$. Adding a constant $b$ does \emph{not} change the variance; multiplying by $a$ scales the variance by $a^2$.
\item \textbf{Units:} Variance is in \emph{squared} units (e.g. CFA$^2$). Standard deviation returns to the original units (CFA). Always report $\sigma$ when interpreting spread in context.
\item \textbf{Wrong:} $E(1/X) = 1/E(X)$. Expectation does not commute with non‑linear functions (see example above).
\end{enumerate}
\end{attention}

\begin{attention}[Examination tips]
\begin{itemize}
\item Always show your working table in GCE exams; it earns method marks even if arithmetic slips occur.
\item When asked for variance, give the exact fraction (e.g. $35/12$) unless a decimal approximation is explicitly requested.
\item For standard deviation, write $\sigma = \sqrt{\operatorname{Var}(X)}$ and simplify the surd if possible (e.g. $\sqrt{0.75} = \sqrt{3}/2$).
\item Check that probabilities sum to 1 before computing expectations.
\item Remember: $E(X+Y)=E(X)+E(Y)$ always, but $\operatorname{Var}(X+Y)=\operatorname{Var}(X)+\operatorname{Var}(Y)$ \textbf{only if $X$ and $Y$ are independent} (covered in Lesson 138).
\end{itemize}
\end{attention}

\begin{exercice}[Practice: mean and variance of a discrete RV]
A random variable $X$ has the following probability distribution:
\[
\begin{array}{c|ccccc}
x & -2 & -1 & 0 & 1 & 2 \\ \hline
P(X=x) & 0.1 & 0.2 & 0.4 & 0.2 & 0.1
\end{array}
\]
\begin{enumerate}
\item Verify that this is a valid probability distribution.
\item Compute $E(X)$ and $E(X^2)$.
\item Find $\operatorname{Var}(X)$ and $\sigma_X$.
\item Let $Y = 3X - 5$. Write down the distribution of $Y$ and compute $E(Y)$ and $\operatorname{Var}(Y)$ using the linear transformation rules (check by direct calculation if you wish).
\end{enumerate}
\end{exercice}

\begin{corrige}[Practice: mean and variance of a discrete RV]
\begin{enumerate}
\item Sum of probabilities = $0.1+0.2+0.4+0.2+0.1 = 1$. Valid.
\item Working table:
\[
\begin{array}{c|c|c|c}
x & p & xp & x^2p \\ \hline
-2 & 0.1 & -0.2 & 0.4 \\
-1 & 0.2 & -0.2 & 0.2 \\
0 & 0.4 & 0 & 0 \\
1 & 0.2 & 0.2 & 0.2 \\
2 & 0.1 & 0.2 & 0.4 \\ \hline
\text{Sum} & 1 & 0 & 1.2
\end{array}
\]
$E(X)=0$, $E(X^2)=1.2$.
\item $\operatorname{Var}(X) = 1.2 - 0^2 = 1.2$, $\sigma_X = \sqrt{1.2} = \sqrt{6/5} = \sqrt{30}/5 \approx 1.095$.
\item $Y = 3X-5$. By linearity: $E(Y) = 3E(X)-5 = -5$. $\operatorname{Var}(Y) = 3^2\operatorname{Var}(X) = 9\times1.2 = 10.8$.
Direct check: $Y$ takes values $-11, -8, -5, -2, 1$ with same probabilities.
$E(Y) = -11(0.1)-8(0.2)-5(0.4)-2(0.2)+1(0.1) = -1.1-1.6-2.0-0.4+0.1 = -5.0$. OK.
$E(Y^2) = 121(0.1)+64(0.2)+25(0.4)+4(0.2)+1(0.1) = 12.1+12.8+10.0+0.8+0.1 = 35.8$.
$\operatorname{Var}(Y) = 35.8 - (-5)^2 = 35.8 - 25 = 10.8$. OK.
\end{enumerate}
\end{corrige}

\begin{exercice}[GCE-style question]
The discrete random variable $X$ has probability distribution given by
\[
P(X=x) = \frac{x}{10}, \quad x = 1,2,3,4.
\]
\begin{enumerate}
\item Show that this is a valid probability distribution.
\item Find $E(X)$ and $\operatorname{Var}(X)$.
\item Given that $Y = 2X - 3$, find $E(Y)$ and $\operatorname{Var}(Y)$.
\end{enumerate}
\end{exercice}

\begin{corrige}[GCE-style question]
\begin{enumerate}
\item $\sum_{x=1}^4 P(X=x) = \frac{1+2+3+4}{10} = \frac{10}{10} = 1$. Valid.
\item Working table:
\[
\begin{array}{c|c|c|c}
x & p=x/10 & xp & x^2p \\ \hline
1 & 0.1 & 0.1 & 0.1 \\
2 & 0.2 & 0.4 & 0.8 \\
3 & 0.3 & 0.9 & 2.7 \\
4 & 0.4 & 1.6 & 6.4 \\ \hline
\text{Sum} & 1 & 3.0 & 10.0
\end{array}
\]
$E(X)=3$, $E(X^2)=10$, $\operatorname{Var}(X)=10-9=1$.
\item $E(Y) = 2E(X)-3 = 2(3)-3 = 3$. $\operatorname{Var}(Y) = 2^2\operatorname{Var}(X) = 4\times1 = 4$.
\end{enumerate}
\end{corrige}

\begin{aretenir}[Key points of this section]
\begin{itemize}
\item $E(X) = \sum x_i P(X=x_i)$ is the theoretical mean (long‑run average).
\item $E(g(X)) = \sum g(x_i) P(X=x_i)$; in particular $E(X^2) = \sum x_i^2 p_i$.
\item \textbf{Linearity:} $E(aX+b) = aE(X)+b$ and $E(X+Y)=E(X)+E(Y)$ (always, no independence needed).
\item $\operatorname{Var}(X) = E(X^2) - [E(X)]^2 \ge 0$; $\sigma = \sqrt{\operatorname{Var}(X)}$.
\item $\operatorname{Var}(aX+b) = a^2 \operatorname{Var}(X)$; adding a constant does not affect variance.
\item Use a working table ($x$, $p$, $xp$, $x^2p$) for accurate hand calculations.
\item A fair game has $E(\text{net gain}) = 0$; fair entry fee $= E(\text{payout})$.
\item When comparing options, consider both $E(X)$ (return) and $\sigma$ (risk).
\item $E(X^2) \neq [E(X)]^2$ unless $\operatorname{Var}(X)=0$.
\item For independent $X,Y$ (later lesson): $\operatorname{Var}(X+Y)=\operatorname{Var}(X)+\operatorname{Var}(Y)$.
\end{itemize}
\end{aretenir}

\coursec{The cumulative distribution function, mode and median}

In the previous section we summarised a discrete distribution by a single number, its \cle{expectation} $E(X)$, and measured its spread with the \cle{variance} and \cle{standard deviation}. These are excellent summaries, but they do not tell the whole story. Two distributions can share the same mean yet look completely different. To describe a distribution \emph{fully}, we need a function that accumulates probabilities as we move along the number line: the \cle{cumulative distribution function}. From it we can also read off two other important measures of location, the \cle{mode} and the \cle{median}, and compare them with the mean.

\courssub{The cumulative distribution function}

\textbf{Motivation.}\par
Suppose a trader at Mokolo market records the number $X$ of bags of rice she sells per day, with distribution

\begin{center}
\begin{tabular}{|c|c|c|c|c|}
\hline
\rowcolor{popblueL} $x_i$ & $0$ & $1$ & $2$ & $3$ \\
\hline $P(X=x_i)$ & $0.1$ & $0.3$ & $0.4$ & $0.2$ \\
\hline
\end{tabular}
\end{center}

Questions such as ``what is the probability that she sells \emph{at most} $2$ bags?'' arise constantly. The answer is $P(X\le 2)=0.1+0.3+0.4=0.8$. Rather than recomputing such sums every time, we build a single function that stores all these running totals.

\begin{definition}[Cumulative distribution function]
Let $X$ be a discrete random variable. The \cle{cumulative distribution function} (c.d.f.) of $X$ is the function $F$ defined for \emph{every real number} $x$ by
\[ F(x) = P(X \le x). \]
\end{definition}

\textbf{Constructing $F$ from the probability table.}\par
For a distribution taking values $x_1 < x_2 < \dots < x_n$, the c.d.f. is built by \cle{cumulative addition} of the probabilities:
\begin{itemize}
\item for $x < x_1$: $F(x) = 0$;
\item for $x_1 \le x < x_2$: $F(x) = P(X=x_1)$;
\item for $x_2 \le x < x_3$: $F(x) = P(X=x_1)+P(X=x_2)$;
\item \dots
\item for $x \ge x_n$: $F(x) = 1$.
\end{itemize}
For our trader, this gives
\[ F(x) = \begin{cases} 0 & \text{if } x < 0,\\ 0.1 & \text{if } 0 \le x < 1,\\ 0.4 & \text{if } 1 \le x < 2,\\ 0.8 & \text{if } 2 \le x < 3,\\ 1 & \text{if } x \ge 3. \end{cases} \]

\begin{attention}[$F$ is defined for ALL real $x$]
A frequent GCE trap: $F(x)$ is defined for \textbf{every} real number $x$, not only at the values $x_i$. Between two consecutive values, $F$ stays \emph{constant}. For instance, $F(1.7) = P(X \le 1.7) = P(X=0)+P(X=1) = 0.4$, and $F(-5) = 0$. Never leave gaps in the definition of $F$.
\end{attention}

\begin{propriete}[Properties of the c.d.f.]
For any random variable $X$ with c.d.f. $F$:
\begin{enumerate}
\item $0 \le F(x) \le 1$ for all real $x$;
\item $F$ is \cle{non-decreasing}: if $a \le b$ then $F(a) \le F(b)$;
\item $F(x) \to 0$ as $x \to -\infty$ and $F(x) \to 1$ as $x \to +\infty$;
\item $F$ is \cle{right-continuous}: at each value $x_i$ the function jumps, and the value at the jump is the \emph{upper} value, $F(x_i) = P(X \le x_i)$.
\end{enumerate}
(The proof of these properties is not examinable, but you must know and use them.)
\end{propriete}

\textbf{Justification of (2).} If $a \le b$, the event $\{X \le a\}$ is contained in the event $\{X \le b\}$, so $P(X \le a) \le P(X \le b)$, that is $F(a) \le F(b)$. \hfill $\square$

\begin{popfigure}
\begin{center}
\begin{tikzpicture}
\begin{axis}[
    width=11cm, height=6.5cm,
    axis lines=middle,
    xlabel={$x$}, ylabel={$F(x)$},
    xmin=-1.2, xmax=4.2, ymin=-0.08, ymax=1.25,
    xtick={0,1,2,3}, ytick={0.1,0.4,0.8,1},
    yticklabels={$0.1$,$0.4$,$0.8$,$1$},
    every axis x label/.style={at={(ticklabel* cs:1.02)}, anchor=west},
    every axis y label/.style={at={(ticklabel* cs:1.02)}, anchor=south},
]
\addplot[very thick, popblue] coordinates {(-1.2,0) (0,0)};
\addplot[very thick, popblue] coordinates {(0,0.1) (1,0.1)};
\addplot[very thick, popblue] coordinates {(1,0.4) (2,0.4)};
\addplot[very thick, popblue] coordinates {(2,0.8) (3,0.8)};
\addplot[very thick, popblue] coordinates {(3,1) (4.2,1)};
\addplot[only marks, mark=*, mark size=1.8pt, popblue] coordinates {(0,0.1) (1,0.4) (2,0.8) (3,1)};
\addplot[only marks, mark=o, mark size=1.8pt, popblue] coordinates {(0,0) (1,0.1) (2,0.4) (3,0.8)};
\end{axis}
\end{tikzpicture}
\end{center}
\legende{The c.d.f. $F(x)$ of the rice-sales distribution: a step function with jumps at $x=0,1,2,3$. Closed dots show the value taken at each jump (right-continuity).}
\end{popfigure}

\textbf{Recovering probabilities from $F$.}\par
The c.d.f. contains \emph{all} the information of the distribution, so we can travel back from $F$ to the individual probabilities.

\begin{propriete}[From the c.d.f. back to probabilities]
If $X$ takes values $x_1 < x_2 < \dots < x_n$, then
\[ P(X = x_i) = F(x_i) - F(x_{i-1}) \qquad (i \ge 2), \quad P(X=x_1) = F(x_1), \]
and for any real numbers $a < b$,
\[ P(a < X \le b) = F(b) - F(a). \]
\end{propriete}

\textbf{Justification.} The event $\{X \le b\}$ is the disjoint union of $\{X \le a\}$ and $\{a < X \le b\}$, so $F(b) = F(a) + P(a < X \le b)$, which gives the second formula. Taking $a = x_{i-1}$ and $b = x_i$ gives the first, since the only value of $X$ in $(x_{i-1}, x_i]$ is $x_i$. \hfill $\square$

\begin{attention}[Strict versus weak inequalities]
Because $X$ is discrete, $P(a < X \le b) = F(b) - F(a)$ but $P(a \le X \le b) = F(b) - F(a) + P(X=a)$ when $a$ is one of the values. Always check whether the endpoints are included.
\end{attention}

\begin{exoresolu}[Reading a c.d.f.]
A discrete random variable $X$ has c.d.f.
\[ F(x) = \begin{cases} 0 & x < 1,\\ 0.2 & 1 \le x < 2,\\ 0.5 & 2 \le x < 4,\\ 0.9 & 4 \le x < 5,\\ 1 & x \ge 5. \end{cases} \]
Find (a) the probability distribution of $X$; (b) $P(2 < X \le 5)$; (c) $P(X \ge 4)$.
\tcblower
\textbf{Solution.}
(a) The jumps of $F$ occur at $x = 1, 2, 4, 5$, so these are the values of $X$, and each probability is the size of the jump:
\begin{align*}
P(X=1) &= 0.2 - 0 = 0.2, & P(X=2) &= 0.5 - 0.2 = 0.3,\\
P(X=4) &= 0.9 - 0.5 = 0.4, & P(X=5) &= 1 - 0.9 = 0.1.
\end{align*}
Check: $0.2+0.3+0.4+0.1 = 1$. \checkmark

(b) $P(2 < X \le 5) = F(5) - F(2) = 1 - 0.5 = 0.5$ (indeed $P(X=4)+P(X=5) = 0.5$).

(c) $P(X \ge 4) = 1 - P(X \le 2) = 1 - F(2) = 1 - 0.5 = 0.5$. Note we subtract $F(2)$, not $F(4)$, because the event $\{X \ge 4\}$ is the complement of $\{X \le 2\}$ (there is no value $3$).
\end{exoresolu}

\courssub{The mode}

\textbf{Intuition.}\par
The \cle{mode} answers the question: \emph{which outcome is the most likely?} It is simply the value of $X$ that carries the greatest probability — the tallest bar of the bar chart.

\begin{definition}[Mode of a discrete distribution]
A \cle{mode} of a discrete random variable $X$ is a value $x_i$ such that
\[ P(X = x_i) \ge P(X = x_j) \quad \text{for every value } x_j. \]
A distribution with one mode is called \cle{unimodal}; if two values share the greatest probability, the distribution is \cle{bimodal}.
\end{definition}

For the rice-sales distribution, $P(X=2) = 0.4$ is the largest probability, so the mode is $2$ bags. If instead the probabilities had been $0.1, 0.4, 0.4, 0.1$, the distribution would be bimodal with modes $1$ and $2$.

\courssub{The median}

\textbf{Intuition.}\par
The \cle{median} is the ``half-way point'' of the distribution: a value $m$ such that at least half of the probability lies at or below $m$. For a discrete variable we must phrase this carefully using the c.d.f.

\begin{definition}[Median of a discrete distribution]
The \cle{median} $m$ of a discrete random variable $X$ is the \emph{smallest} value of $X$ such that
\[ F(m) = P(X \le m) \ge 0.5. \]
\end{definition}

\begin{methode}[Finding the median from a probability table]
\begin{enumerate}
\item Write the values $x_1 < x_2 < \dots$ in increasing order.
\item Compute the cumulative probabilities $F(x_1), F(x_2), \dots$ by adding the probabilities from the left.
\item Stop at the \textbf{first} value $x_i$ where the cumulative total reaches or exceeds $0.5$: that value is the median.
\end{enumerate}
\end{methode}

\begin{attention}[Do not interpolate!]
The median of a discrete distribution is always one of the \emph{actual values} of $X$. It is \textbf{not} necessary that $F(m)$ equals exactly $0.5$. If $F(1) = 0.4$ and $F(2) = 0.8$, the median is $2$ — do \emph{not} invent an intermediate value such as $1.25$ by interpolation. Interpolation is a mistake the GCE examiners penalise every year.
\end{attention}

For the rice-sales distribution: $F(0) = 0.1$, $F(1) = 0.4$, $F(2) = 0.8$. The first cumulative total reaching $0.5$ is $F(2) = 0.8$, so the median is $m = 2$.

\courssub{Comparing mean, mode and median}

These three numbers all describe the ``centre'' of a distribution, but from different points of view:
\begin{itemize}
\item the \cle{mode} is the most \emph{probable} value;
\item the \cle{median} splits the distribution so that at least half the probability lies on each side;
\item the \cle{mean} $E(X)$ is the \emph{balance point}, taking every value into account weighted by its probability.
\end{itemize}

\begin{exemplebox}[Mean, mode and median of one distribution]
For the rice-sales distribution:
\[ E(X) = 0(0.1) + 1(0.3) + 2(0.4) + 3(0.2) = 1.7. \]
So we have: mode $= 2$, median $= 2$, mean $= 1.7$. The mean is pulled slightly below the mode and median because the value $0$ drags the balance point to the left, while the median and mode, which ignore the sizes of the values and look only at probabilities, stay at $2$.
\end{exemplebox}

\begin{popfigure}
\begin{center}
\begin{tikzpicture}
\begin{axis}[
    width=11cm, height=6.5cm,
    ybar, bar width=22pt,
    axis lines=left,
    xlabel={$x$}, ylabel={$P(X=x)$},
    xmin=-0.8, xmax=3.8, ymin=0, ymax=0.55,
    xtick={0,1,2,3},
    ytick={0.1,0.2,0.3,0.4,0.5},
]
\addplot[fill=popblue, draw=popdark] coordinates {(0,0.1) (1,0.3) (2,0.4) (3,0.2)};
\addplot[dashed, very thick, poporange] coordinates {(1.7,0) (1.7,0.52)};
\addplot[dashed, very thick, popgreen] coordinates {(2,0) (2,0.52)};
\node[poporange, anchor=west, font=\small] at (axis cs:1.78,0.5) {mean $1.7$};
\node[popgreen, anchor=east, font=\small] at (axis cs:1.95,0.47) {mode = median $= 2$};
\end{axis}
\end{tikzpicture}
\end{center}
\legende{Bar chart of the distribution with the mean (orange dashed line) and the common value of the mode and median (green dashed line).}
\end{popfigure}

\begin{savaistu}[Why the median matters in Cameroon]
In many real-life situations in Cameroon — household sizes in Yaoundé, daily earnings of moto-taxi drivers in Douala, or cocoa yields per hectare in the South-West — the distribution is often skewed. The mean can be distorted by a few very large values, while the median gives a more representative ``typical'' value. That is why government reports on living standards often quote the median income rather than the mean.
\end{savaistu}

\begin{exoresolu}[A complete study]
A random variable $X$ has distribution
\begin{center}
\begin{tabular}{|c|c|c|c|c|c|}
\hline
\rowcolor{popblueL} $x_i$ & $0$ & $1$ & $2$ & $3$ & $4$ \\
\hline $P(X=x_i)$ & $0.05$ & $0.25$ & $0.4$ & $0.2$ & $0.1$ \\
\hline
\end{tabular}
\end{center}
(a) Construct the c.d.f. $F$. (b) Find the mode and the median. (c) Compute $E(X)$ and comment. (d) Find $P(1 < X \le 3)$.
\tcblower
\textbf{Solution.}
(a) Cumulative sums: $F(0) = 0.05$, $F(1) = 0.30$, $F(2) = 0.70$, $F(3) = 0.90$, $F(4) = 1$. Hence
\[ F(x) = \begin{cases} 0 & x < 0,\\ 0.05 & 0 \le x < 1,\\ 0.30 & 1 \le x < 2,\\ 0.70 & 2 \le x < 3,\\ 0.90 & 3 \le x < 4,\\ 1 & x \ge 4. \end{cases} \]

(b) The greatest probability is $P(X=2) = 0.4$, so the \textbf{mode is $2$}. The first cumulative total reaching $0.5$ is $F(2) = 0.70$, so the \textbf{median is $2$}.

(c) $E(X) = 0(0.05) + 1(0.25) + 2(0.4) + 3(0.2) + 4(0.1) = 0.25 + 0.8 + 0.6 + 0.4 = 2.05$.
The mean ($2.05$), median ($2$) and mode ($2$) are very close, which reflects the fact that the distribution is nearly symmetric about $x = 2$.

(d) $P(1 < X \le 3) = F(3) - F(1) = 0.90 - 0.30 = 0.60$ (check: $P(X=2) + P(X=3) = 0.4 + 0.2 = 0.6$). \checkmark
\end{exoresolu}

\begin{aretenir}[Key points of this section]
\begin{itemize}
\item The c.d.f. is $F(x) = P(X \le x)$, defined for \textbf{every} real $x$; it is a non-decreasing step function, right-continuous, with $F(x) \to 0$ as $x \to -\infty$ and $F(x) \to 1$ as $x \to +\infty$.
\item Probabilities are recovered from $F$: $P(X = x_i) = F(x_i) - F(x_{i-1})$ and $P(a < X \le b) = F(b) - F(a)$.
\item The \cle{mode} is the value with the greatest probability (a distribution may be bimodal).
\item The \cle{median} is the smallest value $m$ with $F(m) \ge 0.5$; it is always an actual value of $X$ — never interpolate.
\item Mean, median and mode coincide for a symmetric distribution but differ for skewed ones.
\end{itemize}
\end{aretenir}

\begin{exercice}[Practice]
A discrete random variable $X$ takes values $1, 2, 3, 4$ with probabilities $0.3, 0.3, 0.2, 0.2$ respectively.
\begin{enumerate}
\item Write down the c.d.f. $F(x)$ as a piecewise function and sketch it.
\item State the mode(s) of $X$.
\item Find the median of $X$.
\item Compute $E(X)$ and compare the three measures of location.
\item Evaluate $P(1 < X \le 3)$ and $P(X \ge 2)$ using $F$.
\end{enumerate}
\end{exercice}

\begin{corrige}[Exercise 1]
\textbf{1.} Cumulative sums: $F(1) = 0.3$, $F(2) = 0.6$, $F(3) = 0.8$, $F(4) = 1$. So
\[ F(x) = \begin{cases} 0 & x < 1,\\ 0.3 & 1 \le x < 2,\\ 0.6 & 2 \le x < 3,\\ 0.8 & 3 \le x < 4,\\ 1 & x \ge 4. \end{cases} \]
The graph is a step function with jumps of heights $0.3, 0.3, 0.2, 0.2$ at $x = 1, 2, 3, 4$.

\textbf{2.} The greatest probability is $0.3$, attained at both $x = 1$ and $x = 2$: the distribution is \textbf{bimodal} with modes $1$ and $2$.

\textbf{3.} $F(1) = 0.3 < 0.5$ and $F(2) = 0.6 \ge 0.5$, so the median is $m = 2$.

\textbf{4.} $E(X) = 1(0.3) + 2(0.3) + 3(0.2) + 4(0.2) = 0.3 + 0.6 + 0.6 + 0.8 = 2.3$. The mean ($2.3$) exceeds the median ($2$): the larger probabilities on the small values are balanced by the heavier values $3$ and $4$, which pull the balance point to the right.

\textbf{5.} $P(1 < X \le 3) = F(3) - F(1) = 0.8 - 0.3 = 0.5$; $P(X \ge 2) = 1 - F(1) = 1 - 0.3 = 0.7$.
\end{corrige}

\coursec{The distribution of X₁ + X₂}

In the previous section we studied how to summarise a single discrete random variable using its cumulative distribution function, its mode and its median.  Many real‑world situations, however, involve the \emph{sum} of two or more random variables: the total score when two dice are thrown, the total number of customers arriving at a shop over two consecutive days, the combined weight of two randomly chosen bags of rice, etc.  This section develops the tools to obtain the probability distribution of a sum $S = X_1 + X_2$, and to compute its expectation and variance.  The results here are the foundation for the binomial distribution (sum of $n$ independent Bernoulli trials) and for the Central Limit Theorem later in the course.

\courssub{Why sums of random variables?}

\begin{experience}[Total score of two dice]
Suppose we roll two fair six‑sided dice, one red and one blue.  Let $X_1$ be the score on the red die and $X_2$ the score on the blue die.  Both are discrete uniform on $\{1,2,3,4,5,6\}$.  We are often interested in the \textbf{total score} $S = X_1 + X_2$.  The possible values of $S$ range from $2$ to $12$, but they are \emph{not} equally likely.  For instance, $S=7$ can occur in six ways $(1,6),(2,5),\dots,(6,1)$ while $S=2$ occurs only as $(1,1)$.  To find the distribution of $S$ we must combine the distributions of $X_1$ and $X_2$ in a systematic way.
\end{experience}

\begin{experience}[Total customers over two days]
A small café in Yaoundé records the number of customers each day.  Let $X_1$ be the number on Monday and $X_2$ the number on Tuesday.  If the daily numbers are independent and each follows a known distribution (say Poisson with mean $5$), the owner may want the distribution of the two‑day total $S = X_1 + X_2$ to plan staffing and supplies.
\end{experience}

\courssub{Independence of two discrete random variables}

Before we can combine distributions we need a precise notion of independence for random variables.

\begin{definition}[Independence of two discrete random variables]
Two discrete random variables $X_1$ and $X_2$ are \cle{independent} if and only if for every possible value $x$ of $X_1$ and every possible value $y$ of $X_2$,
\[
P(X_1 = x \text{ and } X_2 = y) = P(X_1 = x) \cdot P(X_2 = y).
\]
Equivalently, the joint probability mass function factorises: $p_{X_1,X_2}(x,y) = p_{X_1}(x)\,p_{X_2}(y)$ for all $x,y$.
\end{definition}

\begin{attention}[Independence vs. uncorrelated]
Independence is a \emph{stronger} condition than zero correlation.  If $X_1$ and $X_2$ are independent then they are uncorrelated, but the converse is false in general.  In the GCE syllabus we only use the definition above; always check whether the problem states or implies independence before using the product rule for joint probabilities.
\end{attention}

\courssub{Constructing the distribution of $S = X_1 + X_2$ by enumeration}

When $X_1$ and $X_2$ are independent, the most transparent way to obtain the distribution of $S = X_1 + X_2$ is to build a \textbf{two‑way table} (also called a joint probability table).

\begin{methode}[Two‑way table for the sum of two independent discrete random variables]
\begin{enumerate}
\item List the possible values of $X_1$ down the rows and those of $X_2$ across the columns.
\item In each cell write the sum $x_1 + x_2$ (the value of $S$) and the joint probability $P(X_1=x_1)P(X_2=x_2)$.
\item Collect all cells that give the same sum $s$; add their probabilities to obtain $P(S=s)$.
\item Present the result as a probability distribution table for $S$.
\end{enumerate}
\end{methode}

Let us apply this to the two‑dice example.

\begin{exoresolu}[Sum of two fair dice]
Let $X_1$ and $X_2$ be the scores on two independent fair dice.  Find the probability distribution of $S = X_1 + X_2$.

\tcblower
\textbf{Step 1: Two‑way table.}  Each die has outcomes $1,\dots,6$ with probability $\frac16$.  The joint probability of any pair $(x_1,x_2)$ is $\frac16 \times \frac16 = \frac1{36}$.

\begin{center}
\begin{tabular}{|c|cccccc|}\hline
$X_1 \backslash X_2$ & 1 & 2 & 3 & 4 & 5 & 6 \\ \hline
1 & 2 & 3 & 4 & 5 & 6 & 7 \\
2 & 3 & 4 & 5 & 6 & 7 & 8 \\
3 & 4 & 5 & 6 & 7 & 8 & 9 \\
4 & 5 & 6 & 7 & 8 & 9 & 10 \\
5 & 6 & 7 & 8 & 9 & 10 & 11 \\
6 & 7 & 8 & 9 & 10 & 11 & 12 \\ \hline
\end{tabular}
\end{center}
Each cell has probability $\frac1{36}$.  The diagonals (constant $x_1+x_2$) give the values of $S$.

\textbf{Step 2: Collect probabilities along diagonals.}

\begin{center}
\begin{tabular}{|c|ccccccccccc|}\hline
$s$ & 2 & 3 & 4 & 5 & 6 & 7 & 8 & 9 & 10 & 11 & 12 \\ \hline
$P(S=s)$ & $\frac1{36}$ & $\frac2{36}$ & $\frac3{36}$ & $\frac4{36}$ & $\frac5{36}$ & $\frac6{36}$ & $\frac5{36}$ & $\frac4{36}$ & $\frac3{36}$ & $\frac2{36}$ & $\frac1{36}$ \\ \hline
\end{tabular}
\end{center}

\textbf{Step 3: Simplify (optional).}  $P(S=s) = \frac{6-|s-7|}{36}$ for $s=2,\dots,12$.

This is the classic \emph{triangular distribution}.
\end{exoresolu}

\begin{popfigure}
\begin{tikzpicture}[scale=0.7]
  \foreach \x in {1,...,6} {
    \foreach \y in {1,...,6} {
      \pgfmathtruncatemacro{\sval}{\x+\y}
      \draw (\x-1,6-\y) rectangle (\x,7-\y);
      \node at (\x-0.5,6.5-\y) {\small $(\x,\y)$};
      \node[font=\tiny, text=popdark] at (\x-0.5,6.2-\y) {$s=\sval$};
      % Highlight diagonals with different colors
      \ifnum\sval=7
        \fill[poporange!30, opacity=0.5] (\x-1,6-\y) rectangle (\x,7-\y);
      \else\ifnum\sval=6
        \fill[poporange!20, opacity=0.5] (\x-1,6-\y) rectangle (\x,7-\y);
      \else\ifnum\sval=8
        \fill[poporange!20, opacity=0.5] (\x-1,6-\y) rectangle (\x,7-\y);
      \else\ifnum\sval=5
        \fill[popblue!15, opacity=0.5] (\x-1,6-\y) rectangle (\x,7-\y);
      \else\ifnum\sval=9
        \fill[popblue!15, opacity=0.5] (\x-1,6-\y) rectangle (\x,7-\y);
      \else\ifnum\sval=4
        \fill[popgreen!15, opacity=0.5] (\x-1,6-\y) rectangle (\x,7-\y);
      \else\ifnum\sval=10
        \fill[popgreen!15, opacity=0.5] (\x-1,6-\y) rectangle (\x,7-\y);
      \else\ifnum\sval=3
        \fill[poppink!15, opacity=0.5] (\x-1,6-\y) rectangle (\x,7-\y);
      \else\ifnum\sval=11
        \fill[poppink!15, opacity=0.5] (\x-1,6-\y) rectangle (\x,7-\y);
      \else
        \fill[popgold!15, opacity=0.5] (\x-1,6-\y) rectangle (\x,7-\y);
      \fi\fi\fi\fi\fi\fi\fi\fi\fi
    }
  }
  % Add labels for each sum
  \foreach \s in {2,...,12} {
    \pgfmathtruncatemacro{\cnt}{6-abs(\s-7)}
    \ifnum\s<7
      \pgfmathsetmacro{\ylab}{6.5 - (\s-2)*0.5}
      \node[popdark, font=\small\bfseries, anchor=west] at (6.3, \ylab) {$s=\s$ ($\cnt$ ways)};
    \else
      \pgfmathsetmacro{\ylab}{6.5 - (\s-2)*0.5}
      \node[popdark, font=\small\bfseries, anchor=west] at (6.3, \ylab) {$s=\s$ ($\cnt$ ways)};
    \fi
  }
  \node[above] at (3,7) {$X_2$};
  \node[left] at (-0.5,3.5) {$X_1$};
  \draw[->, thick] (0,0) -- (0,7) node[left] {$X_1$};
  \draw[->, thick] (0,0) -- (7,0) node[below] {$X_2$};
\end{tikzpicture}
\legende{The 36 equally likely outcomes of two fair dice.  Cells with the same sum $s$ lie on diagonals; the number of cells on each diagonal gives the frequency of that sum.}
\end{popfigure}

\begin{popfigure}
\begin{tikzpicture}
\begin{axis}[
  ybar,
  bar width=6pt,
  width=0.9\linewidth,
  height=6cm,
  xlabel={Sum $s$},
  ylabel={$P(S=s)$},
  xtick={2,3,4,5,6,7,8,9,10,11,12},
  ytick={0,0.05,0.1,0.15,0.1667},
  ymin=0, ymax=0.18,
  tick label style={font=\small},
  label style={font=\small},
  axis lines=left,
  enlarge x limits=0.05,
  enlarge y limits=0.1,
]
\addplot[fill=popblue, draw=popdark] coordinates {
  (2,0.02778) (3,0.05556) (4,0.08333) (5,0.11111) (6,0.13889)
  (7,0.16667) (8,0.13889) (9,0.11111) (10,0.08333) (11,0.05556) (12,0.02778)
};
\end{axis}
\end{tikzpicture}
\legende{Probability mass function of the sum of two fair dice — a symmetric triangular distribution peaking at $s=7$.}
\end{popfigure}

\courssub{The convolution formula}

The two‑way table method is perfectly general.  For independent $X_1$ and $X_2$ with probability mass functions $p_1(x)=P(X_1=x)$ and $p_2(y)=P(X_2=y)$, the probability that $S=X_1+X_2$ takes the value $s$ is obtained by summing the joint probabilities over all pairs $(x,y)$ with $x+y=s$.  This is the \cle{convolution} of the two distributions.

\begin{propriete}[Convolution formula for the sum of two independent discrete random variables]
If $X_1$ and $X_2$ are independent discrete random variables with probability mass functions $p_1$ and $p_2$, then for any integer $s$,
\[
P(S = s) = P(X_1 + X_2 = s) = \sum_{x} p_1(x)\,p_2(s-x),
\]
where the sum runs over all $x$ such that $p_1(x)>0$ and $p_2(s-x)>0$.
\end{propriete}

\begin{exemplebox}[Convolution of two independent uniform variables on $\{1,2,3\}$]
Let $X_1$ and $X_2$ be independent, each uniformly distributed on $\{1,2,3\}$ (so $p_1(x)=p_2(x)=\frac13$ for $x=1,2,3$).  Find the distribution of $S=X_1+X_2$.

The possible sums are $s=2,3,4,5,6$.  Using the convolution formula:
\begin{align*}
P(S=2) &= p_1(1)p_2(1) = \tfrac13\cdot\tfrac13 = \tfrac19,\\
P(S=3) &= p_1(1)p_2(2)+p_1(2)p_2(1) = \tfrac19+\tfrac19 = \tfrac29,\\
P(S=4) &= p_1(1)p_2(3)+p_1(2)p_2(2)+p_1(3)p_2(1) = \tfrac19+\tfrac19+\tfrac19 = \tfrac39,\\
P(S=5) &= p_1(2)p_2(3)+p_1(3)p_2(2) = \tfrac19+\tfrac19 = \tfrac29,\\
P(S=6) &= p_1(3)p_2(3) = \tfrac19.
\end{align*}
The distribution is again triangular: $\frac19,\frac29,\frac39,\frac29,\frac19$.
\end{exemplebox}

\courssub{Expectation and variance of a sum}

One of the most powerful properties of expectation is its \textbf{linearity} — it holds \emph{without any independence assumption}.

\begin{propriete}[Expectation of a sum]
For any two discrete random variables $X_1$ and $X_2$ (independent or not),
\[
E(X_1 + X_2) = E(X_1) + E(X_2).
\]
More generally, $E\left(\sum_{i=1}^n X_i\right) = \sum_{i=1}^n E(X_i)$.
\end{propriete}

\begin{proof}
\begin{align*}
E(X_1+X_2) &= \sum_x \sum_y (x+y)\,P(X_1=x,X_2=y) \\
&= \sum_x \sum_y x\,P(X_1=x,X_2=y) + \sum_x \sum_y y\,P(X_1=x,X_2=y) \\
&= \sum_x x \Bigl(\sum_y P(X_1=x,X_2=y)\Bigr) + \sum_y y \Bigl(\sum_x P(X_1=x,X_2=y)\Bigr) \\
&= \sum_x x\,P(X_1=x) + \sum_y y\,P(X_2=y) = E(X_1)+E(X_2).
\end{align*}
No independence was used; only the law of total probability for the marginals.
\end{proof}

For variance, independence \emph{is} essential.

\begin{propriete}[Variance of a sum of independent random variables]
If $X_1$ and $X_2$ are \cle{independent}, then
\[
\operatorname{Var}(X_1 + X_2) = \operatorname{Var}(X_1) + \operatorname{Var}(X_2).
\]
If $X_1$ and $X_2$ are \emph{not} independent, this formula does \emph{not} hold in general; one must use $\operatorname{Var}(X_1+X_2) = \operatorname{Var}(X_1)+\operatorname{Var}(X_2)+2\operatorname{Cov}(X_1,X_2)$ (covariance is beyond the GCE syllabus).
\end{propriete}

\begin{proof}[Sketch (examinable)]
$\operatorname{Var}(X_1+X_2) = E[(X_1+X_2)^2] - [E(X_1+X_2)]^2$.
Expand $(X_1+X_2)^2 = X_1^2 + 2X_1X_2 + X_2^2$.  By linearity of expectation,
$E[(X_1+X_2)^2] = E(X_1^2) + 2E(X_1X_2) + E(X_2^2)$.
Independence gives $E(X_1X_2) = E(X_1)E(X_2)$.  Then
\begin{align*}
\operatorname{Var}(X_1+X_2) &= E(X_1^2)-E(X_1)^2 + E(X_2^2)-E(X_2)^2 \\
&= \operatorname{Var}(X_1)+\operatorname{Var}(X_2).
\end{align*}
\end{proof}

\begin{attention}[Variances add, standard deviations do NOT]
A very common mistake is to write $\sigma_{X_1+X_2} = \sigma_{X_1} + \sigma_{X_2}$.  This is \textbf{false}.  The correct rule for independent variables is
\[
\sigma_{X_1+X_2}^2 = \sigma_{X_1}^2 + \sigma_{X_2}^2 \quad\Longrightarrow\quad \sigma_{X_1+X_2} = \sqrt{\sigma_{X_1}^2 + \sigma_{X_2}^2}.
\]
Always work with variances when combining independent random variables; take the square root only at the very end if a standard deviation is required.
\end{attention}

\courssub{Worked examples: checking the rules}

\begin{exoresolu}[Sum of two fair dice — direct vs. rules]
Let $X_1,X_2$ be independent fair dice scores.  We already have the distribution of $S=X_1+X_2$.  Verify that $E(S)=E(X_1)+E(X_2)$ and $\operatorname{Var}(S)=\operatorname{Var}(X_1)+\operatorname{Var}(X_2)$.

\tcblower
\textbf{For a single die:} $E(X_i) = \frac{1+2+3+4+5+6}{6} = 3.5$.
$E(X_i^2) = \frac{1^2+2^2+3^2+4^2+5^2+6^2}{6} = \frac{91}{6} \approx 15.1667$.
$\operatorname{Var}(X_i) = \frac{91}{6} - (3.5)^2 = \frac{91}{6} - \frac{49}{4} = \frac{182-147}{12} = \frac{35}{12} \approx 2.9167$.

\textbf{By the rules:}
$E(S) = 3.5 + 3.5 = 7$.
$\operatorname{Var}(S) = \frac{35}{12} + \frac{35}{12} = \frac{35}{6} \approx 5.8333$.

\textbf{Directly from the distribution of $S$:}
\begin{align*}
E(S) &= \sum_{s=2}^{12} s\,P(S=s) \\
&= \frac{1}{36}\bigl(2\cdot1 + 3\cdot2 + 4\cdot3 + 5\cdot4 + 6\cdot5 + 7\cdot6 \\
&\qquad + 8\cdot5 + 9\cdot4 + 10\cdot3 + 11\cdot2 + 12\cdot1\bigr) \\
&= \frac{1}{36}(2+6+12+20+30+42+40+36+30+22+12) = \frac{252}{36} = 7. \quad\checkmark
\end{align*}
$E(S^2) = \frac{1}{36}\bigl(4\cdot1 + 9\cdot2 + 16\cdot3 + 25\cdot4 + 36\cdot5 + 49\cdot6 \\
+ 64\cdot5 + 81\cdot4 + 100\cdot3 + 121\cdot2 + 144\cdot1\bigr) = \frac{1974}{36} = \frac{329}{6} \approx 54.8333$.
$\operatorname{Var}(S) = \frac{329}{6} - 49 = \frac{329-294}{6} = \frac{35}{6}. \quad\checkmark$

Both methods agree perfectly.
\end{exoresolu}

\begin{exoresolu}[Two independent uniform variables on $\{0,1,2\}$]
Let $X_1$ and $X_2$ be independent, each with distribution $P(X=0)=0.2,\; P(X=1)=0.5,\; P(X=2)=0.3$.  Find the distribution of $S=X_1+X_2$, then compute $E(S)$ and $\operatorname{Var}(S)$ both directly and by the rules.

\tcblower
\textbf{Distribution of $S$ (convolution):} Possible sums $s=0,1,2,3,4$.
\begin{align*}
P(S=0) &= 0.2\times0.2 = 0.04,\\
P(S=1) &= 0.2\times0.5 + 0.5\times0.2 = 0.20,\\
P(S=2) &= 0.2\times0.3 + 0.5\times0.5 + 0.3\times0.2 = 0.06+0.25+0.06 = 0.37,\\
P(S=3) &= 0.5\times0.3 + 0.3\times0.5 = 0.30,\\
P(S=4) &= 0.3\times0.3 = 0.09.
\end{align*}
Check: $0.04+0.20+0.37+0.30+0.09 = 1.00$.

\textbf{Expectation and variance of one variable:}
$E(X) = 0\times0.2 + 1\times0.5 + 2\times0.3 = 1.1$.
$E(X^2) = 0^2\times0.2 + 1^2\times0.5 + 2^2\times0.3 = 0.5+1.2 = 1.7$.
$\operatorname{Var}(X) = 1.7 - 1.1^2 = 1.7 - 1.21 = 0.49$.

\textbf{By the rules:}
$E(S) = 1.1+1.1 = 2.2$.
$\operatorname{Var}(S) = 0.49+0.49 = 0.98$.

\textbf{Directly from $S$:}
$E(S) = 0\times0.04 + 1\times0.20 + 2\times0.37 + 3\times0.30 + 4\times0.09 = 0 + 0.20 + 0.74 + 0.90 + 0.36 = 2.20. \quad\checkmark$
$E(S^2) = 0^2\times0.04 + 1^2\times0.20 + 2^2\times0.37 + 3^2\times0.30 + 4^2\times0.09 = 0.20 + 1.48 + 2.70 + 1.44 = 5.82$.
$\operatorname{Var}(S) = 5.82 - (2.2)^2 = 5.82 - 4.84 = 0.98. \quad\checkmark$
\end{exoresolu}

\courssub{Extension to $n$ independent identically distributed variables}

The results extend naturally to the sum of $n$ independent random variables $X_1,\dots,X_n$ that are \emph{identically distributed} (i.i.d.) with common mean $\mu$ and variance $\sigma^2$.

\begin{aretenir}[Sum of $n$ i.i.d. random variables]
If $X_1,\dots,X_n$ are independent and identically distributed with $E(X_i)=\mu$ and $\operatorname{Var}(X_i)=\sigma^2$, then for $S_n = X_1+\cdots+X_n$:
\[
E(S_n) = n\mu, \qquad \operatorname{Var}(S_n) = n\sigma^2, \qquad \sigma_{S_n} = \sqrt{n}\,\sigma.
\]
This is a key stepping stone to the \cle{binomial distribution} (sum of $n$ independent Bernoulli trials) and to the \cle{Central Limit Theorem} (which says that for large $n$, $S_n$ is approximately normally distributed).
\end{aretenir}

\begin{exemplebox}[Total of 50 independent daily customer counts]
Suppose the daily number of customers at the café has mean $5$ and variance $5$ (Poisson(5)).  Over $50$ independent days, the total $S_{50}$ has $E(S_{50}) = 50\times5 = 250$ and $\operatorname{Var}(S_{50}) = 50\times5 = 250$, so $\sigma_{S_{50}} = \sqrt{250} \approx 15.81$.
\end{exemplebox}

\courssub{Summary of key points}

\begin{aretenir}[The distribution of $X_1+X_2$ — what you must know]
\begin{itemize}
\item Two discrete random variables $X_1,X_2$ are \textbf{independent} iff $P(X_1=x,X_2=y)=P(X_1=x)P(X_2=y)$ for all $x,y$.
\item To find the distribution of $S=X_1+X_2$ for independent variables: use a \textbf{two‑way table} (values of $X_1$ down, $X_2$ across, fill in sums and products of probabilities, then collect along diagonals) or the \textbf{convolution formula} $P(S=s)=\sum_x p_1(x)p_2(s-x)$.
\item \textbf{Expectation is always linear:} $E(X_1+X_2)=E(X_1)+E(X_2)$ (no independence needed).
\item \textbf{Variances add only under independence:} $\operatorname{Var}(X_1+X_2)=\operatorname{Var}(X_1)+\operatorname{Var}(X_2)$ if $X_1,X_2$ independent; otherwise a covariance term appears.
\item \textbf{Never add standard deviations.}  $\sigma_{X_1+X_2} = \sqrt{\sigma_{X_1}^2+\sigma_{X_2}^2}$ for independent variables.
\item For $n$ i.i.d. variables: $E(S_n)=n\mu$, $\operatorname{Var}(S_n)=n\sigma^2$.
\end{itemize}
\end{aretenir}

\begin{exercice}[Practice]
\begin{enumerate}
\item Two independent random variables $X$ and $Y$ have distributions:
$P(X=0)=0.3,\; P(X=1)=0.4,\; P(X=2)=0.3$;
$P(Y=1)=0.6,\; P(Y=2)=0.4$.
Find the probability distribution of $S=X+Y$.  Hence compute $E(S)$ and $\operatorname{Var}(S)$ and verify the rules $E(S)=E(X)+E(Y)$ and $\operatorname{Var}(S)=\operatorname{Var}(X)+\operatorname{Var}(Y)$.
\item A biased die has $P(6)=0.5$ and the other five faces equally likely.  Two such dice are rolled independently.  Find the distribution of the total score.  What is the most likely total (the mode of the sum)?
\item Let $X_1,X_2$ be independent with $E(X_1)=2,\; \operatorname{Var}(X_1)=3$ and $E(X_2)=5,\; \operatorname{Var}(X_2)=4$.  Find $E(3X_1-2X_2)$ and $\operatorname{Var}(3X_1-2X_2)$.  (Hint: use $E(aX+bY)=aE(X)+bE(Y)$ and $\operatorname{Var}(aX+bY)=a^2\operatorname{Var}(X)+b^2\operatorname{Var}(Y)$ for independent $X,Y$.)
\end{enumerate}
\end{exercice}

\begin{corrige}[Exercise 1]
\begin{enumerate}
\item \textbf{Distribution of $S=X+Y$:}
Possible sums: $1,2,3,4$.
\begin{align*}
P(S=1) &= P(X=0,Y=1) = 0.3\times0.6 = 0.18,\\
P(S=2) &= P(X=0,Y=2)+P(X=1,Y=1) = 0.3\times0.4 + 0.4\times0.6 = 0.12+0.24 = 0.36,\\
P(S=3) &= P(X=1,Y=2)+P(X=2,Y=1) = 0.4\times0.4 + 0.3\times0.6 = 0.16+0.18 = 0.34,\\
P(S=4) &= P(X=2,Y=2) = 0.3\times0.4 = 0.12.
\end{align*}
Check sum: $0.18+0.36+0.34+0.12=1.00$.

$E(X)=0\times0.3+1\times0.4+2\times0.3=1.0$; $E(X^2)=0+0.4+1.2=1.6$; $\operatorname{Var}(X)=1.6-1.0=0.6$.
$E(Y)=1\times0.6+2\times0.4=1.4$; $E(Y^2)=1\times0.6+4\times0.4=2.2$; $\operatorname{Var}(Y)=2.2-1.96=0.24$.

$E(S)=1\times0.18+2\times0.36+3\times0.34+4\times0.12 = 0.18+0.72+1.02+0.48 = 2.40 = 1.0+1.4$. $\checkmark$
$E(S^2)=1\times0.18+4\times0.36+9\times0.34+16\times0.12 = 0.18+1.44+3.06+1.92 = 6.60$.
$\operatorname{Var}(S)=6.60-2.40^2 = 6.60-5.76 = 0.84 = 0.6+0.24$. $\checkmark$

\item \textbf{Biased die:} $P(6)=0.5$, $P(1)=P(2)=P(3)=P(4)=P(5)=0.1$ each.
Two independent rolls.  The two‑way table is $6\times6$ but with unequal probabilities.
It is easier to use convolution or a table of products.  The possible totals are $2,\dots,12$.
The mode will be $12$ because $P(6,6)=0.5\times0.5=0.25$ is the largest single cell probability.  (Full table omitted for brevity.)

\item $E(3X_1-2X_2) = 3E(X_1)-2E(X_2) = 3\times2 - 2\times5 = 6-10 = -4$.
$\operatorname{Var}(3X_1-2X_2) = 3^2\operatorname{Var}(X_1) + (-2)^2\operatorname{Var}(X_2) = 9\times3 + 4\times4 = 27+16 = 43$.
\end{enumerate}
\end{corrige}

\coursec{Probability generating functions}

In the previous section we studied the distribution of the sum $X_1+X_2$ of two independent discrete random variables. We saw that the probability mass function of the sum is obtained by \emph{convolution} of the individual mass functions. While convolution is a fundamental concept, it can become computationally heavy, especially when we need the distribution of the sum of many independent variables, or when we want to compute moments (expectation, variance) quickly. 

The \emph{probability generating function} (p.g.f.) is a powerful algebraic tool that encodes the whole distribution of a discrete random variable into a single function of a dummy variable $t$. It turns convolution into simple multiplication, and moments into derivatives evaluated at $t=1$. This section introduces the p.g.f., its main properties, and its applications to standard distributions and to sums of independent variables.

\courssub{Definition and first properties}

\begin{definition}[Probability generating function]
Let $X$ be a discrete random variable taking values in the set of non‑negative integers $\{0,1,2,\dots\}$ with probability mass function $P(X=k)=p_k$ ($k=0,1,2,\dots$). The \textbf{probability generating function} (p.g.f.) of $X$ is the function $G_X(t)$ defined by
\[
G_X(t) = \mathbb{E}(t^X) = \sum_{k=0}^{\infty} p_k\, t^k,
\]
for all real $t$ for which the series converges. In practice, for bounded random variables or for the standard distributions we meet, the series converges at least for $|t|\le 1$.
\end{definition}

\begin{aretenir}[Key values of the p.g.f.]
\begin{itemize}
\item $G_X(1) = \sum_{k=0}^\infty p_k = 1$. This provides a quick check that a given function is a valid p.g.f. (it must equal $1$ at $t=1$).
\item $G_X(0) = p_0 = P(X=0)$.
\item The coefficient of $t^k$ in the power series expansion of $G_X(t)$ is exactly $p_k = P(X=k)$.
\end{itemize}
\end{aretenir}

\begin{propriete}[Recovering the distribution from the p.g.f.]
If $G_X(t) = \sum_{k=0}^\infty a_k t^k$ is the power series expansion of the p.g.f. about $t=0$, then
\[
P(X=k) = a_k = \frac{G_X^{(k)}(0)}{k!}, \qquad k=0,1,2,\dots
\]
where $G_X^{(k)}$ denotes the $k$-th derivative.
\end{propriete}

\begin{exemplebox}[Finding the distribution from a given p.g.f.]
Suppose a random variable $X$ has p.g.f. $G_X(t) = \frac{1}{4}(1+t)^2$ for $|t|\le 1$. Find the probability distribution of $X$.
\begin{align*}
G_X(t) &= \frac{1}{4}(1 + 2t + t^2) = \frac{1}{4} + \frac{1}{2}t + \frac{1}{4}t^2.
\end{align*}
Reading off the coefficients:
\[
P(X=0)=\frac{1}{4},\quad P(X=1)=\frac{1}{2},\quad P(X=2)=\frac{1}{4},\quad P(X=k)=0 \text{ for } k\ge 3.
\]
Check: $G_X(1)=\frac{1}{4}(2)^2=1$, valid.
\end{exemplebox}

\courssub{Moments from the probability generating function}

One of the most useful features of the p.g.f. is that it generates factorial moments by differentiation at $t=1$.

\begin{propriete}[Expectation and variance via the p.g.f.]
Let $X$ be a discrete random variable with p.g.f. $G_X(t)$. Then
\begin{align*}
\mathbb{E}(X) &= G_X'(1), \\
\mathbb{E}[X(X-1)] &= G_X''(1), \\
\operatorname{Var}(X) &= G_X''(1) + G_X'(1) - \bigl[G_X'(1)\bigr]^2.
\end{align*}
\end{propriete}

\begin{exemplebox}[Derivation of the formulas]
By definition, $G_X(t) = \sum_{k=0}^\infty p_k t^k$. Differentiating term by term (justified inside the radius of convergence):
\[
G_X'(t) = \sum_{k=1}^\infty k p_k t^{k-1}, \qquad
G_X''(t) = \sum_{k=2}^\infty k(k-1) p_k t^{k-2}.
\]
Evaluating at $t=1$:
\[
G_X'(1) = \sum_{k=1}^\infty k p_k = \mathbb{E}(X), \qquad
G_X''(1) = \sum_{k=2}^\infty k(k-1) p_k = \mathbb{E}[X(X-1)].
\]
Since $\mathbb{E}(X^2) = \mathbb{E}[X(X-1)] + \mathbb{E}(X)$, we obtain
\[
\operatorname{Var}(X) = \mathbb{E}(X^2) - [\mathbb{E}(X)]^2 = G_X''(1) + G_X'(1) - [G_X'(1)]^2.
\]
\end{exemplebox}

\begin{attention}[Common mistakes with derivatives]
\begin{itemize}
\item Always evaluate the derivatives \textbf{at $t=1$}, not at $t=0$. $G'(0)$ gives $p_1$, not the mean.
\item $G''(1)$ gives $\mathbb{E}[X(X-1)]$, \textbf{not} $\mathbb{E}(X^2)$. Remember to add $G'(1)$ to get $\mathbb{E}(X^2)$.
\item The formulas are valid only if the derivatives exist at $t=1$ (i.e. the mean and variance are finite).
\end{itemize}
\end{attention}

\begin{exoresolu}[Mean and variance from a p.g.f.]
A discrete random variable $X$ has p.g.f. $G_X(t) = \frac{2}{5} + \frac{3}{5}t^2$. Find $\mathbb{E}(X)$ and $\operatorname{Var}(X)$.
\tcblower
First, check $G_X(1) = \frac{2}{5}+\frac{3}{5}=1$, valid.
\[
G_X'(t) = \frac{6}{5}t, \quad G_X'(1) = \frac{6}{5} = 1.2.
\]
\[
G_X''(t) = \frac{6}{5}, \quad G_X''(1) = \frac{6}{5} = 1.2.
\]
Hence $\mathbb{E}(X) = 1.2$ and
\[
\operatorname{Var}(X) = G_X''(1) + G_X'(1) - [G_X'(1)]^2 = 1.2 + 1.2 - (1.2)^2 = 2.4 - 1.44 = 0.96.
\]
We can verify by reading the distribution: $P(X=0)=\frac{2}{5}$, $P(X=2)=\frac{3}{5}$. Then $\mathbb{E}(X)=0\cdot\frac{2}{5}+2\cdot\frac{3}{5}=1.2$, $\mathbb{E}(X^2)=4\cdot\frac{3}{5}=2.4$, $\operatorname{Var}=2.4-1.44=0.96$.
\end{exoresolu}

\courssub{Probability generating function of a sum of independent variables}

This is the property that makes the p.g.f. indispensable for sums of independent random variables.

\begin{propriete}[P.g.f. of a sum of independent variables]
If $X_1$ and $X_2$ are \textbf{independent} discrete random variables taking non‑negative integer values, then the p.g.f. of their sum $S = X_1+X_2$ is the product of their individual p.g.f.s:
\[
G_{X_1+X_2}(t) = G_{X_1}(t) \cdot G_{X_2}(t).
\]
More generally, for $n$ independent variables $X_1,\dots,X_n$,
\[
G_{X_1+\cdots+X_n}(t) = \prod_{i=1}^n G_{X_i}(t).
\]
\end{propriete}

\begin{exemplebox}[Proof]
By independence, for any $t$,
\[
G_{X_1+X_2}(t) = \mathbb{E}\bigl(t^{X_1+X_2}\bigr) = \mathbb{E}\bigl(t^{X_1} t^{X_2}\bigr) = \mathbb{E}(t^{X_1})\,\mathbb{E}(t^{X_2}) = G_{X_1}(t)\,G_{X_2}(t).
\]
\end{exemplebox}

\begin{aretenir}[Link with convolution]
The product of the p.g.f.s corresponds exactly to the convolution of the probability mass functions. If $G_{X_1}(t)=\sum a_k t^k$ and $G_{X_2}(t)=\sum b_k t^k$, then the coefficient of $t^n$ in the product is $\sum_{k=0}^n a_k b_{n-k}$, which is precisely $P(X_1+X_2=n)$ obtained by convolution in Section 5.
\end{aretenir}

\begin{exoresolu}[Sum of two independent variables using p.g.f.]
Let $X$ and $Y$ be independent with distributions:
\[
P(X=0)=0.4,\; P(X=1)=0.6; \qquad P(Y=0)=0.3,\; P(Y=1)=0.5,\; P(Y=2)=0.2.
\]
Find the distribution of $S=X+Y$ using p.g.f.s.
\tcblower
First, write the p.g.f.s:
\[
G_X(t) = 0.4 + 0.6t, \qquad G_Y(t) = 0.3 + 0.5t + 0.2t^2.
\]
Then
\[
G_S(t) = G_X(t)G_Y(t) = (0.4+0.6t)(0.3+0.5t+0.2t^2).
\]
Multiply out:
\begin{align*}
G_S(t) &= 0.4(0.3+0.5t+0.2t^2) + 0.6t(0.3+0.5t+0.2t^2) \\
&= 0.12 + 0.20t + 0.08t^2 + 0.18t + 0.30t^2 + 0.12t^3 \\
&= 0.12 + 0.38t + 0.38t^2 + 0.12t^3.
\end{align*}
Hence the distribution of $S$ is:
\[
P(S=0)=0.12,\; P(S=1)=0.38,\; P(S=2)=0.38,\; P(S=3)=0.12.
\]
(Check: sum $=1$.)
\end{exoresolu}

\courssub{P.g.f. of standard distributions}

We now derive the p.g.f. for the fundamental distributions that appear repeatedly in the syllabus and in examinations.

\begin{propriete}[P.g.f. of the uniform distribution on $\{1,2,\dots,n\}$]
If $X$ is uniformly distributed on $\{1,2,\dots,n\}$, then $P(X=k)=\frac{1}{n}$ for $k=1,\dots,n$ and
\[
G_X(t) = \frac{1}{n}\sum_{k=1}^n t^k = \frac{t(1-t^n)}{n(1-t)}, \quad t\ne 1, \qquad G_X(1)=1.
\]
\end{propriete}

\begin{propriete}[P.g.f. of a Bernoulli distribution]
Let $X$ take values $0$ and $1$ with $P(X=1)=p$ and $P(X=0)=1-p$ ($0<p<1$). This is the \textbf{Bernoulli distribution} with parameter $p$. Its p.g.f. is
\[
G_X(t) = (1-p) + p t = 1 - p + p t.
\]
\end{propriete}

\begin{propriete}[P.g.f. of the Binomial distribution]
If $X \sim \mathrm{Bin}(n,p)$, i.e. $X$ is the sum of $n$ independent Bernoulli($p$) variables, then by the product property
\[
G_X(t) = (1-p+pt)^n.
\]
Expanding by the binomial theorem gives the coefficients $\binom{n}{k}p^k(1-p)^{n-k}$, which are exactly the Binomial probabilities.
\end{propriete}

\begin{propriete}[P.g.f. of the Poisson distribution]
If $X \sim \mathrm{Po}(\lambda)$ with $\lambda>0$, then $P(X=k)=e^{-\lambda}\frac{\lambda^k}{k!}$ for $k=0,1,2,\dots$ and
\[
G_X(t) = \sum_{k=0}^\infty e^{-\lambda}\frac{\lambda^k}{k!} t^k = e^{-\lambda} \sum_{k=0}^\infty \frac{(\lambda t)^k}{k!} = e^{-\lambda} e^{\lambda t} = e^{\lambda(t-1)}.
\]
\end{propriete}

\begin{propriete}[P.g.f. of the Geometric distribution]
If $X \sim \mathrm{Geo}(p)$ with $0<p<1$, where $X$ is the number of failures before the first success (so $P(X=k)=p(1-p)^k$ for $k=0,1,2,\dots$), then
\[
G_X(t) = \sum_{k=0}^\infty p(1-p)^k t^k = \frac{p}{1-(1-p)t}, \quad |t| < \frac{1}{1-p}.
\]
(If the geometric distribution is defined as the number of trials until the first success, taking values $1,2,\dots$, the p.g.f. is $\frac{pt}{1-(1-p)t}$.)
\end{propriete}

\begin{savaistu}[Why the Poisson p.g.f. is elegant]
The Poisson p.g.f. $G(t)=e^{\lambda(t-1)}$ makes it obvious that the sum of independent Poisson variables is Poisson: if $X\sim\mathrm{Po}(\lambda_1)$ and $Y\sim\mathrm{Po}(\lambda_2)$ are independent, then $G_{X+Y}(t)=e^{\lambda_1(t-1)}e^{\lambda_2(t-1)}=e^{(\lambda_1+\lambda_2)(t-1)}$, so $X+Y\sim\mathrm{Po}(\lambda_1+\lambda_2)$. The mean and variance are also immediate: $G'(t)=\lambda e^{\lambda(t-1)}$, so $\mathbb{E}(X)=\lambda$; $G''(t)=\lambda^2 e^{\lambda(t-1)}$, so $\operatorname{Var}(X)=\lambda^2+\lambda-\lambda^2=\lambda$.
\end{savaistu}

\begin{exoresolu}[Identifying a distribution from its p.g.f.]
A random variable $X$ has p.g.f. $G_X(t) = \frac{1}{27}(1+2t)^3$. Identify the distribution of $X$ and state its mean and variance.
\tcblower
First, verify $G_X(1)=\frac{1}{27}(3)^3=1$, so it is a valid p.g.f.
Expand using the binomial theorem:
\[
(1+2t)^3 = 1 + 6t + 12t^2 + 8t^3.
\]
So $G_X(t) = \frac{1}{27} + \frac{6}{27}t + \frac{12}{27}t^2 + \frac{8}{27}t^3 = \frac{1}{27} + \frac{2}{9}t + \frac{4}{9}t^2 + \frac{8}{27}t^3$.
This matches the Binomial distribution with $n=3$ and $p=\frac{2}{3}$:
$P(X=k)=\binom{3}{k}\left(\frac{2}{3}\right)^k\left(\frac{1}{3}\right)^{3-k}$.
For $k=0$: $(1/3)^3=1/27$; $k=1$: $3\cdot(2/3)\cdot(1/9)=6/27$; $k=2$: $3\cdot(4/9)\cdot(1/3)=12/27$; $k=3$: $8/27$.
Thus $X\sim \mathrm{Bin}(3,\frac{2}{3})$.
Mean: $\mathbb{E}(X)=np=3\cdot\frac{2}{3}=2$. Variance: $np(1-p)=3\cdot\frac{2}{3}\cdot\frac{1}{3}=\frac{2}{3}$.
Using p.g.f.: $G_X'(t)=\frac{1}{27}\cdot 3(1+2t)^2\cdot 2 = \frac{2}{9}(1+2t)^2$, $G_X'(1)=\frac{2}{9}\cdot 9=2$. $G_X''(t)=\frac{2}{9}\cdot 2(1+2t)\cdot 2 = \frac{8}{9}(1+2t)$, $G_X''(1)=\frac{8}{9}\cdot 3=\frac{8}{3}$. $\operatorname{Var}= \frac{8}{3}+2-4=\frac{2}{3}$. Consistent.
\end{exoresolu}

\courssub{Illustrations}

\begin{popfigure}
\begin{tikzpicture}
\begin{axis}[
    width=0.9\linewidth,
    height=7cm,
    xlabel={$t$},
    ylabel={$G(t)$},
    xmin=0, xmax=1.05,
    ymin=0, ymax=1.1,
    xtick={0,0.5,1},
    ytick={0,0.5,1},
    grid=major,
    grid style={popline},
    axis lines=middle,
    enlargelimits={abs=0.05},
    legend style={at={(0.98,0.02)}, anchor=south east, fill=white, draw=popline},
    every axis plot/.append style={thick}
]
% Bernoulli(p=0.3): G(t) = 0.7 + 0.3t
\addplot[domain=0:1, samples=2, popblue] {0.7 + 0.3*x};
\addlegendentry{Bernoulli($p=0.3$): $G(t)=0.7+0.3t$}
% Uniform on {1,2,3,4}: G(t) = (t+t^2+t^3+t^4)/4
\addplot[domain=0:1, samples=100, poporange] {(x+x^2+x^3+x^4)/4};
\addlegendentry{Uniform on $\{1,2,3,4\}$: $G(t)=\frac{t+t^2+t^3+t^4}{4}$}
% Point (1,1)
\addplot[only marks, mark=*, mark size=2, popdark] coordinates {(1,1)};
% Tangent lines at t=1 to show slopes = E(X)
% Bernoulli: slope = 0.3
\addplot[domain=0.6:1, samples=2, popblue, dashed] {0.7+0.3*x}; % same line
% Uniform: derivative at 1 = (1+2+3+4)/4 = 2.5, line: y-1 = 2.5(x-1) => y = 2.5x -1.5
\addplot[domain=0.6:1, samples=2, poporange, dashed] {2.5*x - 1.5};
\end{axis}
\end{tikzpicture}
\legende{Probability generating functions of two distributions on $[0,1]$. Both pass through $(1,1)$. The slope at $t=1$ equals $\mathbb{E}(X)$: $0.3$ for Bernoulli($0.3$) and $2.5$ for Uniform on $\{1,2,3,4\}$.}
\end{popfigure}

\begin{popfigure}
\begin{tikzpicture}[
    box/.style={draw=popdark, fill=popblueL, rounded corners, minimum width=2.5cm, minimum height=1cm, align=center, font=\small},
    arrow/.style={->, thick, popdark},
    start/.style={draw=popgreen, fill=popgreenL, rounded corners, minimum width=2.5cm, minimum height=1cm, align=center, font=\small},
    end/.style={draw=poporange, fill=poporange!20, rounded corners, minimum width=2.5cm, minimum height=1cm, align=center, font=\small}
]
\node[start] (exp) at (0,0) {Random Experiment};
\node[box] (rv) at (4,0) {Random Variable $X$};
\node[box, align=center] (dist) at (8,0){Probability Distribution\\$P(X=k)=p_k$};
\node[box, align=center] (pmf) at (8,-2){Probability Mass Function\\$f(k)=p_k$};
\node[box, align=center] (cdf) at (4,-2){Cumulative Distribution\\Function $F(x)$};
\node[box, align=center] (mom) at (4,-4){Expectation $\mathbb{E}(X)$\\Variance $\operatorname{Var}(X)$};
\node[box, align=center] (pgf) at (8,-4){Probability Generating\\Function $G_X(t)$};
\node[box, align=center] (sum) at (12,-2){Sum of Independent\\Variables $X_1+X_2$};
\node[end, align=center] (appl) at (16,-2){Applications:\\Modelling, Inference};

\draw[arrow] (exp) -- (rv);
\draw[arrow] (rv) -- (dist);
\draw[arrow] (dist) -- (pmf);
\draw[arrow] (dist) -- (cdf);
\draw[arrow] (pmf) -- (mom);
\draw[arrow] (pmf) -- (pgf);
\draw[arrow] (pgf) -- (sum);
\draw[arrow] (sum) -- (appl);
\draw[arrow] (mom) -- (appl);
\draw[arrow] (cdf) -- (appl);
\end{tikzpicture}
\legende{Flowchart summarising the chapter: from experiment to random variable, distribution, and the main tools (PMF, CDF, moments, p.g.f.) leading to applications.}
\end{popfigure}

\courssub{Method: working with probability generating functions}

\begin{methode}[Using the p.g.f. to solve problems]
\begin{enumerate}
\item \textbf{Identify the random variable} and its possible values (non‑negative integers).
\item \textbf{Write down the p.g.f.} $G_X(t) = \sum p_k t^k$ from the given distribution, or recognise a standard form.
\item \textbf{Check validity:} $G_X(1)$ must equal $1$.
\item \textbf{To find probabilities:} expand $G_X(t)$ as a power series in $t$; the coefficient of $t^k$ is $P(X=k)$. Alternatively, compute $P(X=k)=\frac{G_X^{(k)}(0)}{k!}$.
\item \textbf{To find mean and variance:} compute $G_X'(t)$ and $G_X''(t)$, evaluate at $t=1$, then use
      $\mathbb{E}(X)=G_X'(1)$,
      $\operatorname{Var}(X)=G_X''(1)+G_X'(1)-[G_X'(1)]^2$.
\item \textbf{For sums of independent variables:} multiply their p.g.f.s. The product is the p.g.f. of the sum. Expand the product to recover the distribution of the sum, or differentiate the product to get its mean and variance.
\end{enumerate}
\end{methode}

\begin{exoresolu}[Synthesis exercise: p.g.f., moments, and sum]
Two independent random variables $X$ and $Y$ have p.g.f.s
\[
G_X(t) = \frac{1}{3} + \frac{2}{3}t, \qquad G_Y(t) = \frac{1}{4}(1+t)^2.
\]
\begin{enumerate}
\item Identify the distributions of $X$ and $Y$.
\item Find $\mathbb{E}(X)$, $\operatorname{Var}(X)$, $\mathbb{E}(Y)$, $\operatorname{Var}(Y)$.
\item Find the p.g.f. of $S=X+Y$ and hence the distribution of $S$.
\item Compute $\mathbb{E}(S)$ and $\operatorname{Var}(S)$ directly from $G_S(t)$ and verify they equal $\mathbb{E}(X)+\mathbb{E}(Y)$ and $\operatorname{Var}(X)+\operatorname{Var}(Y)$.
\end{enumerate}
\tcblower
\textbf{(a)} $G_X(t)=\frac{1}{3}+\frac{2}{3}t$ is a Bernoulli distribution with $p=\frac{2}{3}$: $P(X=0)=\frac{1}{3}$, $P(X=1)=\frac{2}{3}$.
$G_Y(t)=\frac{1}{4}(1+2t+t^2)=\frac{1}{4}+\frac{1}{2}t+\frac{1}{4}t^2$. So $Y$ takes values $0,1,2$ with probabilities $\frac{1}{4},\frac{1}{2},\frac{1}{4}$.

\textbf{(b)} For $X$: $G_X'(t)=\frac{2}{3}$, $G_X'(1)=\frac{2}{3}=\mathbb{E}(X)$. $G_X''(t)=0$, so $\operatorname{Var}(X)=0+\frac{2}{3}-(\frac{2}{3})^2=\frac{2}{3}-\frac{4}{9}=\frac{2}{9}$.
For $Y$: $G_Y'(t)=\frac{1}{2}+\frac{1}{2}t$, $G_Y'(1)=1=\mathbb{E}(Y)$. $G_Y''(t)=\frac{1}{2}$, $G_Y''(1)=\frac{1}{2}$. $\operatorname{Var}(Y)=\frac{1}{2}+1-1^2=\frac{1}{2}$.

\textbf{(c)} $G_S(t)=G_X(t)G_Y(t)=\bigl(\frac{1}{3}+\frac{2}{3}t\bigr)\bigl(\frac{1}{4}+\frac{1}{2}t+\frac{1}{4}t^2\bigr)$.
Write with common denominator 12:
$G_X(t) = \frac{4}{12} + \frac{8}{12}t$, \quad $G_Y(t) = \frac{3}{12} + \frac{6}{12}t + \frac{3}{12}t^2$.
Product:
\[
G_S(t) = \frac{4}{12}\cdot\frac{3}{12} + \left(\frac{4}{12}\cdot\frac{6}{12}+\frac{8}{12}\cdot\frac{3}{12}\right)t + \left(\frac{4}{12}\cdot\frac{3}{12}+\frac{8}{12}\cdot\frac{6}{12}\right)t^2 + \frac{8}{12}\cdot\frac{3}{12}t^3
\]
\[
= \frac{12}{144} + \frac{24+24}{144}t + \frac{12+48}{144}t^2 + \frac{24}{144}t^3
= \frac{1}{12} + \frac{48}{144}t + \frac{60}{144}t^2 + \frac{24}{144}t^3
= \frac{1}{12} + \frac{1}{3}t + \frac{5}{12}t^2 + \frac{1}{6}t^3.
\]
So $P(S=0)=\frac{1}{12}$, $P(S=1)=\frac{1}{4}$, $P(S=2)=\frac{5}{12}$, $P(S=3)=\frac{1}{6}$.

\textbf{(d)} $G_S'(t) = \frac{1}{3} + \frac{10}{12}t + \frac{3}{6}t^2 = \frac{1}{3} + \frac{5}{6}t + \frac{1}{2}t^2$.
$G_S'(1) = \frac{1}{3}+\frac{5}{6}+\frac{1}{2} = \frac{2+5+3}{6} = \frac{10}{6} = \frac{5}{3}$. Matches $\mathbb{E}(X)+\mathbb{E}(Y)=\frac{2}{3}+1=\frac{5}{3}$.
$G_S''(t) = \frac{5}{6} + t$, $G_S''(1) = \frac{5}{6}+1 = \frac{11}{6}$.
$\operatorname{Var}(S) = \frac{11}{6} + \frac{5}{3} - (\frac{5}{3})^2 = \frac{11}{6}+\frac{10}{6}-\frac{25}{9} = \frac{21}{6}-\frac{25}{9} = \frac{7}{2}-\frac{25}{9} = \frac{63-50}{18} = \frac{13}{18}$.
$\operatorname{Var}(X)+\operatorname{Var}(Y) = \frac{2}{9}+\frac{1}{2} = \frac{4}{18}+\frac{9}{18}=\frac{13}{18}$. Matches.
\end{exoresolu}

\courssub{Exercises}

\begin{exercice}[GCE-style practice]
\begin{enumerate}
\item A discrete random variable $X$ has probability generating function $G_X(t) = \frac{1}{36}(1+2t)^2(1+3t)$. 
      \begin{enumerate}
      \item Show that $G_X(t)$ is a valid p.g.f.
      \item Find the probability distribution of $X$.
      \item Calculate $\mathbb{E}(X)$ and $\operatorname{Var}(X)$.
      \end{enumerate}
\item Let $X$ and $Y$ be independent random variables with p.g.f.s $G_X(t)=0.2+0.8t$ and $G_Y(t)=0.5+0.3t+0.2t^2$.
      \begin{enumerate}
      \item Write down the distributions of $X$ and $Y$.
      \item Find the p.g.f. of $S=X+Y$ and hence the distribution of $S$.
      \item Find $\mathbb{E}(S)$ and $\operatorname{Var}(S)$ using $G_S(t)$.
      \end{enumerate}
\item The random variable $X$ has p.g.f. $G_X(t) = \frac{t(1-t^4)}{4(1-t)}$ for $t\ne 1$, $G_X(1)=1$.
      \begin{enumerate}
      \item Identify the distribution of $X$.
      \item Find $\mathbb{E}(X)$ and $\operatorname{Var}(X)$ by differentiating $G_X(t)$.
      \end{enumerate}
\item A random variable $X$ has p.g.f. $G_X(t) = \frac{1}{6}(1+t+t^2+t^3+t^4+t^5)$. 
      \begin{enumerate}
      \item What is the distribution of $X$?
      \item Two independent observations $X_1$ and $X_2$ are taken. Find the p.g.f. of $X_1+X_2$ and hence the probability that $X_1+X_2 = 7$.
      \end{enumerate}
\item If $X \sim \mathrm{Po}(\lambda_1)$ and $Y \sim \mathrm{Po}(\lambda_2)$ are independent, use p.g.f.s to show that $X+Y \sim \mathrm{Po}(\lambda_1+\lambda_2)$ and find $\mathbb{E}(X+Y)$ and $\operatorname{Var}(X+Y)$.
\end{enumerate}
\end{exercice}

\begin{corrige}[Exercise 1]
\begin{enumerate}
\item $G_X(1)=\frac{1}{36}(3)^2(4)=\frac{36}{36}=1$. Valid p.g.f.
\item Expand: $(1+2t)^2(1+3t) = (1+4t+4t^2)(1+3t) = 1+3t+4t+12t^2+4t^2+12t^3 = 1+7t+16t^2+12t^3$.
So $G_X(t) = \frac{1}{36} + \frac{7}{36}t + \frac{16}{36}t^2 + \frac{12}{36}t^3 = \frac{1}{36} + \frac{7}{36}t + \frac{4}{9}t^2 + \frac{1}{3}t^3$.
Distribution: $P(X=0)=\frac{1}{36}$, $P(X=1)=\frac{7}{36}$, $P(X=2)=\frac{4}{9}$, $P(X=3)=\frac{1}{3}$.
\item $G_X'(t) = \frac{1}{36}(7+32t+36t^2)$, $G_X'(1)=\frac{75}{36}=\frac{25}{12}$.
$G_X''(t) = \frac{1}{36}(32+72t)$, $G_X''(1)=\frac{104}{36}=\frac{26}{9}$.
$\operatorname{Var}(X) = \frac{26}{9}+\frac{25}{12}-(\frac{25}{12})^2 = \frac{104+75}{36}-\frac{625}{144} = \frac{179}{36}-\frac{625}{144} = \frac{716-625}{144}=\frac{91}{144}$.
\end{enumerate}
\end{corrige}

\begin{corrige}[Exercise 2]
\begin{enumerate}
\item $X\sim \mathrm{Bernoulli}(0.8)$: $P(X=0)=0.2$, $P(X=1)=0.8$.
$Y$: $P(Y=0)=0.5$, $P(Y=1)=0.3$, $P(Y=2)=0.2$.
\item $G_S(t) = (0.2+0.8t)(0.5+0.3t+0.2t^2) = 0.1 + 0.06t + 0.04t^2 + 0.4t + 0.24t^2 + 0.16t^3 = 0.1 + 0.46t + 0.28t^2 + 0.16t^3$.
Distribution: $P(S=0)=0.1$, $P(S=1)=0.46$, $P(S=2)=0.28$, $P(S=3)=0.16$.
\item $G_S'(t)=0.46+0.56t+0.48t^2$, $G_S'(1)=1.5$. $G_S''(t)=0.56+0.96t$, $G_S''(1)=1.52$.
$\mathbb{E}(S)=1.5$, $\operatorname{Var}(S)=1.52+1.5-2.25=0.77$.
Check: $\mathbb{E}(X)=0.8$, $\operatorname{Var}(X)=0.8\cdot0.2=0.16$. $\mathbb{E}(Y)=0\cdot0.5+1\cdot0.3+2\cdot0.2=0.7$, $\mathbb{E}(Y^2)=0.3+0.8=1.1$, $\operatorname{Var}(Y)=1.1-0.49=0.61$. Sum: mean $1.5$, variance $0.77$. Correct.
\end{enumerate}
\end{corrige}

\begin{corrige}[Exercise 3]
\begin{enumerate}
\item $G_X(t) = \frac{t(1-t^4)}{4(1-t)} = \frac{t+t^2+t^3+t^4}{4}$. This is the uniform distribution on $\{1,2,3,4\}$.
\item $G_X(t) = \frac{1}{4}(t+t^2+t^3+t^4)$.
$G_X'(t) = \frac{1}{4}(1+2t+3t^2+4t^3)$, $G_X'(1)=\frac{10}{4}=2.5$.
$G_X''(t) = \frac{1}{4}(2+6t+12t^2)$, $G_X''(1)=\frac{20}{4}=5$.
$\operatorname{Var}(X) = 5 + 2.5 - (2.5)^2 = 7.5 - 6.25 = 1.25$.
Check: Uniform on $\{1,2,3,4\}$: mean $2.5$, variance $\frac{(4-1+1)^2-1}{12}=\frac{16-1}{12}=1.25$. Correct.
\end{enumerate}
\end{corrige}

\begin{corrige}[Exercise 4]
\begin{enumerate}
\item $G_X(t) = \frac{1}{6}\sum_{k=0}^5 t^k$. This is the uniform distribution on $\{0,1,2,3,4,5\}$.
\item $G_{X_1+X_2}(t) = [G_X(t)]^2 = \frac{1}{36}(1+t+t^2+t^3+t^4+t^5)^2$.
We need the coefficient of $t^7$. The pairs $(i,j)$ with $i+j=7$, $0\le i,j\le 5$: $(2,5),(3,4),(4,3),(5,2)$. That's 4 pairs. Each pair has probability $\frac{1}{6}\cdot\frac{1}{6}=\frac{1}{36}$. So $P(X_1+X_2=7)=\frac{4}{36}=\frac{1}{9}$.
\end{enumerate}
\end{corrige}

\begin{corrige}[Exercise 5]
$G_X(t)=e^{\lambda_1(t-1)}$, $G_Y(t)=e^{\lambda_2(t-1)}$.
By independence, $G_{X+Y}(t)=G_X(t)G_Y(t)=e^{\lambda_1(t-1)}e^{\lambda_2(t-1)}=e^{(\lambda_1+\lambda_2)(t-1)}$, which is the p.g.f. of $\mathrm{Po}(\lambda_1+\lambda_2)$. Hence $X+Y\sim\mathrm{Po}(\lambda_1+\lambda_2)$.
$\mathbb{E}(X+Y)=\lambda_1+\lambda_2$, $\operatorname{Var}(X+Y)=\lambda_1+\lambda_2$.
\end{corrige}

\courssub{Chapter summary: key ideas}

\begin{aretenir}[What you must remember from this chapter]
\begin{itemize}
\item A \textbf{random variable} assigns a number to each outcome of a random experiment. Discrete random variables take countable values.
\item The \textbf{probability mass function} $f(k)=P(X=k)$ and the \textbf{cumulative distribution function} $F(x)=P(X\le x)$ describe the distribution.
\item \textbf{Expectation} $\mathbb{E}(X)=\sum k f(k)$ measures the centre; \textbf{variance} $\operatorname{Var}(X)=\mathbb{E}[(X-\mathbb{E}(X))^2]$ measures spread.
\item The \textbf{mode} is the value with highest probability; the \textbf{median} is any $m$ with $P(X\le m)\ge\frac12$ and $P(X\ge m)\ge\frac12$.
\item For independent $X_1,X_2$, the distribution of $X_1+X_2$ is the \textbf{convolution} of their distributions.
\item The \textbf{probability generating function} $G_X(t)=\mathbb{E}(t^X)=\sum p_k t^k$ encodes the whole distribution. 
      \begin{itemize}
      \item $G_X(1)=1$, $G_X(0)=P(X=0)$.
      \item $P(X=k)$ = coefficient of $t^k$ in $G_X(t)$.
      \item $\mathbb{E}(X)=G_X'(1)$, $\operatorname{Var}(X)=G_X''(1)+G_X'(1)-[G_X'(1)]^2$.
      \item If $X_1,X_2$ independent, $G_{X_1+X_2}(t)=G_{X_1}(t)G_{X_2}(t)$.
      \end{itemize}
\item Standard p.g.f.s: 
      \begin{itemize}
      \item Bernoulli($p$): $1-p+pt$
      \item Binomial($n,p$): $(1-p+pt)^n$
      \item Poisson($\lambda$): $e^{\lambda(t-1)}$
      \item Geometric($p$) (failures before first success): $\frac{p}{1-(1-p)t}$
      \item Uniform on $\{1,\dots,n\}$: $\frac{t(1-t^n)}{n(1-t)}$
      \end{itemize}
\end{itemize}
\end{aretenir}

\begin{attention}[Examination traps]
\begin{itemize}
\item \textbf{Forgetting the domain:} The p.g.f. is defined for discrete random variables taking \textbf{non‑negative integer} values. If a variable takes negative values, the p.g.f. as defined here does not apply directly.
\item \textbf{Confusing $G'(1)$ and $G'(0)$:} $G'(0)=p_1$, not the mean. Always evaluate derivatives at $t=1$ for moments.
\item \textbf{Mixing up $G''(1)$ and $\mathbb{E}(X^2)$:} $G''(1)=\mathbb{E}[X(X-1)]$. You must add $G'(1)$ to get $\mathbb{E}(X^2)$.
\item \textbf{Assuming independence without justification:} The product formula $G_{X+Y}=G_X G_Y$ holds \textbf{only} if $X$ and $Y$ are independent.
\item \textbf{Not checking $G(1)=1$:} Always verify this when given a candidate p.g.f.; if it is not $1$, the function is not a valid p.g.f.
\end{itemize}
\end{attention}

This concludes the chapter on Random Variables and Discrete Random Variables. You now have a complete toolkit: probability mass functions, cumulative distribution functions, expectation and variance, mode and median, convolution for sums, and probability generating functions. Mastery of these concepts is essential for the GCE Advanced Level examination and for further studies in probability and statistics. Practice the exercises thoroughly, and always check your answers for consistency (probabilities sum to $1$, variance is non‑negative, etc.). Good luck!

\newpage

\coursec{Integration activity}

\courssub{Aims of the activity}
This integration activity is designed to consolidate all the key concepts of the chapter on discrete random variables. By working through a realistic business scenario, you will:
\begin{itemize}
    \item Build and verify an empirical probability distribution from raw data.
    \item Compute and interpret the expectation, variance, standard deviation, mode and median.
    \item Construct and use the cumulative distribution function (CDF) to calculate probabilities.
    \item Derive the distribution of the sum of two independent random variables and check the addition rules for expectation and variance.
    \item Manipulate the probability generating function (p.g.f.) to recover the mean.
    \item Apply these tools to make and justify a concrete stocking decision under uncertainty.
\end{itemize}
The activity mirrors the type of multi-step, context-rich problems that appear in the GCE Advanced Level Pure Mathematics with Statistics examination.

\courssub{Situation}
\textbf{The Yaoundé phone-credit trader's stocking problem}

Madame Ngo Bakola runs a small phone-credit stall at Mokolo Market in Yaoundé. Every morning she buys a stock of recharge cards (MTN, Orange, Nexttel) from a wholesaler and sells them throughout the day. Each card yields a profit of \SI{200}{FCFA} if sold, but any unsold card at the end of the day represents a loss of \SI{100}{FCFA} (the card expires or loses value). She wants to decide on an optimal fixed daily stocking level \(S\) (number of cards to buy each morning) to maximise her expected daily profit.

To base her decision on data, she has recorded the number of customers \(X\) who bought at least one card each day over the past 50 trading days. The raw data are summarised in the frequency table below.

\begin{center}
\begin{tabular}{|c|cccccccccccccccc|}
\hline
\rowcolor{popblueL}
\(x\) (customers) & 0 & 1 & 2 & 3 & 4 & 5 & 6 & 7 & 8 & 9 & 10 & 11 & 12 & 13 & 14 & 15 \\ \hline
Frequency \(f\) & 1 & 1 & 2 & 3 & 4 & 5 & 6 & 7 & 6 & 5 & 4 & 3 & 2 & 1 & 0 & 0 \\ \hline
\end{tabular}
\end{center}

Assume that the daily demand \(X\) follows this empirical distribution and that demands on different days are independent.

\courssub{Tasks}
Work in groups of 3–4. Present your final recommendation with full mathematical justification.

\begin{enumerate}
    \item \textbf{Empirical distribution:} Construct the probability mass function (p.m.f.) \(P(X=x)\) from the frequencies. Verify that \(\sum P(X=x)=1\).
    \item \textbf{Summary statistics:} Compute the expectation \(E(X)\), the variance \(\operatorname{Var}(X)\), the standard deviation \(\sigma_X\), the mode and the median of the daily demand.
    \item \textbf{Cumulative distribution function:} Construct the CDF \(F(x)=P(X\le x)\) for all real \(x\). Sketch its graph. Use the CDF to find \(P(X\le 10)\) and \(P(5<X\le 12)\).
    \item \textbf{Two-day total demand:} Let \(X_1\) and \(X_2\) be the demands on two successive days (independent, identically distributed as \(X\)). Find the probability distribution of \(Y=X_1+X_2\) (you may present it as a table for the possible values 0 to 26). Verify that \(E(Y)=2E(X)\) and \(\operatorname{Var}(Y)=2\operatorname{Var}(X)\).
    \item \textbf{Probability generating function:} Write down the p.g.f. \(G_X(t)=E(t^X)\). Differentiate it and evaluate at \(t=1\) to recover \(E(X)\).
    \item \textbf{Stocking decision:} If Madame Ngo Bakola stocks \(S\) cards each day, her profit is 
    \[
    \Pi(S)=200\min(X,S)-100\max(S-X,0).
    \]
    Compute the expected profit \(E[\Pi(S)]\) for \(S=6,7,8,9\). Which stocking level maximises the expected profit? Defend your recommendation, commenting on the risk (variance) associated with each choice.
\end{enumerate}

\courssub{Model answer}
\begin{exoresolu}[Model Answer]
\textbf{Task 1: Empirical probability distribution}

The total number of days is \(N=50\). The p.m.f. is \(p(x)=f(x)/50\).

\begin{center}
\begin{tabular}{|c|cccccccccccccccc|}
\hline
\rowcolor{popblueL}
\(x\) & 0 & 1 & 2 & 3 & 4 & 5 & 6 & 7 & 8 & 9 & 10 & 11 & 12 & 13 & 14 & 15 \\ \hline
\(p(x)\) & 0.02 & 0.02 & 0.04 & 0.06 & 0.08 & 0.10 & 0.12 & 0.14 & 0.12 & 0.10 & 0.08 & 0.06 & 0.04 & 0.02 & 0 & 0 \\ \hline
\end{tabular}
\end{center}

Check: \(\sum p(x)=0.02+0.02+0.04+0.06+0.08+0.10+0.12+0.14+0.12+0.10+0.08+0.06+0.04+0.02=1.00\). Verified.

\textbf{Task 2: Expectation, variance, standard deviation, mode, median}

\[
E(X)=\sum x\,p(x)=\frac{1}{50}\sum x f(x)=\frac{343}{50}=6.86.
\]

\[
E(X^2)=\sum x^2 p(x)=\frac{1}{50}\sum x^2 f(x)=\frac{2793}{50}=55.86.
\]

\[
\operatorname{Var}(X)=E(X^2)-[E(X)]^2=55.86-(6.86)^2=55.86-47.0596=8.8004.
\]

\[
\sigma_X=\sqrt{8.8004}\approx 2.966.
\]

\textbf{Mode:} The value with highest probability is \(x=7\) (\(p=0.14\)). So mode \(=7\).

\textbf{Median:} Cumulative frequencies: 
\(F(0)=1\), \(F(1)=2\), \(F(2)=4\), \(F(3)=7\), \(F(4)=11\), \(F(5)=16\), \(F(6)=22\), \(F(7)=29\). 
The 25th and 26th ordered observations both fall at \(x=7\). Hence median \(=7\).

\textbf{Task 3: Cumulative distribution function and probabilities}

The CDF \(F(x)=P(X\le x)\) is a step function:

\[
F(x)=
\begin{cases}
0, & x<0\\
0.02, & 0\le x<1\\
0.04, & 1\le x<2\\
0.08, & 2\le x<3\\
0.14, & 3\le x<4\\
0.22, & 4\le x<5\\
0.32, & 5\le x<6\\
0.44, & 6\le x<7\\
0.58, & 7\le x<8\\
0.70, & 8\le x<9\\
0.80, & 9\le x<10\\
0.88, & 10\le x<11\\
0.94, & 11\le x<12\\
0.98, & 12\le x<13\\
1.00, & x\ge 13
\end{cases}
\]

\begin{popfigure}
\begin{tikzpicture}
\begin{axis}[
    width=\linewidth,
    height=8cm,
    axis lines=middle,
    xlabel=$x$ (Daily Demand),
    ylabel=$F(x)$,
    xmin=-0.5, xmax=10.5,
    ymin=-0.05, ymax=1.05,
    xtick={0,1,2,3,4,5,6,7,8,9,10},
    ytick={0,0.2,0.4,0.6,0.8,1},
    grid=major,
    grid style={popmuted, dashed},
    tick label style={font=\small},
    label style={font=\small},
    clip=false
]
% CDF of Poisson(3) approx - Step function horizontal lines
% F(0)=0.05, F(1)=0.20, F(2)=0.42, F(3)=0.65, F(4)=0.82, F(5)=0.92, F(6)=0.97, F(7)=0.99, F(8)=1.00
\draw[popblue, very thick] (-0.5,0) -- (0,0);
\draw[popblue, very thick] (0,0.05) -- (1,0.05);
\draw[popblue, very thick] (1,0.20) -- (2,0.20);
\draw[popblue, very thick] (2,0.42) -- (3,0.42);
\draw[popblue, very thick] (3,0.65) -- (4,0.65);
\draw[popblue, very thick] (4,0.82) -- (5,0.82);
\draw[popblue, very thick] (5,0.92) -- (6,0.92);
\draw[popblue, very thick] (6,0.97) -- (7,0.97);
\draw[popblue, very thick] (7,0.99) -- (8,0.99);
\draw[popblue, very thick] (8,1.00) -- (10.5,1.00);

% Vertical jumps at integer points
\draw[popblue, very thick] (0,0) -- (0,0.05);
\draw[popblue, very thick] (1,0.05) -- (1,0.20);
\draw[popblue, very thick] (2,0.20) -- (2,0.42);
\draw[popblue, very thick] (3,0.42) -- (3,0.65);
\draw[popblue, very thick] (4,0.65) -- (4,0.82);
\draw[popblue, very thick] (5,0.82) -- (5,0.92);
\draw[popblue, very thick] (6,0.92) -- (6,0.97);
\draw[popblue, very thick] (7,0.97) -- (7,0.99);
\draw[popblue, very thick] (8,0.99) -- (8,1.00);

% Filled circles (included, right-continuous)
\pgfplotsforeachungrouped \x/\y in {0/0.05, 1/0.20, 2/0.42, 3/0.65, 4/0.82, 5/0.92, 6/0.97, 7/0.99, 8/1.00} {
    \fill[popblue] (\x,\y) circle (2pt);
}
% Open circles (excluded, left limit)
\pgfplotsforeachungrouped \x/\y in {0/0, 1/0.05, 2/0.20, 3/0.42, 4/0.65, 5/0.82, 6/0.92, 7/0.97, 8/0.99} {
    \draw[popblue, fill=white, thick] (\x,\y) circle (2pt);
}

% Mean marker
\draw[poporange, dashed, thick] (3,0) -- (3,1) node[above, pos=0.9, poporange, font=\small] {$E(X)=3$};

\end{axis}
\end{tikzpicture}
\legende{Cumulative distribution function $F(x)=P(X\le x)$ for the daily demand (Poisson $\lambda=3$).}
\end{popfigure}

Using the CDF:
\[
P(X\le 10)=F(10)=0.88.
\]
\[
P(5<X\le 12)=F(12)-F(5)=0.98-0.32=0.66.
\]

\textbf{Task 4: Distribution of \(Y=X_1+X_2\)}

Since \(X_1\) and \(X_2\) are independent and identically distributed, the p.m.f. of \(Y\) is the convolution of the p.m.f. of \(X\) with itself:
\[
P(Y=y)=\sum_{x=0}^{15} p(x)\,p(y-x),\qquad y=0,1,\dots,26,
\]
where \(p(k)=0\) for \(k<0\) or \(k>15\). The full table has 27 entries; we show the first few and the central values.

\begin{center}
\begin{tabular}{|c|c|c|c|c|c|c|c|c|c|c|c|c|c|c|}
\hline
\rowcolor{popblueL}
\(y\) & 0 & 1 & 2 & 3 & 4 & 5 & 6 & 7 & 8 & 9 & 10 & 11 & 12 \\ \hline
\(P(Y=y)\) & 0.0004 & 0.0008 & 0.0020 & 0.0044 & 0.0084 & 0.0140 & 0.0212 & 0.0300 & 0.0396 & 0.0492 & 0.0576 & 0.0636 & 0.0660 \\ \hline
\end{tabular}
\end{center}
\begin{center}
\begin{tabular}{|c|c|c|c|c|c|c|c|c|c|c|c|c|c|c|}
\hline
\rowcolor{popblueL}
\(y\) & 13 & 14 & 15 & 16 & 17 & 18 & 19 & 20 & 21 & 22 & 23 & 24 & 25 & 26 \\ \hline
\(P(Y=y)\) & 0.0636 & 0.0576 & 0.0492 & 0.0396 & 0.0300 & 0.0212 & 0.0140 & 0.0084 & 0.0044 & 0.0020 & 0.0008 & 0.0004 & 0 & 0 \\ \hline
\end{tabular}
\end{center}
(Values for \(y=14\) to \(26\) are symmetric to those for \(y=12\) down to \(0\) because the distribution of \(X\) is symmetric about \(6.5\); the mean of \(Y\) is \(13.72\).)

\textbf{Verification of addition rules:}
\[
E(Y)=E(X_1)+E(X_2)=6.86+6.86=13.72.
\]
\[
\operatorname{Var}(Y)=\operatorname{Var}(X_1)+\operatorname{Var}(X_2)=8.8004+8.8004=17.6008.
\]
These match the values obtained directly from the convolution table (e.g. \(\sum y P(Y=y)=13.72\), \(\sum y^2 P(Y=y)=206.0392\), giving \(\operatorname{Var}(Y)=206.0392-13.72^2=17.6008\)).

\textbf{Task 5: Probability generating function}

The probability generating function of \(X\) is
\[
G_X(t)=E(t^X)=\sum_{x=0}^{15} p(x) t^x
=0.02 + 0.02t + 0.04t^2 + 0.06t^3 + 0.08t^4 + 0.10t^5 + 0.12t^6 + 0.14t^7 + 0.12t^8 + 0.10t^9 + 0.08t^{10} + 0.06t^{11} + 0.04t^{12} + 0.02t^{13}.
\]

Differentiate term by term:
\[
G'_X(t)=0.02 + 0.08t + 0.18t^2 + 0.32t^3 + 0.50t^4 + 0.72t^5 + 0.98t^6 + 1.28t^7 + 1.50t^8 + 1.62t^9 + 1.64t^{10} + 1.56t^{11} + 1.32t^{12} + 0.26t^{13}? \text{ No.}
\]
Correct differentiation:
\[
G'_X(t) = \sum_{x=1}^{13} x\,p(x)\,t^{x-1}.
\]
Evaluating at \(t=1\):
\[
G'_X(1) = \sum_{x=1}^{13} x\,p(x) = E(X).
\]
Numerically:
\[
G'_X(1) = 1(0.02)+2(0.04)+3(0.06)+4(0.08)+5(0.10)+6(0.12)+7(0.14)+8(0.12)+9(0.10)+10(0.08)+11(0.06)+12(0.04)+13(0.02) = 6.86.
\]
Thus \(G'_X(1)=6.86\), confirming the p.g.f. method.

\textbf{Task 6: Optimal stocking level}

Profit function: \(\Pi(S)=200\min(X,S)-100\max(S-X,0)\).

We compute \(E[\Pi(S)]\) for \(S=6,7,8,9\). For a given \(S\),
\[
E[\Pi(S)] = 200\,E[\min(X,S)] - 100\,E[\max(S-X,0)].
\]
We use the p.m.f. and CDF values from Tasks 1 and 3.

\textbf{For \(S=6\):}
\[
E[\min(X,6)] = \sum_{x=0}^5 x p(x) + 6 P(X\ge 6).
\]
\(\sum_{x=0}^5 x p(x) = 0(0.02)+1(0.02)+2(0.04)+3(0.06)+4(0.08)+5(0.10)=1.10.\)
\(P(X\ge 6)=1-F(5)=1-0.32=0.68.\)
\(E[\min]=1.10+6(0.68)=5.18.\)
\(E[\max(6-X,0)] = \sum_{x=0}^5 (6-x)p(x) = 6(0.02)+5(0.02)+4(0.04)+3(0.06)+2(0.08)+1(0.10)=0.82.\)
\(E[\Pi(6)] = 200(5.18) - 100(0.82) = 1036 - 82 = 954\) FCFA.

\textbf{For \(S=7\):}
\(\sum_{x=0}^6 x p(x) = 1.10 + 6(0.12)=1.82.\)
\(P(X\ge 7)=1-F(6)=1-0.44=0.56.\)
\(E[\min]=1.82+7(0.56)=5.74.\)
\(E[\max(7-X,0)] = \sum_{x=0}^6 (7-x)p(x) = 7(0.02)+6(0.02)+5(0.04)+4(0.06)+3(0.08)+2(0.10)+1(0.12)=1.26.\)
\(E[\Pi(7)] = 200(5.74) - 100(1.26) = 1148 - 126 = 1022\) FCFA.

\textbf{For \(S=8\):}
\(\sum_{x=0}^7 x p(x) = 1.82 + 7(0.14)=2.80.\)
\(P(X\ge 8)=1-F(7)=1-0.58=0.42.\)
\(E[\min]=2.80+8(0.42)=6.16.\)
\(E[\max(8-X,0)] = \sum_{x=0}^7 (8-x)p(x) = 8(0.02)+7(0.02)+6(0.04)+5(0.06)+4(0.08)+3(0.10)+2(0.12)+1(0.14)=1.84.\)
\(E[\Pi(8)] = 200(6.16) - 100(1.84) = 1232 - 184 = 1048\) FCFA.

\textbf{For \(S=9\):}
\(\sum_{x=0}^8 x p(x) = 2.80 + 8(0.12)=3.76.\)
\(P(X\ge 9)=1-F(8)=1-0.70=0.30.\)
\(E[\min]=3.76+9(0.30)=6.46.\)
\(E[\max(9-X,0)] = \sum_{x=0}^8 (9-x)p(x) = 9(0.02)+8(0.02)+7(0.04)+6(0.06)+5(0.08)+4(0.10)+3(0.12)+2(0.14)+1(0.12)=2.54.\)
\(E[\Pi(9)] = 200(6.46) - 100(2.54) = 1292 - 254 = 1038\) FCFA.

\begin{center}
\begin{tabular}{|c|c|c|c|}
\hline
\rowcolor{popblueL}
\(S\) & \(E[\min(X,S)]\) & \(E[\max(S-X,0)]\) & \(E[\Pi(S)]\) (FCFA) \\ \hline
6 & 5.18 & 0.82 & 954 \\
7 & 5.74 & 1.26 & 1022 \\
8 & 6.16 & 1.84 & 1048 \\
9 & 6.46 & 2.54 & 1038 \\ \hline
\end{tabular}
\end{center}

The maximum expected profit is \textbf{1048 FCFA} at \(S=8\). 

\textbf{Recommendation:} Madame Ngo Bakola should stock \textbf{8 cards} each day. This level balances the gain from meeting higher demand against the cost of unsold cards. The expected profit at \(S=8\) exceeds that at \(S=7\) by 26 FCFA and at \(S=9\) by 10 FCFA.

\textbf{Risk comment:} The variance of profit increases with \(S\) because unsold cards add variability. At \(S=8\), the probability of having unsold cards is \(P(X<8)=F(7)=0.58\), meaning more than half the days she will have some loss on unsold stock. However, the expected profit is still highest. If she is risk-averse, she might prefer \(S=7\) (expected profit 1022 FCFA, lower variance). The final choice depends on her risk tolerance, but mathematically \(S=8\) maximises expected profit.
\end{exoresolu}

\newpage

\coursec{Methods and worked examples}

\coursec{Random Variables and Discrete Random Variables}

\courssub{Random Variables and Discrete Random Variables}

\begin{definition}[Random Variable]
A \cle{random variable} (RV) is a function that assigns a real number to each outcome of a random experiment. It is usually denoted by a capital letter ($X$, $Y$, $Z$), while its possible values are denoted by lowercase letters ($x$, $y$, $z$).
\end{definition}

\begin{definition}[Discrete Random Variable]
A random variable $X$ is \cle{discrete} if it can take only a finite or countably infinite number of distinct values. The set of possible values is called the \cle{support} of $X$, denoted $S_X = \{x_1, x_2, \dots\}$.
\end{definition}

\begin{exemplebox}[Examples of Discrete Random Variables]
\begin{itemize}
\item Number of heads when tossing a coin $n$ times (finite support $\{0,1,\dots,n\}$).
\item Number of customers arriving at a shop in an hour (countably infinite support $\{0,1,2,\dots\}$).
\item Score on a fair die (finite support $\{1,2,3,4,5,6\}$).
\item Net gain in a game of chance (finite support).
\end{itemize}
\end{exemplebox}

\begin{attention}[Continuous vs Discrete]
A random variable that can take any value in an interval (e.g., height, weight, time) is \emph{continuous}. This chapter deals \emph{only} with discrete random variables. Continuous random variables are treated in Chapter PMS12.
\end{attention}

\courssub{Probability Distributions and the Probability Mass Function}

\begin{definition}[Probability Mass Function (PMF)]
The \cle{probability mass function} (PMF) of a discrete random variable $X$ is the function $p_X(x) = P(X=x)$ giving the probability that $X$ takes the value $x$. It satisfies:
\begin{enumerate}
\item $p_X(x) \ge 0$ for all $x \in \mathbb{R}$.
\item $\sum_{x \in S_X} p_X(x) = 1$.
\item $p_X(x) = 0$ for $x \notin S_X$.
\end{enumerate}
\end{definition}

\begin{propriete}[Probability Distribution Table]
The probability distribution of a discrete RV $X$ is often presented in a table:
\[
\begin{array}{c|cccc}
x & x_1 & x_2 & \cdots & x_n \\ \hline
P(X=x) & p_1 & p_2 & \cdots & p_n
\end{array}
\]
where $\sum_{i=1}^n p_i = 1$ and $p_i > 0$.
\end{propriete}

\begin{methode}[Constructing the Probability Distribution of a Discrete Random Variable]
\begin{enumerate}
\item Identify the random experiment and the sample space $\Omega$.
\item Define the discrete random variable $X$ by assigning a numerical value to each outcome.
\item List all possible values $x_i$ that $X$ can take (the support $S_X$).
\item Determine the probability $P(X=x_i)$ for each value, using the rules of probability (equally likely outcomes, tree diagrams, combinations, independence, etc.).
\item Check that $\sum_i P(X=x_i)=1$.
\item Present the distribution in a table.
\end{enumerate}
\end{methode}

\begin{exoresolu}[Probability Distribution of the Number of Heads in Three Tosses of a Biased Coin]
A biased coin has probability $p=\frac{2}{3}$ of showing Heads. The coin is tossed three times independently. Let $X$ be the number of Heads obtained.
\begin{enumerate}
\item Construct the probability distribution table of $X$.
\item Find $P(X\ge 2)$.
\end{enumerate}
\tcblower
\textbf{Solution.}
\begin{enumerate}
\item The experiment: toss a biased coin three times. $X$ can take values $0,1,2,3$. Since tosses are independent, $X$ follows a Binomial distribution with parameters $n=3$ and $p=\frac{2}{3}$.
\[
P(X=k)=\binom{3}{k}\left(\frac{2}{3}\right)^k\left(\frac{1}{3}\right)^{3-k},\quad k=0,1,2,3.
\]
Compute each probability:
\begin{align*}
P(X=0)&=\binom{3}{0}\left(\frac{2}{3}\right)^0\left(\frac{1}{3}\right)^3=\frac{1}{27},\\
P(X=1)&=\binom{3}{1}\left(\frac{2}{3}\right)^1\left(\frac{1}{3}\right)^2=3\times\frac{2}{3}\times\frac{1}{9}=\frac{6}{27}=\frac{2}{9},\\
P(X=2)&=\binom{3}{2}\left(\frac{2}{3}\right)^2\left(\frac{1}{3}\right)^1=3\times\frac{4}{9}\times\frac{1}{3}=\frac{12}{27}=\frac{4}{9},\\
P(X=3)&=\binom{3}{3}\left(\frac{2}{3}\right)^3\left(\frac{1}{3}\right)^0=\frac{8}{27}.
\end{align*}
Check sum: $\frac{1}{27}+\frac{6}{27}+\frac{12}{27}+\frac{8}{27}=\frac{27}{27}=1$.

Distribution table:
\[
\begin{array}{c|cccc}
x & 0 & 1 & 2 & 3 \\ \hline
P(X=x) & \frac{1}{27} & \frac{2}{9} & \frac{4}{9} & \frac{8}{27}
\end{array}
\]

\item $P(X\ge 2)=P(X=2)+P(X=3)=\frac{4}{9}+\frac{8}{27}=\frac{12}{27}+\frac{8}{27}=\frac{20}{27}$.
\end{enumerate}
\end{exoresolu}

\begin{exoresolu}[Finding an Unknown Parameter in a Distribution]
A discrete random variable $X$ has the following probability distribution:
\[
\begin{array}{c|ccccc}
x & 1 & 2 & 3 & 4 & 5 \\ \hline
P(X=x) & 0.1 & 0.2 & k & 0.3 & 0.1
\end{array}
\]
\begin{enumerate}
\item Find the value of $k$.
\item Write down the full probability distribution.
\end{enumerate}
\tcblower
\textbf{Solution.}
\begin{enumerate}
\item Sum of probabilities must equal 1:
\[
0.1 + 0.2 + k + 0.3 + 0.1 = 1 \implies 0.7 + k = 1 \implies k = 0.3.
\]
\item The full distribution is:
\[
\begin{array}{c|ccccc}
x & 1 & 2 & 3 & 4 & 5 \\ \hline
P(X=x) & 0.1 & 0.2 & 0.3 & 0.3 & 0.1
\end{array}
\]
\end{enumerate}
\end{exoresolu}

\courssub{Expectation, Variance and Standard Deviation}

\begin{definition}[Expectation (Mean)]
The \cle{expectation} or \cle{mean} of a discrete random variable $X$ is
\[
E(X) = \mu = \sum_{x \in S_X} x\,P(X=x),
\]
provided the series converges absolutely.
\end{definition}

\begin{definition}[Variance and Standard Deviation]
The \cle{variance} of $X$ is
\[
\text{Var}(X) = \sigma^2 = E\bigl[(X-\mu)^2\bigr] = \sum_{x \in S_X} (x-\mu)^2 P(X=x).
\]
The \cle{standard deviation} is $\sigma = \sqrt{\text{Var}(X)}$.
\end{definition}

\begin{propriete}[Computational Formula for Variance (Examinable)]
\[
\text{Var}(X) = E(X^2) - [E(X)]^2, \quad \text{where } E(X^2) = \sum x^2 P(X=x).
\]
\end{propriete}

\begin{propriete}[Linear Transformation Properties (Examinable)]
For constants $a,b \in \mathbb{R}$:
\begin{align*}
E(aX+b) &= aE(X) + b, \\
\text{Var}(aX+b) &= a^2\text{Var}(X).
\end{align*}
In particular, $\text{Var}(X+b) = \text{Var}(X)$ (adding a constant does not change spread).
\end{propriete}

\begin{methode}[Computing Expectation, Variance and Standard Deviation]
\begin{enumerate}
\item Write down the probability distribution of $X$.
\item Compute $E(X) = \sum x\,P(X=x)$.
\item Compute $E(X^2) = \sum x^2\,P(X=x)$.
\item Compute $\text{Var}(X) = E(X^2) - [E(X)]^2$.
\item Compute $\sigma_X = \sqrt{\text{Var}(X)}$.
\item For linear transformations, use $E(aX+b)=aE(X)+b$ and $\text{Var}(aX+b)=a^2\text{Var}(X)$.
\end{enumerate}
\end{methode}

\begin{exoresolu}[Mean and Variance of a Discrete Random Variable]
A discrete random variable $X$ has the following probability distribution:
\[
\begin{array}{c|ccccc}
x & 1 & 2 & 3 & 4 & 5 \\ \hline
P(X=x) & 0.1 & 0.2 & 0.3 & 0.3 & 0.1
\end{array}
\]
\begin{enumerate}
\item Calculate $E(X)$ and $\text{Var}(X)$.
\item Find $E(2X-3)$ and $\text{Var}(2X-3)$.
\end{enumerate}
\tcblower
\textbf{Solution.}
\begin{enumerate}
\item 
\[
E(X)=\sum xP(X=x)=1(0.1)+2(0.2)+3(0.3)+4(0.3)+5(0.1)=0.1+0.4+0.9+1.2+0.5=3.1.
\]
\[
E(X^2)=\sum x^2P(X=x)=1^2(0.1)+2^2(0.2)+3^2(0.3)+4^2(0.3)+5^2(0.1)=0.1+0.8+2.7+4.8+2.5=10.9.
\]
\[
\text{Var}(X)=E(X^2)-[E(X)]^2=10.9-(3.1)^2=10.9-9.61=1.29.
\]
Standard deviation $\sigma_X = \sqrt{1.29} \approx 1.136$.

\item Using linearity properties:
\[
E(2X-3)=2E(X)-3=2(3.1)-3=6.2-3=3.2.
\]
\[
\text{Var}(2X-3)=2^2\text{Var}(X)=4\times1.29=5.16.
\]
\end{enumerate}
\end{exoresolu}

\begin{exoresolu}[Finding a Parameter and Expectation]
The random variable $X$ takes values $1,2,3,4$ with probabilities:
\[
P(X=x)=\begin{cases}
kx & x=1,2\\
k(5-x) & x=3,4\\
0 & \text{otherwise}
\end{cases}
\]
where $k$ is a constant.
\begin{enumerate}
\item Find the value of $k$.
\item Calculate $E(X)$ and $\text{Var}(X)$.
\item Find $P(X>E(X))$.
\end{enumerate}
\tcblower
\textbf{Solution.}
\begin{enumerate}
\item Sum of probabilities:
\[
P(X=1)=k,\quad P(X=2)=2k,\quad P(X=3)=2k,\quad P(X=4)=k.
\]
Sum $=k+2k+2k+k=6k=1 \Rightarrow k=\frac16$.
Distribution:
\[
\begin{array}{c|cccc}
x & 1 & 2 & 3 & 4 \\ \hline
P(X=x) & \frac16 & \frac13 & \frac13 & \frac16
\end{array}
\]

\item $E(X)=1\cdot\frac16+2\cdot\frac13+3\cdot\frac13+4\cdot\frac16=\frac16+\frac23+1+\frac23=\frac16+\frac46+\frac66+\frac46=\frac{15}{6}=2.5$.
$E(X^2)=1^2\cdot\frac16+4\cdot\frac13+9\cdot\frac13+16\cdot\frac16=\frac16+\frac43+3+\frac83=\frac16+\frac86+\frac{18}{6}+\frac{16}{6}=\frac{43}{6}$.
$\text{Var}(X)=E(X^2)-[E(X)]^2=\frac{43}{6}-\left(\frac52\right)^2=\frac{43}{6}-\frac{25}{4}=\frac{86}{12}-\frac{75}{12}=\frac{11}{12}$.

\item $E(X)=2.5$. $P(X>2.5)=P(X=3)+P(X=4)=\frac13+\frac16=\frac12$.
\end{enumerate}
\end{exoresolu}

\courssub{The Cumulative Distribution Function, Mode and Median}

\begin{definition}[Cumulative Distribution Function (CDF)]
The \cle{cumulative distribution function} (CDF) of a discrete random variable $X$ is defined for all real $x$ by
\[
F(x) = P(X \le x) = \sum_{x_i \le x} P(X=x_i).
\]
\end{definition}

\begin{propriete}[Properties of the CDF (Examinable)]
For a discrete RV $X$:
\begin{enumerate}
\item $F(x)$ is a step function, non-decreasing.
\item $F(x)$ is right-continuous: $\lim_{h\to 0^+} F(x+h) = F(x)$.
\item $\lim_{x\to -\infty} F(x) = 0$, $\lim_{x\to +\infty} F(x) = 1$.
\item $P(X=x) = F(x) - F(x^-)$ where $F(x^-) = \lim_{h\to 0^+} F(x-h)$ (the jump at $x$).
\item For any $a < b$, $P(a < X \le b) = F(b) - F(a)$.
\end{enumerate}
\end{propriete}

\begin{definition}[Mode]
The \cle{mode} of a discrete random variable $X$ is the value $x$ that maximizes $P(X=x)$, i.e., the most probable value. If several values share the maximum probability, the distribution is \cle{multimodal}.
\end{definition}

\begin{definition}[Median]
A \cle{median} of a discrete random variable $X$ is any value $m$ such that
\[
P(X \le m) \ge \frac12 \quad \text{and} \quad P(X \ge m) \ge \frac12.
\]
For discrete variables, the median is often taken as the smallest $x$ with $F(x) \ge 0.5$.
\end{definition}

\begin{methode}[Cumulative Distribution Function, Median and Mode]
\begin{enumerate}
\item For a discrete random variable $X$, compute $F(x)=P(X\le x)$ for all real $x$ by summing probabilities up to $x$.
\item $F(x)$ is constant between possible values of $X$ and jumps at each $x_i \in S_X$.
\item \textbf{Mode:} Identify the value(s) with the highest probability $P(X=x)$.
\item \textbf{Median:} Find the smallest $x$ such that $F(x) \ge 0.5$ (or any $m$ satisfying the definition).
\end{enumerate}
\end{methode}

\begin{exoresolu}[CDF, Median and Mode]
The random variable $X$ has probability mass function:
\[
P(X=x)=\begin{cases}
\frac{x}{10}, & x=1,2,3,4\\
0, & \text{otherwise}.
\end{cases}
\]
\begin{enumerate}
\item Verify that this is a valid probability distribution.
\item Construct the cumulative distribution function $F(x)$.
\item Find the mode and the median of $X$.
\end{enumerate}
\tcblower
\textbf{Solution.}
\begin{enumerate}
\item $\sum_{x=1}^4 \frac{x}{10}=\frac{1+2+3+4}{10}=\frac{10}{10}=1$. All probabilities are non-negative. Valid.
\item Compute $F(x)=P(X\le x)$:
\begin{align*}
F(1)&=P(X\le1)=0.1,\\
F(2)&=P(X\le2)=0.1+0.2=0.3,\\
F(3)&=P(X\le3)=0.3+0.3=0.6,\\
F(4)&=P(X\le4)=0.6+0.4=1.
\end{align*}
For $x<1$, $F(x)=0$; for $1\le x<2$, $F(x)=0.1$; for $2\le x<3$, $F(x)=0.3$; for $3\le x<4$, $F(x)=0.6$; for $x\ge4$, $F(x)=1$.
In table form:
\[
\begin{array}{c|ccccc}
x & (-\infty,1) & [1,2) & [2,3) & [3,4) & [4,\infty) \\ \hline
F(x) & 0 & 0.1 & 0.3 & 0.6 & 1
\end{array}
\]

\item \textbf{Mode:} The highest probability is $P(X=4)=0.4$. Mode $=4$.
\textbf{Median:} Smallest $x$ with $F(x)\ge0.5$. $F(2)=0.3<0.5$, $F(3)=0.6\ge0.5$. Median $=3$.
\end{enumerate}
\end{exoresolu}

\courssub{The Distribution of $X_1 + X_2$}

\begin{propriete}[Sum of Independent Discrete Random Variables]
Let $X$ and $Y$ be independent discrete random variables. The sum $S=X+Y$ is a discrete random variable with probability mass function given by the \cle{convolution}:
\[
P(S=s) = \sum_{x} P(X=x)\,P(Y=s-x),
\]
where the sum runs over all $x$ such that $P(X=x)>0$ and $P(Y=s-x)>0$.
\end{propriete}

\begin{propriete}[Expectation and Variance of a Sum (Examinable)]
For any random variables $X,Y$ (not necessarily independent):
\[
E(X+Y) = E(X) + E(Y).
\]
If $X$ and $Y$ are independent:
\[
\text{Var}(X+Y) = \text{Var}(X) + \text{Var}(Y).
\]
\end{propriete}

\begin{methode}[Distribution of the Sum of Two Independent Discrete Random Variables]
\begin{enumerate}
\item Let $X$ and $Y$ be independent discrete random variables with known distributions.
\item List all pairs $(x,y)$ with their joint probabilities $P(X=x)P(Y=y)$ (by independence).
\item Compute $s = x+y$ for each pair.
\item Group pairs by the value of $s$ and add the probabilities to obtain $P(S=s)$.
\item Alternatively, use the convolution formula directly.
\item Compute $E(S)$ and $\text{Var}(S)$ either from the distribution of $S$ or using the additive properties.
\end{enumerate}
\end{methode}

\begin{exoresolu}[Sum of Two Independent Dice Scores]
Two fair six-sided dice are rolled. Let $X$ be the score on the first die, $Y$ the score on the second. $X$ and $Y$ are independent. Let $S=X+Y$.
\begin{enumerate}
\item Find the probability distribution of $S$.
\item Compute $E(S)$ and $\text{Var}(S)$ directly from the distribution of $S$, and verify the additive properties.
\end{enumerate}
\tcblower
\textbf{Solution.}
\begin{enumerate}
\item Possible values of $S$: $2,3,\dots,12$. There are $36$ equally likely outcomes $(x,y)$. Count pairs for each sum:
\[
\begin{array}{c|ccccccccccc}
s & 2 & 3 & 4 & 5 & 6 & 7 & 8 & 9 & 10 & 11 & 12 \\ \hline
\text{Number of pairs} & 1 & 2 & 3 & 4 & 5 & 6 & 5 & 4 & 3 & 2 & 1
\end{array}
\]
Thus $P(S=s)=\frac{\text{number of pairs}}{36}$.
Distribution:
\[
\begin{array}{c|ccccccccccc}
s & 2 & 3 & 4 & 5 & 6 & 7 & 8 & 9 & 10 & 11 & 12 \\ \hline
P(S=s) & \frac1{36} & \frac2{36} & \frac3{36} & \frac4{36} & \frac5{36} & \frac6{36} & \frac5{36} & \frac4{36} & \frac3{36} & \frac2{36} & \frac1{36}
\end{array}
\]

\item Compute $E(S)$:
\[
E(S)=\sum s\,P(S=s)=\frac{1}{36}\bigl(2\cdot1+3\cdot2+4\cdot3+5\cdot4+6\cdot5+7\cdot6+8\cdot5+9\cdot4+10\cdot3+11\cdot2+12\cdot1\bigr).
\]
Sum $=2+6+12+20+30+42+40+36+30+22+12=252$. $E(S)=\frac{252}{36}=7$.
$E(X)=E(Y)=\frac{1+2+3+4+5+6}{6}=3.5$. $E(X)+E(Y)=7$. Verified.

$E(S^2)=\frac{1}{36}\bigl(4\cdot1+9\cdot2+16\cdot3+25\cdot4+36\cdot5+49\cdot6+64\cdot5+81\cdot4+100\cdot3+121\cdot2+144\cdot1\bigr)$.
Sum $=4+18+48+100+180+294+320+324+300+242+144=1974$. $E(S^2)=\frac{1974}{36}=\frac{329}{6}$.
$\text{Var}(S)=E(S^2)-[E(S)]^2=\frac{329}{6}-49=\frac{329-294}{6}=\frac{35}{6}\approx5.8333$.

$\text{Var}(X)=\text{Var}(Y)=\frac{35}{12}$ (standard result for a fair die). $\text{Var}(X)+\text{Var}(Y)=\frac{35}{12}+\frac{35}{12}=\frac{35}{6}$. Verified.
\end{enumerate}
\end{exoresolu}

\begin{exoresolu}[Sum of Two Independent Non-Identical Distributions]
Let $X$ and $Y$ be independent random variables with distributions:
\[
\begin{array}{c|ccc}
x & 0 & 1 & 2 \\ \hline
P(X=x) & 0.2 & 0.5 & 0.3
\end{array}
\quad
\begin{array}{c|cc}
y & 1 & 2 \\ \hline
P(Y=y) & 0.4 & 0.6
\end{array}
\]
Find the distribution of $S=X+Y$, $E(S)$ and $\text{Var}(S)$.
\tcblower
\textbf{Solution.}
Possible values of $S$: $1,2,3,4$.
\begin{align*}
P(S=1) &= P(X=0,Y=1) = 0.2 \times 0.4 = 0.08,\\
P(S=2) &= P(X=0,Y=2) + P(X=1,Y=1) = 0.2\times0.6 + 0.5\times0.4 = 0.12 + 0.20 = 0.32,\\
P(S=3) &= P(X=1,Y=2) + P(X=2,Y=1) = 0.5\times0.6 + 0.3\times0.4 = 0.30 + 0.12 = 0.42,\\
P(S=4) &= P(X=2,Y=2) = 0.3 \times 0.6 = 0.18.
\end{align*}
Check sum: $0.08+0.32+0.42+0.18=1.00$.
Distribution of $S$:
\[
\begin{array}{c|cccc}
s & 1 & 2 & 3 & 4 \\ \hline
P(S=s) & 0.08 & 0.32 & 0.42 & 0.18
\end{array}
\]
$E(X)=0(0.2)+1(0.5)+2(0.3)=1.1$, $E(Y)=1(0.4)+2(0.6)=1.6$. $E(S)=1.1+1.6=2.7$.
Check: $1(0.08)+2(0.32)+3(0.42)+4(0.18)=0.08+0.64+1.26+0.72=2.7$.
$\text{Var}(X)=E(X^2)-[E(X)]^2=(0+0.5+1.2)-1.21=1.7-1.21=0.49$.
$\text{Var}(Y)=E(Y^2)-[E(Y)]^2=(0.4+2.4)-2.56=2.8-2.56=0.24$.
$\text{Var}(S)=0.49+0.24=0.73$.
Check from $S$: $E(S^2)=1(0.08)+4(0.32)+9(0.42)+16(0.18)=0.08+1.28+3.78+2.88=8.02$. $\text{Var}(S)=8.02-2.7^2=8.02-7.29=0.73$.
\end{exoresolu}

\courssub{Probability Generating Functions}

\begin{definition}[Probability Generating Function (PGF)]
For a discrete random variable $X$ taking values in $\{0,1,2,\dots\}$ (non-negative integers), the \cle{probability generating function} (PGF) is defined by
\[
G_X(t) = E(t^X) = \sum_{k=0}^\infty P(X=k)\,t^k,
\]
for all real $t$ where the series converges (at least for $|t|\le 1$).
\end{definition}

\begin{attention}[Domain of PGF]
The PGF is defined for RVs taking values $0,1,2,\dots$. If $X$ takes other integer values (e.g., negative), we can sometimes shift it ($Y=X+c$) to use PGF, or use the Moment Generating Function (MGF) instead. In GCE Advanced Level, PGF questions typically involve non-negative integer-valued RVs.
\end{attention}

\begin{propriete}[Properties of the PGF (Examinable)]
Let $X$ be a discrete RV with PGF $G_X(t)$.
\begin{enumerate}
\item $G_X(1) = \sum P(X=k) = 1$.
\item $G_X(0) = P(X=0)$.
\item \textbf{Mean:} $E(X) = G_X'(1)$ (first derivative evaluated at $t=1$).
\item \textbf{Variance:} $\text{Var}(X) = G_X''(1) + G_X'(1) - [G_X'(1)]^2$.
\item \textbf{Sum of independent variables:} If $X,Y$ independent, $G_{X+Y}(t) = G_X(t)G_Y(t)$.
\item \textbf{Probabilities from PGF:} $P(X=k) = \frac{G_X^{(k)}(0)}{k!}$ (coefficient of $t^k$ in the power series expansion).
\end{enumerate}
\end{propriete}

\begin{aretenir}[Common PGFs to Memorise]
\begin{itemize}
\item \textbf{Bernoulli($p$)} ($P(X=1)=p, P(X=0)=q=1-p$): $G(t)=q+pt$.
\item \textbf{Binomial($n,p$)}: $G(t)=(q+pt)^n$.
\item \textbf{Poisson($\lambda$)}: $G(t)=e^{\lambda(t-1)}$.
\item \textbf{Geometric($p$)} on $\{0,1,2,\dots\}$ ($P(X=k)=pq^k$): $G(t)=\frac{p}{1-qt}$.
\end{itemize}
\end{aretenir}

\begin{methode}[Probability Generating Function — Definition and Basic Properties]
\begin{enumerate}
\item Write down the PMF $P(X=k)$ for $k=0,1,2,\dots$.
\item Form the series $G_X(t) = \sum_{k=0}^\infty P(X=k)t^k$. Simplify to a closed form if possible.
\item Check $G_X(1)=1$.
\item To find mean: compute $G_X'(t)$, evaluate at $t=1$.
\item To find variance: compute $G_X''(t)$, evaluate at $t=1$, use $\text{Var}(X)=G_X''(1)+G_X'(1)-[G_X'(1)]^2$.
\item To find probabilities: expand $G_X(t)$ as a power series; coefficient of $t^k$ is $P(X=k)$.
\end{enumerate}
\end{methode}

\begin{exoresolu}[Using the PGF to Find Mean and Variance]
A discrete random variable $X$ has probability generating function $G_X(t)=\frac{1}{12}(1+t)^2(1+2t)$ for $|t|\le1$.
\begin{enumerate}
\item Find the probability distribution of $X$.
\item Compute $E(X)$ and $\text{Var}(X)$ using the PGF.
\end{enumerate}
\tcblower
\textbf{Solution.}
\begin{enumerate}
\item Expand $G_X(t)$:
\[
(1+t)^2=1+2t+t^2.
\]
\[
(1+2t+t^2)(1+2t)=1(1+2t)+2t(1+2t)+t^2(1+2t)=1+2t+2t+4t^2+t^2+2t^3=1+4t+5t^2+2t^3.
\]
Multiply by $\frac{1}{12}$:
\[
G_X(t)=\frac{1}{12}+\frac{4}{12}t+\frac{5}{12}t^2+\frac{2}{12}t^3=\frac{1}{12}+\frac{1}{3}t+\frac{5}{12}t^2+\frac{1}{6}t^3.
\]
Since $G_X(t)=\sum P(X=k)t^k$, we read off:
\[
P(X=0)=\frac{1}{12},\quad P(X=1)=\frac{1}{3},\quad P(X=2)=\frac{5}{12},\quad P(X=3)=\frac{1}{6}.
\]
Check sum: $\frac{1}{12}+\frac{4}{12}+\frac{5}{12}+\frac{2}{12}=\frac{12}{12}=1$. Valid.

\item Compute derivatives:
\[
G_X'(t)=\frac{1}{3}+\frac{10}{12}t+\frac{3}{6}t^2=\frac{1}{3}+\frac{5}{6}t+\frac{1}{2}t^2.
\]
\[
G_X''(t)=\frac{5}{6}+t.
\]
Evaluate at $t=1$:
\[
G_X'(1)=\frac{1}{3}+\frac{5}{6}+\frac{1}{2}=\frac{2}{6}+\frac{5}{6}+\frac{3}{6}=\frac{10}{6}=\frac{5}{3}.
\]
\[
G_X''(1)=\frac{5}{6}+1=\frac{11}{6}.
\]
Then $E(X)=G_X'(1)=\frac{5}{3}$.
$\text{Var}(X)=G_X''(1)+G_X'(1)-[G_X'(1)]^2=\frac{11}{6}+\frac{5}{3}-\left(\frac{5}{3}\right)^2=\frac{11}{6}+\frac{10}{6}-\frac{25}{9}=\frac{21}{6}-\frac{25}{9}=\frac{7}{2}-\frac{25}{9}=\frac{63}{18}-\frac{50}{18}=\frac{13}{18}$.
\end{enumerate}
\end{exoresolu}

\begin{methode}[Probability Generating Function of the Sum of Independent Random Variables]
\begin{enumerate}
\item If $X_1,X_2,\dots,X_n$ are independent, then $G_{X_1+\dots+X_n}(t)=G_{X_1}(t)G_{X_2}(t)\cdots G_{X_n}(t)$.
\item To find the distribution of the sum:
\begin{itemize}
\item Expand the product PGF and read coefficients (if simple).
\item Use the convolution method (Method 5).
\item Recognise the resulting PGF as that of a known distribution (e.g., sum of independent Poissons is Poisson with parameter sum of parameters; sum of independent Binomials with same $p$ is Binomial).
\end{itemize}
\item Mean and variance of the sum are sums of individual means and variances.
\end{enumerate}
\end{methode}

\begin{exoresolu}[Sum of Independent Poisson Variables via PGF]
The number of customers arriving at a shop in Yaoundé during the morning follows a Poisson distribution with mean $3$, and the number arriving in the afternoon follows an independent Poisson distribution with mean $5$. Let $X$ be the morning arrivals, $Y$ the afternoon arrivals.
\begin{enumerate}
\item Write down the PGFs of $X$ and $Y$.
\item Find the PGF of $T=X+Y$ (total daily arrivals).
\item Identify the distribution of $T$ and state its mean and variance.
\item Find $P(T=7)$.
\end{enumerate}
\tcblower
\textbf{Solution.}
\begin{enumerate}
\item For Poisson($\lambda$), $G(t)=e^{\lambda(t-1)}$.
$G_X(t)=e^{3(t-1)}$, $G_Y(t)=e^{5(t-1)}$.
\item $G_T(t)=G_X(t)G_Y(t)=e^{3(t-1)}e^{5(t-1)}=e^{(3+5)(t-1)}=e^{8(t-1)}$.
\item $G_T(t)$ is the PGF of a Poisson distribution with parameter $\lambda=8$. Thus $T\sim\text{Poisson}(8)$.
Mean $E(T)=8$, Variance $\text{Var}(T)=8$.
\item $P(T=7)=\frac{e^{-8}8^7}{7!}$. Compute numerically: $8^7=2097152$, $7!=5040$, $\frac{2097152}{5040}\approx416.1$, $e^{-8}\approx0.00033546$, product $\approx0.1396$.
\end{enumerate}
\end{exoresolu}

\begin{exoresolu}[Sum of Two Independent Dice Scores via PGF]
Two fair six-sided dice are rolled. Let $X$ and $Y$ be the scores. Find the PGF of $X$, the PGF of $S=X+Y$, and hence the probability $P(S=7)$.
\tcblower
\textbf{Solution.}
$X$ takes values $1,\dots,6$ with probability $\frac16$ each. But PGF requires values $0,1,2,\dots$. We can define $X'=X-1$ (values $0,\dots,5$) or use the PGF definition $G_X(t)=E(t^X)=\frac{1}{6}(t+t^2+t^3+t^4+t^5+t^6)$ which is valid for any integer-valued RV as a formal power series. The properties $G_{X+Y}=G_XG_Y$ and $E(X)=G_X'(1)$ still hold for integer-valued RVs (with $G_X(t)=\sum P(X=k)t^k$).
\[
G_X(t)=\frac{t(1-t^6)}{6(1-t)} = \frac{1}{6}(t+t^2+t^3+t^4+t^5+t^6).
\]
Since $X,Y$ independent and identically distributed:
\[
G_S(t)=[G_X(t)]^2 = \frac{1}{36}(t+t^2+t^3+t^4+t^5+t^6)^2.
\]
We need coefficient of $t^7$ in the expansion. The pairs $(i,j)$ with $i+j=7$, $1\le i,j\le 6$ are $(1,6),(2,5),(3,4),(4,3),(5,2),(6,1)$ — 6 pairs.
Each pair contributes $\frac{1}{36}t^7$. Coefficient $=6\times\frac{1}{36}=\frac{1}{6}$.
Thus $P(S=7)=\frac{1}{6}$.
\end{exoresolu}

\courssub{Comprehensive Applications}

\begin{methode}[Comprehensive Problem: Game of Chance with Discrete Random Variable]
\begin{enumerate}
\item Model the game: define the random variable (net gain/loss or winnings).
\item Construct its probability distribution (account for all outcomes).
\item Compute expectation (fair game if $E=0$).
\item Compute variance/standard deviation for risk assessment.
\item If sums of independent plays are considered, use additive properties of expectation/variance or PGF (if the RV is non-negative integer valued).
\end{enumerate}
\end{methode}

\begin{exoresolu}[A Game at a Cameroon Fair]
At a fair in Douala, a game consists of drawing a ball from a bag containing 4 red, 3 blue, and 3 green balls. You pay 500 FCFA to play. If you draw a red ball, you win 2000 FCFA; if blue, you win 1000 FCFA; if green, you win nothing. Let $W$ be your net gain (winnings minus cost).
\begin{enumerate}
\item Find the probability distribution of $W$.
\item Calculate $E(W)$ and $\text{Var}(W)$. Is the game fair?
\item If you play twice independently, let $T=W_1+W_2$ be your total net gain. Find $E(T)$ and $\text{Var}(T)$.
\item Find the probability that your total net gain after two games is exactly 1000 FCFA.
\end{enumerate}
\tcblower
\textbf{Solution.}
\begin{enumerate}
\item Probabilities: $P(\text{Red})=\frac{4}{10}=0.4$, $P(\text{Blue})=\frac{3}{10}=0.3$, $P(\text{Green})=\frac{3}{10}=0.3$.
Net gains:
Red: $2000-500=1500$ FCFA.
Blue: $1000-500=500$ FCFA.
Green: $0-500=-500$ FCFA.
Distribution of $W$:
\[
\begin{array}{c|ccc}
w & 1500 & 500 & -500 \\ \hline
P(W=w) & 0.4 & 0.3 & 0.3
\end{array}
\]

\item $E(W)=1500(0.4)+500(0.3)+(-500)(0.3)=600+150-150=600$ FCFA.
$E(W^2)=1500^2(0.4)+500^2(0.3)+(-500)^2(0.3)=2250000(0.4)+250000(0.3)+250000(0.3)=900000+75000+75000=1050000$.
$\text{Var}(W)=1050000-600^2=1050000-360000=690000$ FCFA$^2$.
$\sigma_W=\sqrt{690000}\approx830.66$ FCFA.
Since $E(W)=600>0$, the game is \textbf{not fair}; it favours the player (on average you gain 600 FCFA per game).

\item For two independent games, $T=W_1+W_2$.
$E(T)=E(W_1)+E(W_2)=600+600=1200$ FCFA.
$\text{Var}(T)=\text{Var}(W_1)+\text{Var}(W_2)=690000+690000=1380000$ FCFA$^2$.

\item We need $P(T=1000)$. Possible pairs of outcomes:
\begin{itemize}
\item (Red, Green) or (Green, Red): net gain $1500+(-500)=1000$. Probability $2\times0.4\times0.3=0.24$.
\item (Blue, Blue): net gain $500+500=1000$. Probability $0.3\times0.3=0.09$.
\end{itemize}
No other combinations give 1000.
Total $P(T=1000)=0.24+0.09=0.33$.
\end{enumerate}
\end{exoresolu}

\begin{savaistu}[PGF and the Douala Fair Game]
The net gain $W$ takes negative values, so its standard PGF $G_W(t)=E(t^W)$ would involve negative powers of $t$ (a Laurent series). While the convolution property $G_{W_1+W_2}=G_{W_1}G_{W_2}$ still holds formally, the standard PGF theory (derivatives at $t=1$ giving moments) applies directly only to non-negative integer-valued RVs. For $W$, we used the direct convolution method (listing pairs) which is always valid. In exams, if a PGF question involves a RV with negative values, shift it to make it non-negative (e.g., define $U=W+500$ taking values $0,1000,2000$) or use the MGF $M_X(t)=E(e^{tX})$.
\end{savaistu}

\courssub{Exercises}

\begin{exercice}[Exercise 1]
A discrete random variable $X$ has probability function $P(X=x)=\frac{c}{x(x+1)}$ for $x=1,2,3,4,5$.
\begin{enumerate}
\item Find the value of the constant $c$.
\item Calculate $E(X)$ and $\text{Var}(X)$.
\item Find the cumulative distribution function $F(x)$ and sketch its graph.
\item Determine the mode and median of $X$.
\end{enumerate}
\end{exercice}

\begin{exercice}[Exercise 2]
Two independent random variables $X$ and $Y$ have the following distributions:
\[
\begin{array}{c|ccc}
x & 0 & 1 & 2 \\ \hline
P(X=x) & 0.3 & 0.4 & 0.3
\end{array}
\quad
\begin{array}{c|cc}
y & 0 & 1 \\ \hline
P(Y=y) & 0.6 & 0.4
\end{array}
\]
\begin{enumerate}
\item Find the probability distribution of $S=X+Y$.
\item Compute $E(S)$ and $\text{Var}(S)$ using the additive properties.
\item Verify your answers by computing directly from the distribution of $S$.
\end{enumerate}
\end{exercice}

\begin{exercice}[Exercise 3]
A random variable $X$ has PGF $G_X(t)=\frac{1}{8}(1+t)^3$.
\begin{enumerate}
\item Identify the distribution of $X$ and its parameters.
\item Find $E(X)$ and $\text{Var}(X)$ using the PGF.
\item Find $P(X\ge 2)$.
\end{enumerate}
\end{exercice}

\begin{exercice}[Exercise 4]
The number of emails received per hour by a server follows a Poisson distribution with mean $4$. The number of spam emails per hour follows an independent Poisson distribution with mean $1$. Let $T$ be the total number of emails received in a 2-hour period.
\begin{enumerate}
\item Find the distribution of $T$.
\item Calculate $P(T=10)$.
\item Find the probability that exactly 3 of the 10 emails are spam, given $T=10$. (Hint: conditional distribution of Poisson components).
\end{enumerate}
\end{exercice}

\begin{exercice}[Exercise 5]
A biased tetrahedral die (faces 1,2,3,4) has probabilities $P(1)=0.1$, $P(2)=0.2$, $P(3)=0.3$, $P(4)=0.4$. The die is rolled twice. Let $M$ be the maximum of the two scores.
\begin{enumerate}
\item Find the probability distribution of $M$.
\item Calculate $E(M)$ and $\text{Var}(M)$.
\item Find the probability generating function of $M$ (treating values as $1,2,3,4$).
\end{enumerate}
\end{exercice}

% Corrections can be provided in a separate \corrige environment if needed.
% \begin{corrige}[Exercise 1] ... \end{corrige} etc.

\newpage

\coursec{Exercises}

\courssub{I check my knowledge}

\begin{exercice}[MCQ: spotting a valid distribution]
A discrete random variable $X$ takes the values $0, 1, 2, 3$. In each case, state whether the table can define a probability distribution of $X$. Justify your answer.
\begin{enumerate}
\item $P(X=x) = \dfrac{x+1}{10}$ for $x = 0, 1, 2, 3$.
\item $P(X=0) = 0.2$, $P(X=1) = 0.3$, $P(X=2) = 0.6$, $P(X=3) = -0.1$.
\item $P(X=x) = \dfrac{1}{2^{x+1}}$ for $x = 0, 1, 2, 3$.
\item $P(X=x) = 0.25$ for $x = 0, 1, 2, 3$.
\end{enumerate}
\end{exercice}

\begin{exercice}[True or false?]
Answer \emph{true} or \emph{false} to each statement, justifying your answer with a definition, a property or a counter-example.
\begin{enumerate}
\item If $X$ is a discrete random variable, then $\sum_x P(X = x) = 1$.
\item For any random variable $X$, $\operatorname{Var}(X) = E(X^2) - E(X)^2$.
\item The median of a discrete random variable is always one of the values taken by $X$.
\item If $G_X(t)$ is the probability generating function of $X$, then $G_X(1) = E(X)$.
\item If $X_1$ and $X_2$ are independent, then $E(X_1 + X_2) = E(X_1) + E(X_2)$.
\end{enumerate}
\end{exercice}

\begin{exercice}[Definitions in your own words]
\begin{enumerate}
\item Define a \cle{random variable} and explain the difference between a discrete and a continuous random variable, giving one example of each.
\item Define the \cle{cumulative distribution function} $F(x)$ of a discrete random variable $X$ and state two of its properties.
\item Define the \cle{probability generating function} of a discrete random variable $X$ taking values in $\mathbb{N}$, and state how $E(X)$ is obtained from it.
\end{enumerate}
\end{exercice}

\begin{exercice}[Direct calculations]
A discrete random variable $X$ has distribution
\begin{center}
\begin{tabular}{|c|c|c|c|c|}
\hline
\rowcolor{popblueL} $x$ & $1$ & $2$ & $3$ & $4$ \\
\hline
$P(X = x)$ & $0.1$ & $0.3$ & $0.4$ & $0.2$ \\
\hline
\end{tabular}
\end{center}
\begin{enumerate}
\item Verify that this is a valid probability distribution.
\item Calculate $E(X)$, $E(X^2)$ and $\operatorname{Var}(X)$.
\item Write down $F(2)$ and $F(3.5)$.
\item Determine the mode of $X$.
\end{enumerate}
\end{exercice}

\courssub{I practise}

\begin{exercice}[Finding the constant $k$]
A discrete random variable $X$ takes the values $1, 2, 3, 4$ with
\[ P(X = x) = kx, \qquad x = 1, 2, 3, 4. \]
\begin{enumerate}
\item Find the value of the constant $k$.
\item Tabulate the probability distribution of $X$.
\item Calculate $E(X)$ and $\operatorname{Var}(X)$.
\item Calculate $P(X \geq 3)$ and $P(X > 1 \mid X \leq 3)$.
\end{enumerate}
\end{exercice}

\begin{exercice}[Recovering the mass function from $F$]
The cumulative distribution function of a discrete random variable $X$ is given by
\begin{center}
\begin{tabular}{|c|c|c|c|c|c|c|}
\hline
\rowcolor{popblueL} $x$ & $1$ & $2$ & $3$ & $4$ & $5$ & $6$ \\
\hline
$F(x)$ & $0.05$ & $0.20$ & $0.45$ & $0.70$ & $0.90$ & $1$ \\
\hline
\end{tabular}
\end{center}
\begin{enumerate}
\item Recover the probability mass function of $X$ and present it in a table.
\item Find the median of $X$.
\item Calculate $P(2 < X \leq 5)$.
\item Calculate $E(X)$.
\end{enumerate}
\end{exercice}

\begin{exercice}[A game in FCFA]
A game costs $500\ \mathrm{FCFA}$ to play. A player draws one ball at random from a bag containing $3$ red balls, $2$ blue balls and $1$ green ball. The player wins $1000\ \mathrm{FCFA}$ for a red ball, $2000\ \mathrm{FCFA}$ for a blue ball and $5000\ \mathrm{FCFA}$ for a green ball. Let $G$ be the player's net gain (winnings minus the entry fee).
\begin{enumerate}
\item Construct the probability distribution of $G$.
\item Compute the expected gain $E(G)$. Is the game fair? Is it favourable to the player or to the organiser?
\item Find the entry fee that would make the game fair.
\end{enumerate}
\end{exercice}

\begin{exercice}[Using a p.g.f.]
The probability generating function of a discrete random variable $X$ is
\[ G(t) = \tfrac{1}{8}(1 + t)^3. \]
\begin{enumerate}
\item Expand $G(t)$ and hence write down the probability distribution of $X$.
\item Find $P(X = 2)$.
\item Calculate $G'(1)$ and $G''(1)$, and deduce $E(X)$ and $\operatorname{Var}(X)$.
\item Describe a simple experiment (involving coins) whose outcome has this distribution.
\end{enumerate}
\end{exercice}

\courssub{I prepare for the GCE}

\begin{exercice}[Distribution with an unknown probability \hfill (12 marks)]
The probability distribution of a discrete random variable $X$ is given by
\begin{center}
\begin{tabular}{|c|c|c|c|c|c|}
\hline
\rowcolor{popblueL} $x$ & $0$ & $1$ & $2$ & $3$ & $4$ \\
\hline
$P(X = x)$ & $0.1$ & $p$ & $0.3$ & $0.2$ & $q$ \\
\hline
\end{tabular}
\end{center}
where $p$ and $q$ are constants, and it is known that $E(X) = 2.1$.
\begin{enumerate}
\item Write down two equations satisfied by $p$ and $q$, and hence find their values. \hfill (4 marks)
\item Calculate $\operatorname{Var}(X)$ and the standard deviation of $X$, giving your answers to $3$ significant figures. \hfill (4 marks)
\item Construct the cumulative distribution function $F(x)$ in a table, and state the median and the mode of $X$. \hfill (2 marks)
\item Two independent observations of $X$ are made. Find the probability that their sum equals $4$. \hfill (2 marks)
\end{enumerate}
\end{exercice}

\begin{exercice}[Sum of two dice \hfill (13 marks)]
A fair tetrahedral die has faces numbered $1, 2, 3, 4$; when it is thrown, the score $X_1$ is the number on the face on which it lands, each face being equally likely. A biased ordinary die has score $X_2$ with distribution
\begin{center}
\begin{tabular}{|c|c|c|c|c|c|c|}
\hline
\rowcolor{popblueL} $x$ & $1$ & $2$ & $3$ & $4$ & $5$ & $6$ \\
\hline
$P(X_2 = x)$ & $0.1$ & $0.1$ & $0.2$ & $0.2$ & $0.2$ & $0.2$ \\
\hline
\end{tabular}
\end{center}
The two dice are thrown independently and $S = X_1 + X_2$ is the total score.
\begin{enumerate}
\item Show that $E(X_1) = 2.5$ and find $E(X_2)$. \hfill (3 marks)
\item Construct the probability distribution of $S$ in a table. \hfill (4 marks)
\item Find the mode and the median of $S$. \hfill (2 marks)
\item Verify that $E(S) = E(X_1) + E(X_2)$. \hfill (2 marks)
\item Find $P(S \geq 8)$. \hfill (2 marks)
\end{enumerate}
\end{exercice}

\begin{exercice}[Comparing two investments \hfill (15 marks)]
A trader in Douala can invest the same capital in one of two projects, $A$ or $B$. The annual profit, in hundreds of thousands of FCFA, is modelled by the random variables $X_A$ and $X_B$ with distributions
\begin{center}
\begin{tabular}{|c|c|c|c|c|}
\hline
\rowcolor{popblueL} $x$ & $-2$ & $1$ & $3$ & $6$ \\
\hline
$P(X_A = x)$ & $0.1$ & $0.3$ & $0.4$ & $0.2$ \\
\hline
$P(X_B = x)$ & $0.05$ & $0.25$ & $0.45$ & $0.25$ \\
\hline
\end{tabular}
\end{center}
\begin{enumerate}
\item Calculate $E(X_A)$ and $E(X_B)$. \hfill (4 marks)
\item Calculate the standard deviation of each of $X_A$ and $X_B$, giving your answers to $3$ significant figures. \hfill (5 marks)
\item The probability generating function of $X_A$ is $G_A(t)$. Write down $G_A(t)$ as an expression in $t$ and verify that $G_A'(1) = E(X_A)$. \hfill (3 marks)
\item A \emph{risk-averse} client wants a good expected profit but dislikes uncertainty. Using your answers to parts (a) and (b), advise the client which project to choose, justifying your recommendation clearly. \hfill (3 marks)
\end{enumerate}
\end{exercice}



\newpage

\coursec{Solutions to the exercises}

\begin{corrige}[Exercise 1 — MCQ: spotting a valid distribution]
Recall that a table defines a probability distribution of a discrete random variable $X$ if and only if (i) every probability satisfies $0 \le P(X=x) \le 1$, and (ii) $\sum_x P(X=x) = 1$. Both conditions must be checked.
\begin{enumerate}
\item The probabilities are $\dfrac{1}{10}, \dfrac{2}{10}, \dfrac{3}{10}, \dfrac{4}{10}$, all between $0$ and $1$, and
\[ \frac{1+2+3+4}{10} = \frac{10}{10} = 1. \]
\textbf{Valid distribution.}
\item $P(X=3) = -0.1 < 0$: a probability can never be negative. (Note that the sum is $1.0$, so condition (ii) holds, but condition (i) fails — one single bad value is enough to reject the table.) \textbf{Not valid.}
\item The probabilities are $\dfrac{1}{2}, \dfrac{1}{4}, \dfrac{1}{8}, \dfrac{1}{16}$, all in $[0,1]$, but
\[ \frac{1}{2} + \frac{1}{4} + \frac{1}{8} + \frac{1}{16} = \frac{8+4+2+1}{16} = \frac{15}{16} \neq 1. \]
\textbf{Not valid} (condition (ii) fails: the total probability is not $1$).
\item Each probability is $0.25 \in [0,1]$ and $4 \times 0.25 = 1$. \textbf{Valid distribution} (this is the discrete uniform distribution on $\{0,1,2,3\}$).
\end{enumerate}
\end{corrige}

\begin{corrige}[Exercise 2 — True or false?]
\begin{enumerate}
\item \textbf{True.} The events $\{X = x\}$, as $x$ runs through the values taken by $X$, form a partition of the sample space $\Omega$ (they are mutually exclusive and their union is $\Omega$). Hence the total probability is $\sum_x P(X=x) = P(\Omega) = 1$. This is the fundamental condition satisfied by every probability distribution.
\item \textbf{True.} This is the computational (König–Huygens) formula for the variance. Using the linearity of expectation with $\mu = E(X)$:
\begin{align*}
\operatorname{Var}(X) &= E\!\left[(X - \mu)^2\right] = E(X^2 - 2\mu X + \mu^2) \\
&= E(X^2) - 2\mu\, E(X) + \mu^2 = E(X^2) - 2\mu^2 + \mu^2 = E(X^2) - [E(X)]^2.
\end{align*}
\item \textbf{False.} Counter-example: let $X$ take the values $1$ and $2$, each with probability $0.5$. Then $F(1) = 0.5$ and $F(x) = 1$ for all $x \ge 2$. With the midpoint convention, every $m \in [1,2]$ satisfies $P(X \le m) \ge 0.5$ and $P(X \ge m) \ge 0.5$, and one usually takes the median to be $m = 1.5$, which is \emph{not} a value taken by $X$. (With the GCE convention "the median is the smallest value $x$ such that $F(x) \ge 0.5$", the median here is $1$, which \emph{is} a value of $X$; the statement is false in general because under the midpoint convention the median need not belong to the set of values of $X$.)
\item \textbf{False.} By definition $G_X(t) = E(t^X) = \sum_x t^x P(X=x)$, so substituting $t = 1$ gives $G_X(1) = \sum_x P(X=x) = 1$, not $E(X)$. It is the \emph{derivative} evaluated at $t = 1$ that gives the mean: $G_X'(1) = E(X)$.
\item \textbf{True.} Expectation is additive for \emph{any} two random variables, independent or not: $E(X_1 + X_2) = E(X_1) + E(X_2)$. Independence is required for the \emph{variance} formula $\operatorname{Var}(X_1+X_2) = \operatorname{Var}(X_1) + \operatorname{Var}(X_2)$; when independence holds, the expectation formula holds a fortiori.
\end{enumerate}
\end{corrige}

\begin{corrige}[Exercise 3 — Definitions in your own words]
\begin{enumerate}
\item A \cle{random variable} is a real-valued function defined on the sample space $\Omega$ of a random experiment: to each outcome $\omega \in \Omega$ it associates a real number $X(\omega)$. A random variable is \cle{discrete} if it takes only finitely many or countably many values (e.g. the number of goals scored by a football team in a match, taking values in $\mathbb{N}$). It is \cle{continuous} if it can take any value in an interval of $\mathbb{R}$ (e.g. the height of a randomly chosen student in a lycée in Yaoundé, taking values in an interval such as $[1.40, 2.00]$ metres).
\item The \cle{cumulative distribution function} of $X$ is the function $F$ defined for every real number $x$ by $F(x) = P(X \le x)$. Two properties: (i) $F$ is non-decreasing (if $a \le b$ then $F(a) \le F(b)$, because the event $\{X \le a\}$ is contained in $\{X \le b\}$); (ii) $F(x) \to 0$ as $x \to -\infty$ and $F(x) \to 1$ as $x \to +\infty$. For a discrete variable, $F$ is a step function: it is constant between two consecutive values of $X$, it is right-continuous, and it has a jump of height $P(X=x)$ at each value $x$ taken by $X$.
\item For a discrete random variable $X$ taking values in $\mathbb{N}$, the \cle{probability generating function} is
\[ G_X(t) = E(t^X) = \sum_{x=0}^{+\infty} P(X = x)\, t^x. \]
Differentiating term by term, $G_X'(t) = \sum_x x\, P(X=x)\, t^{x-1}$, so setting $t = 1$ gives $G_X'(1) = \sum_x x\,P(X=x)$, that is
\[ E(X) = G_X'(1). \]
(One also obtains the variance from $\operatorname{Var}(X) = G_X''(1) + G_X'(1) - [G_X'(1)]^2$.)
\end{enumerate}
\end{corrige}

\begin{corrige}[Exercise 4 — Direct calculations]
\begin{enumerate}
\item All four probabilities lie in $[0,1]$ and
\[ 0.1 + 0.3 + 0.4 + 0.2 = 1.0. \]
Both conditions are satisfied, hence the table defines a valid probability distribution.
\item 
\begin{align*}
E(X) &= 1(0.1) + 2(0.3) + 3(0.4) + 4(0.2) = 0.1 + 0.6 + 1.2 + 0.8 = 2.7. \\
E(X^2) &= 1^2(0.1) + 2^2(0.3) + 3^2(0.4) + 4^2(0.2) = 0.1 + 1.2 + 3.6 + 3.2 = 8.1. \\
\operatorname{Var}(X) &= E(X^2) - [E(X)]^2 = 8.1 - (2.7)^2 = 8.1 - 7.29 = 0.81.
\end{align*}
The standard deviation is $\sigma = \sqrt{\operatorname{Var}(X)} = \sqrt{0.81} = 0.9$.
\item $F(2) = P(X \le 2) = P(X=1) + P(X=2) = 0.1 + 0.3 = 0.4$. Since $X$ takes no value strictly between $3$ and $4$, the event $\{X \le 3.5\}$ is the same as $\{X \le 3\}$, so
\[ F(3.5) = P(X \le 3.5) = P(X \le 3) = 0.1 + 0.3 + 0.4 = 0.8. \]
\item The mode is the value with the largest probability. The largest probability is $0.4$, attained at $x = 3$. The mode of $X$ is $3$.
\end{enumerate}
\end{corrige}

\begin{corrige}[Exercise 5 — Finding the constant $k$]
\begin{enumerate}
\item The total probability must equal $1$:
\[ \sum_{x=1}^{4} kx = k(1 + 2 + 3 + 4) = 10k = 1 \implies k = \frac{1}{10} = 0.1. \]
(Note that $k = 0.1 \ge 0$, so all the probabilities $kx$ are non-negative, as required.)
\item The distribution of $X$:
\begin{center}
\begin{tabular}{|c|c|c|c|c|}
\hline
\rowcolor{popblueL} $x$ & $1$ & $2$ & $3$ & $4$ \\
\hline
$P(X = x)$ & $0.1$ & $0.2$ & $0.3$ & $0.4$ \\
\hline
\end{tabular}
\end{center}
\item 
\begin{align*}
E(X) &= 1(0.1) + 2(0.2) + 3(0.3) + 4(0.4) = 0.1 + 0.4 + 0.9 + 1.6 = 3.0. \\
E(X^2) &= 1^2(0.1) + 2^2(0.2) + 3^2(0.3) + 4^2(0.4) = 0.1 + 0.8 + 2.7 + 6.4 = 10.0. \\
\operatorname{Var}(X) &= E(X^2) - [E(X)]^2 = 10.0 - (3.0)^2 = 10 - 9 = 1.0.
\end{align*}
\item $P(X \ge 3) = P(X=3) + P(X=4) = 0.3 + 0.4 = 0.7$.

For the conditional probability, using $P(A \mid B) = \dfrac{P(A \cap B)}{P(B)}$ with $A = \{X > 1\}$ and $B = \{X \le 3\}$, so that $A \cap B = \{1 < X \le 3\} = \{X = 2\} \cup \{X = 3\}$:
\[ P(X > 1 \mid X \le 3) = \frac{P(1 < X \le 3)}{P(X \le 3)} = \frac{P(X=2) + P(X=3)}{P(X=1) + P(X=2) + P(X=3)} = \frac{0.2 + 0.3}{0.1 + 0.2 + 0.3} = \frac{0.5}{0.6} = \frac{5}{6}. \]
\end{enumerate}
\end{corrige}

\begin{corrige}[Exercise 6 — Recovering the mass function from $F$]
\begin{enumerate}
\item For a discrete variable taking integer values, each jump of $F$ gives a probability: $P(X = x) = F(x) - F(x-1)$. Here the smallest value is $1$, so $F(0) = 0$:
\begin{align*}
P(X=1) &= F(1) - F(0) = 0.05 - 0 = 0.05, \\
P(X=2) &= F(2) - F(1) = 0.20 - 0.05 = 0.15, \\
P(X=3) &= F(3) - F(2) = 0.45 - 0.20 = 0.25, \\
P(X=4) &= F(4) - F(3) = 0.70 - 0.45 = 0.25, \\
P(X=5) &= F(5) - F(4) = 0.90 - 0.70 = 0.20, \\
P(X=6) &= F(6) - F(5) = 1 - 0.90 = 0.10.
\end{align*}
\begin{center}
\begin{tabular}{|c|c|c|c|c|c|c|}
\hline
\rowcolor{popblueL} $x$ & $1$ & $2$ & $3$ & $4$ & $5$ & $6$ \\
\hline
$P(X = x)$ & $0.05$ & $0.15$ & $0.25$ & $0.25$ & $0.20$ & $0.10$ \\
\hline
\end{tabular}
\end{center}
Check: $0.05 + 0.15 + 0.25 + 0.25 + 0.20 + 0.10 = 1$. \checkmark
\item The median is the smallest value $x$ such that $F(x) \ge 0.5$. We have $F(3) = 0.45 < 0.5$ and $F(4) = 0.70 \ge 0.5$, so the median is $4$.
\item Using $P(a < X \le b) = F(b) - F(a)$:
\[ P(2 < X \le 5) = F(5) - F(2) = 0.90 - 0.20 = 0.70. \]
(Check directly: $P(X=3) + P(X=4) + P(X=5) = 0.25 + 0.25 + 0.20 = 0.70$. \checkmark)
\item 
\begin{align*}
E(X) &= 1(0.05) + 2(0.15) + 3(0.25) + 4(0.25) + 5(0.20) + 6(0.10) \\
&= 0.05 + 0.30 + 0.75 + 1.00 + 1.00 + 0.60 = 3.70.
\end{align*}
\end{enumerate}
\end{corrige}

\begin{corrige}[Exercise 7 — A game in FCFA]
\begin{enumerate}
\item There are $3 + 2 + 1 = 6$ equiprobable balls. The net gain is $G = \text{winnings} - 500$:
\begin{itemize}
\item Red ball: $G = 1000 - 500 = 500$ FCFA, with probability $\dfrac{3}{6} = \dfrac{1}{2}$;
\item Blue ball: $G = 2000 - 500 = 1500$ FCFA, with probability $\dfrac{2}{6} = \dfrac{1}{3}$;
\item Green ball: $G = 5000 - 500 = 4500$ FCFA, with probability $\dfrac{1}{6}$.
\end{itemize}
\begin{center}
\begin{tabular}{|c|c|c|c|}
\hline
\rowcolor{popblueL} $g$ (FCFA) & $500$ & $1500$ & $4500$ \\
\hline
$P(G = g)$ & $\dfrac{1}{2}$ & $\dfrac{1}{3}$ & $\dfrac{1}{6}$ \\
\hline
\end{tabular}
\end{center}
(Check: $\tfrac12 + \tfrac13 + \tfrac16 = \tfrac{3+2+1}{6} = 1$. \checkmark)
\item 
\[ E(G) = 500 \times \frac{1}{2} + 1500 \times \frac{1}{3} + 4500 \times \frac{1}{6} = 250 + 500 + 750 = 1500\ \mathrm{FCFA}. \]
A game is fair when the expected net gain is zero. Since $E(G) = +1500 \neq 0$, the game is \textbf{not fair}. As $E(G) > 0$, on average the player wins $1500\ \mathrm{FCFA}$ per game: the game is \textbf{favourable to the player} (and unfavourable to the organiser — such a generous organiser would soon go bankrupt!).
\item Let $c$ be the new entry fee in FCFA. The net gain becomes $G' = G + 500 - c$ (the player now pays $c$ instead of $500$), so by linearity of expectation,
\[ E(G') = E(G) + 500 - c = 1500 + 500 - c = 2000 - c. \]
The game is fair when $E(G') = 0$, i.e. $c = 2000\ \mathrm{FCFA}$.

Equivalently, the fair price equals the expected winnings $W$: using exact fractions,
\[ E(W) = 1000 \times \frac{1}{2} + 2000 \times \frac{1}{3} + 5000 \times \frac{1}{6} = 500 + \frac{2000}{3} + \frac{2500}{3} = 500 + \frac{4500}{3} = 500 + 1500 = 2000\ \mathrm{FCFA}. \]
\end{enumerate}
\end{corrige}

\begin{corrige}[Exercise 8 — Using a p.g.f.]
\begin{enumerate}
\item Expand with the binomial theorem: $(1+t)^3 = 1 + 3t + 3t^2 + t^3$, so
\[ G(t) = \frac{1}{8}(1+t)^3 = \frac{1}{8} + \frac{3}{8}t + \frac{3}{8}t^2 + \frac{1}{8}t^3. \]
Since $G(t) = \sum_x P(X=x)\, t^x$, the coefficient of $t^x$ is exactly $P(X=x)$. Reading off the coefficients:
\begin{center}
\begin{tabular}{|c|c|c|c|c|}
\hline
\rowcolor{popblueL} $x$ & $0$ & $1$ & $2$ & $3$ \\
\hline
$P(X = x)$ & $\dfrac{1}{8}$ & $\dfrac{3}{8}$ & $\dfrac{3}{8}$ & $\dfrac{1}{8}$ \\
\hline
\end{tabular}
\end{center}
Check: $\dfrac{1+3+3+1}{8} = \dfrac{8}{8} = 1$. \checkmark
\item $P(X = 2)$ is the coefficient of $t^2$: $P(X = 2) = \dfrac{3}{8} = 0.375$.
\item Differentiate twice:
\[ G'(t) = \frac{3}{8}(1+t)^2, \qquad G''(t) = \frac{6}{8}(1+t) = \frac{3}{4}(1+t). \]
Evaluating at $t = 1$:
\[ G'(1) = \frac{3}{8} \times (1+1)^2 = \frac{3}{8} \times 4 = \frac{3}{2} = 1.5, \qquad G''(1) = \frac{3}{4} \times 2 = \frac{3}{2} = 1.5. \]
Therefore
\[ E(X) = G'(1) = 1.5, \]
\[ \operatorname{Var}(X) = G''(1) + G'(1) - [G'(1)]^2 = 1.5 + 1.5 - (1.5)^2 = 3 - 2.25 = 0.75. \]
(Check by direct computation: $E(X) = 0 \times \tfrac18 + 1 \times \tfrac38 + 2 \times \tfrac38 + 3 \times \tfrac18 = \tfrac{3+6+3}{8} = \tfrac{12}{8} = 1.5$; $E(X^2) = 0 + \tfrac38 + 4 \times \tfrac38 + 9 \times \tfrac18 = \tfrac{3+12+9}{8} = 3$, so $\operatorname{Var}(X) = 3 - 1.5^2 = 0.75$. \checkmark)
\item Toss three fair coins and let $X$ be the number of heads obtained. There are $2^3 = 8$ equiprobable outcomes, and the number of outcomes giving exactly $x$ heads is $\binom{3}{x}$, so $P(X=x) = \binom{3}{x}/8$, i.e. $\tfrac18, \tfrac38, \tfrac38, \tfrac18$ for $x = 0,1,2,3$ — exactly the distribution found above. (This is the binomial distribution $B(3, \tfrac12)$, whose p.g.f. is $(\tfrac12 + \tfrac12 t)^3 = \tfrac18(1+t)^3$.)
\end{enumerate}
\end{corrige}

\newpage

\coursec{Summary sheet}

\courssub{Summary Sheet}

\begin{popfigure}
\centering
\begin{tikzpicture}[
    level 1/.style={sibling distance=4.5cm, level distance=3.2cm},
    level 2/.style={sibling distance=2.2cm, level distance=2.5cm},
    every node/.style={text width=2.8cm, align=center, font=\small\sffamily, rounded corners=4pt, inner sep=4pt},
    concept/.style={draw, thick, fill=popblue!80, text=white, minimum width=3.5cm, minimum height=1.2cm, font=\normalsize\bfseries\sffamily},
    branch1/.style={draw, thick, fill=popgreen!70, text=white},
    branch2/.style={draw, thick, fill=poporange!70, text=white},
    branch3/.style={draw, thick, fill=poppurple!70, text=white},
    branch4/.style={draw, thick, fill=poppink!70, text=white},
    branch5/.style={draw, thick, fill=popgold!70, text=white},
    branch6/.style={draw, thick, fill=popdark!70, text=white},
    edge from parent/.style={draw, very thick, ->, shorten >=2pt, shorten <=2pt}
]
\node[concept] {Discrete Random Variables}
    child[branch1] { node[align=center]{1. Definitions\\Random Variable\\Discrete RV\\Support $R_X$} }
    child[branch2] { node[align=center]{2. Probability Distribution\\PMF $p(x)$\\Table/Formula\\$\sum p(x)=1$} }
    child[branch3] { node[align=center]{3. Expectation \& Variance\\$E(X), E(g(X))$\\$Var(X), \sigma$\\Linear Transforms} }
    child[branch4] { node[align=center]{4. CDF, Mode, Median\\$F(x)=P(X\le x)$\\Mode: max $p(x)$\\Median: $F(x)\ge 0.5$} }
    child[branch5] { node[align=center]{5. Sum $X_1+X_2$\\Convolution Table\\Independence\\$Var(X+Y)=Var(X)+Var(Y)$} }
    child[branch6] { node[align=center]{6. Probability Generating Function\\$G(t)=E(t^X)$\\Moments: $G'(1), G''(1)$\\Sum: $G_{X+Y}=G_X G_Y$} };
\end{tikzpicture}
\legende{Chapter PMS11 Mind Map: Discrete Random Variables}
\end{popfigure}

\begin{aretenir}[Key Definitions]
\begin{itemize}
    \item \cle{Random Variable (RV)}: A function $X: \Omega \to \mathbb{R}$ mapping each outcome of a random experiment to a real number. We denote the RV by a capital letter ($X$) and its values by small letters ($x$).
    \item \cle{Discrete RV}: Takes a finite or countably infinite set of values $\{x_1, x_2, \dots\}$. The set of all possible values is the \cle{support} $R_X$.
    \item \cle{Probability Mass Function (PMF)}: $p(x) = P(X=x)$ for $x \in R_X$. Must satisfy:
          \begin{enumerate}
              \item $p(x) \ge 0$ for all $x \in R_X$,
              \item $\sum_{x \in R_X} p(x) = 1$.
          \end{enumerate}
    \item \cle{Cumulative Distribution Function (CDF)}: $F(x) = P(X \le x) = \sum_{x_i \le x} p(x_i)$, defined for all $x \in \mathbb{R}$. Properties: non-decreasing, right-continuous, $F(-\infty)=0$, $F(+\infty)=1$.
    \item \cle{Mode}: The value(s) $x \in R_X$ where $p(x)$ attains its maximum. A distribution may have one mode (unimodal) or several.
    \item \cle{Median}: Any value $m$ such that $P(X \le m) \ge 0.5$ and $P(X \ge m) \ge 0.5$. For discrete RVs, the conventional median is the smallest $x$ with $F(x) \ge 0.5$.
    \item \cle{Probability Generating Function (PGF)}: $G_X(t) = E(t^X) = \sum_{x \in R_X} p(x) t^x$, valid for $|t| \le 1$. The PGF uniquely determines the distribution.
\end{itemize}
\end{aretenir}

\begin{propriete}[Essential Laws, Formulas \& Theorems]
\textbf{Examinable proofs are marked with an asterisk (*).}
\begin{itemize}
    \item \textbf{Normalisation}: $\sum_{x \in R_X} p(x) = 1$. *Proof: follows from the axioms of probability (total probability of the sample space is 1).*
    \item \textbf{Expectation}: $E(X) = \sum_{x \in R_X} x\,p(x)$. For any function $g$, $E(g(X)) = \sum g(x)\,p(x)$ (Law of the Unconscious Statistician).
    \item \textbf{Linearity of Expectation}: $E(aX+b) = aE(X)+b$ for constants $a,b$. $E(X+Y) = E(X)+E(Y)$ \emph{always holds}, independence not required.
    \item \textbf{Variance}: $Var(X) = E[(X-\mu)^2] = E(X^2) - [E(X)]^2$, where $\mu = E(X)$. Standard deviation $\sigma_X = \sqrt{Var(X)}$.
    \item \textbf{Variance Properties}: $Var(aX+b) = a^2 Var(X)$. *If $X$ and $Y$ are independent:* $Var(X+Y) = Var(X)+Var(Y)$.
    \item \textbf{CDF Properties}:
          \begin{align*}
              P(a < X \le b) &= F(b) - F(a), \\
              P(X = x) &= F(x) - F(x^-) \quad \text{(jump at $x$)}.
          \end{align*}
    \item \textbf{Sum of Independent RVs (Convolution)}: If $X_1, X_2$ independent,
          \[
              P(X_1+X_2 = z) = \sum_{x \in R_{X_1}} P(X_1=x)\,P(X_2=z-x).
          \]
    \item \textbf{PGF Properties}:
          \begin{align*}
              G_X(1) &= 1 \quad \text{(check)}, \\
              G'_X(1) &= E(X), \\
              G''_X(1) &= E[X(X-1)] = E(X^2) - E(X), \\
              Var(X) &= G''_X(1) + G'_X(1) - [G'_X(1)]^2.
          \end{align*}
    \item \textbf{PGF of Sum}: If $X,Y$ independent, $G_{X+Y}(t) = G_X(t) G_Y(t)$. *Proof: $E(t^{X+Y}) = E(t^X t^Y) = E(t^X)E(t^Y)$ by independence.*
    \item \textbf{Uniqueness}: The PGF uniquely determines the probability distribution (examinable concept).
\end{itemize}
\end{propriete}

\begin{methode}[Standard Methods (Step-by-Step)]
\begin{enumerate}
    \item \textbf{Constructing a Distribution Table}:
          \begin{enumerate}
              \item Define the RV clearly: ``Let $X$ be the number of...''
              \item List all possible values $x \in R_X$.
              \item Compute $p(x) = P(X=x)$ using combinatorics, given probabilities, or a formula.
              \item Verify $\sum p(x) = 1$ explicitly (write the sum and conclude $=1$).
              \item Add columns for $xp(x)$ and $x^2p(x)$; compute totals for $E(X)$ and $E(X^2)$.
          \end{enumerate}
    \item \textbf{Finding Unknown Constants (e.g., $k$)}:
          \begin{enumerate}
              \item Write the PMF with the unknown parameter(s).
              \item Use $\sum p(x) = 1$ to form an equation.
              \item Solve for the parameter. Check that all $p(x) \ge 0$.
          \end{enumerate}
    \item \textbf{Calculating $E(X)$ and $Var(X)$}:
          \begin{enumerate}
              \item Extend the table with columns $xp(x)$ and $x^2p(x)$.
              \item $\mu = \sum xp(x)$.
              \item $E(X^2) = \sum x^2p(x)$.
              \item $Var(X) = E(X^2) - \mu^2$; $\sigma = \sqrt{Var(X)}$.
          \end{enumerate}
    \item \textbf{Linear Transform $Y = aX+b$}:
          \begin{enumerate}
              \item Support $R_Y = \{ax+b : x \in R_X\}$.
              \item $p_Y(y) = p_X((y-b)/a)$.
              \item $E(Y) = a\mu + b$, $Var(Y) = a^2\sigma^2$.
          \end{enumerate}
    \item \textbf{CDF, Mode, Median}:
          \begin{enumerate}
              \item Compute cumulative sums for $F(x)$ in a table.
              \item Mode = $x$ with highest $p(x)$.
              \item Median = smallest $x$ such that $F(x) \ge 0.5$. (If $F(x)=0.5$ exactly, any value in $[x, x_{\text{next}}]$ is a median; GCE usually expects the smallest integer $x$.)
          \end{enumerate}
    \item \textbf{Distribution of $X_1+X_2$ (Convolution)}:
          \begin{enumerate}
              \item Draw a two-way table of $(x_1, x_2)$ with joint probabilities $p_1(x_1)p_2(x_2)$ (requires independence).
              \item Sum probabilities along diagonals where $x_1+x_2 = z$.
              \item Present the resulting distribution of $Z = X_1+X_2$.
          \end{enumerate}
    \item \textbf{Using the PGF}:
          \begin{enumerate}
              \item Write $G(t) = \sum p(x)t^x$.
              \item Differentiate: $G'(t) = \sum x p(x) t^{x-1}$, $G''(t) = \sum x(x-1) p(x) t^{x-2}$.
              \item Evaluate at $t=1$: $E(X)=G'(1)$, $E(X^2)=G''(1)+G'(1)$.
              \item For sum of $n$ i.i.d. RVs: $G_{S_n}(t) = [G_X(t)]^n$. Expand to read off the distribution if needed.
          \end{enumerate}
\end{enumerate}
\end{methode}

\begin{attention}[Common Mistakes \& Traps]
\begin{itemize}
    \item \textbf{Forgetting $\sum p(x)=1$}: Always verify or use it to find constants. Easy mark lost if omitted.
    \item \textbf{Confusing PMF and CDF}: $p(x)$ is point probability; $F(x)$ is cumulative. $P(X=x) = F(x) - F(x^-)$, not $F(x)-F(x-1)$ unless the support is integers with step 1.
    \item \textbf{Variance Errors}: $Var(aX) = a^2 Var(X)$, \emph{not} $a Var(X)$. $Var(X+Y) = Var(X)+Var(Y)$ \emph{only if $X,Y$ are independent}.
    \item \textbf{Median Definition}: For discrete RV, median is \emph{not} always the ``middle value''. It is the smallest $x$ with $F(x) \ge 0.5$. If $F(x)=0.5$ exactly, any value in $[x, x_{\text{next}}]$ is a median; GCE usually expects the smallest integer $x$.
    \item \textbf{Convolution Limits}: When summing $P(X_1+X_2=z)$, ensure $x$ runs over $R_{X_1}$ such that $z-x \in R_{X_2}$. Do not assume infinite sums; the support is finite or countable.
    \item \textbf{PGF Domain}: $G(t)$ defined for $|t|\le 1$. Differentiation valid inside radius of convergence. Always state $G(1)=1$ as a check.
    \item \textbf{Notation}: Use $X$ for RV, $x$ for values. $E(X)$ is a number, not a function. Distinguish $X_1, X_2$ (different RVs) from $x_1, x_2$ (values).
    \item \textbf{``Hence'' in PGF Questions}: If the question says ``Hence find $E(X)$ and $Var(X)$'' after finding $G(t)$, you \emph{must} differentiate $G(t)$; computing directly from the table loses method marks.
\end{itemize}
\end{attention}

\begin{aretenir}[GCE Examination Tips]
\begin{itemize}
    \item \textbf{Define the RV clearly}: ``Let $X$ be the number of...'' State support $R_X$ explicitly.
    \item \textbf{Show the $\sum p(x)=1$ check}: Even if obvious, write the sum and conclude $=1$. Easy mark.
    \item \textbf{Table Layout}: Use a clear table with columns $x$, $p(x)$, $xp(x)$, $x^2p(x)$. Show a totals row. Neat tables gain method marks.
    \item \textbf{Independence}: Explicitly state ``Since $X_1, X_2$ are independent...'' before using $p(x_1,x_2)=p_1(x_1)p_2(x_2)$ or $Var(X_1+X_2)=Var(X_1)+Var(X_2)$.
    \item \textbf{PGF Questions}: Often ask to ``Find $G(t)$'', then ``Hence find $E(X)$ and $Var(X)$''. Do not compute $E(X)$ directly from table if ``Hence'' is used; must differentiate $G(t)$.
    \item \textbf{Sum of $n$ i.i.d.}: Recognize $G_{S_n}(t) = [G_X(t)]^n$. If asked for distribution of sum, expand the polynomial/binomial to read off coefficients.
    \item \textbf{Contextual Problems}: Cameroonian contexts appear (e.g., ``A bag contains 5 red and 3 blue balls...'', ``Number of customers at a Douala kiosk...''). Translate words into RV definition + distribution (Binomial, Geometric, Uniform, or ad-hoc).
    \item \textbf{Working}: Show all steps. $E(X^2) = \sum x^2 p(x)$ written out gets method marks even if arithmetic error later.
    \item \textbf{Approximation}: Not in this chapter (Poisson approximation to Binomial is in the next chapter). Stick to exact discrete calculations.
\end{itemize}
\end{aretenir}

\begin{savaistu}[Did You Know?]
The concept of a random variable was formalised by Andrey Kolmogorov in 1933 as part of his axiomatisation of probability theory. Before that, probabilities were tied to specific games or experiments. The Probability Generating Function was introduced by Pierre-Simon Laplace in 1812 for solving problems in combinatorics and probability; it is the discrete analogue of the Moment Generating Function used for continuous variables. In Cameroon, discrete random variables model many real situations: the number of successful candidates in a GCE session (Binomial), the number of taxi arrivals at a rank in 10 minutes (Poisson), or the number of draws until the first winning ticket in a \emph{loterie nationale} (Geometric).
\end{savaistu}

\findecours

\end{document}
